% ---- hand chunk 00-preamble.tex ----
\documentclass[11pt,a4paper,oneside,openany]{book}
% Gaussian Coefficient Calculus: consolidated companion volume to
% Combinatorial_Coefficient_Calculus.  Assembled 2026-09-04 from the three
% q-binomial coefficient-calculus arrivals filed in the 2026-09-04 intake;
% see the provenance chapter and README.md.
\usepackage[T1]{fontenc}
\usepackage[utf8]{inputenc}
\IfFileExists{libertinus.sty}{\usepackage{libertinus}}{\usepackage{lmodern}}
\usepackage{microtype}
\usepackage[a4paper,margin=27mm,headheight=15pt]{geometry}
\usepackage{amsmath,amssymb,amsthm,mathtools,bm,mathrsfs}
\usepackage{booktabs,longtable,array,tabularx,ragged2e}
\usepackage{enumitem,xcolor}
\usepackage[most]{tcolorbox}
\usepackage{fancyhdr,xurl,hyperref}
\usepackage[nameinlink,noabbrev,capitalise]{cleveref}
% Canonical mathematical notation for active FabiusFunction LaTeX documents.
%
% This file is intentionally a plain TeX input rather than a LaTeX package.
% Every standalone document loads it by an explicit path relative to that
% document's directory; included fragments inherit the definitions from their
% root document.  Keep this file independent of document layout and metadata.
%
% Contract:
%   * one command has one mathematical meaning throughout the active corpus;
%   * names encode semantic roles rather than a document-local choice of letter;
%   * visually similar but differently normalized objects have different print;
%   * every command below is defined normatively in the notation catalogue.
%
% The active corpus excludes docs/archive/.  Historical sources may retain
% their original notation only inside an explicitly labelled quotation.

\makeatletter
\@ifundefined{mathbb}{%
  \PackageError{fabius-notation}{The amssymb package must be loaded first}%
    {Load amsmath and amssymb before inputting fabius-notation.tex.}%
}{}
\@ifundefined{operatorname}{%
  \PackageError{fabius-notation}{The amsmath package must be loaded first}%
    {Load amsmath before inputting fabius-notation.tex.}%
}{}
\makeatother

% Version stamp used by the catalogue and by static audits.
\newcommand{\FabiusNotationVersion}{2026-08-31}

% Number systems.  NaturalNumbers is explicitly the nonnegative convention.
\newcommand{\NaturalNumbers}{\mathbb{N}_{0}}
\newcommand{\PositiveIntegers}{\mathbb{N}_{>0}}
\newcommand{\IntegerNumbers}{\mathbb{Z}}
\newcommand{\RationalNumbers}{\mathbb{Q}}
\newcommand{\RealNumbers}{\mathbb{R}}
\newcommand{\ComplexNumbers}{\mathbb{C}}
\newcommand{\ExtendedRealNumbers}{\overline{\mathbb{R}}}

% The three principal Fabius--Rvachev functions.  Subscripts are deliberate:
% a bare F and a calligraphic F have several unrelated uses in the corpus.
\newcommand{\FabiusBounded}{F_{\mathrm{bd}}}
\newcommand{\FabiusGlobal}{F_{\mathrm{gl}}}
\newcommand{\FabiusQuantile}{Q_{F}}
\newcommand{\FabiusClampedQuantile}{Q_{F,\mathrm{cl}}}
\newcommand{\RvachevUp}{\operatorname{up}}

% Constants, elementary operators, and delimiters.
\newcommand{\EulerE}{\mathrm{e}}
\newcommand{\ImaginaryUnit}{\mathrm{i}}
\newcommand{\Differential}{\mathop{}\!\mathrm{d}}
\newcommand{\AbsoluteValue}[1]{\left\lvert #1\right\rvert}
\newcommand{\Norm}[1]{\left\lVert #1\right\rVert}
\newcommand{\Floor}[1]{\left\lfloor #1\right\rfloor}
\newcommand{\Ceiling}[1]{\left\lceil #1\right\rceil}
\newcommand{\FractionalPart}[1]{\left\{#1\right\}}
\newcommand{\PositivePart}[1]{\left(#1\right)_{+}}
\newcommand{\IndicatorOf}[1]{\mathbf{1}_{#1}}
\newcommand{\SetOf}[1]{\left\{#1\right\}}
\newcommand{\InnerProduct}[1]{\left\langle #1\right\rangle}
\newcommand{\CardinalityOf}[1]{\left\lvert #1\right\rvert}
\newcommand{\DefinitionEquals}{\mathrel{\coloneqq}}
\newcommand{\Sign}{\operatorname{sgn}}
\newcommand{\IdentityOperator}{\operatorname{id}}
\newcommand{\SpanOperator}{\operatorname{span}}
\newcommand{\TraceOperator}{\operatorname{tr}}
\newcommand{\RankOperator}{\operatorname{rank}}
\newcommand{\DiagonalOperator}{\operatorname{diag}}
\newcommand{\ResidueOperator}{\operatorname*{Res}}
\newcommand{\LeastCommonMultiple}{\operatorname{lcm}}
\newcommand{\OrderOperator}{\operatorname{ord}}
\newcommand{\PrincipalLogarithm}{\operatorname{Log}}
\newcommand{\FallingFactorial}[2]{\left(#1\right)^{\underline{#2}}}
\newcommand{\RisingFactorial}[2]{\left(#1\right)^{\overline{#2}}}

% Support, topology, convolution, and metric notation.
\newcommand{\SupportOperator}{\operatorname{supp}}
\newcommand{\SupportOf}[1]{\SupportOperator\!\left(#1\right)}
\newcommand{\TopologicalSupportOf}[1]{\operatorname{tsupp}\!\left(#1\right)}
\newcommand{\Convolution}{\mathbin{*}}
\newcommand{\MetricDistance}{\operatorname{dist}}
\newcommand{\TotalVariationDistance}{d_{\mathrm{TV}}}

% Probability.  Equality and convergence in law are not metric distance.
\newcommand{\Probability}{\mathbb{P}}
\newcommand{\Expectation}{\mathbb{E}}
\newcommand{\Variance}{\operatorname{Var}}
\newcommand{\Covariance}{\operatorname{Cov}}
\newcommand{\LawOperator}{\operatorname{Law}}
\newcommand{\LawOf}[1]{\LawOperator\!\left(#1\right)}
\newcommand{\EqualInLaw}{\mathrel{\overset{\mathrm{law}}{=}}}
\newcommand{\ConvergesInLaw}{\mathrel{\xrightarrow{\mathrm{law}}}}

% Fourier and integral transforms.  The printed labels expose normalization.
% FourierTwoPi uses exp(-2*pi*i*x*xi); FourierAngular uses exp(-i*x*omega).
\newcommand{\FourierTwoPi}[1]{\widehat{#1}_{\,2\pi}}
\newcommand{\FourierAngular}[1]{\widehat{#1}_{\,\mathrm{ang}}}
\newcommand{\LaplaceTransformSymbol}{\mathcal{L}}
\newcommand{\MellinTransformSymbol}{\mathcal{M}}
\newcommand{\LaplaceTransformOf}[1]{\LaplaceTransformSymbol\!\left[#1\right]}
\newcommand{\MellinTransformOf}[1]{\MellinTransformSymbol\!\left[#1\right]}
\newcommand{\MomentGeneratingFunction}[1]{M_{#1}}
\newcommand{\CharacteristicFunction}[1]{\varphi_{#1}}

% The two standard sinc functions must never share one symbol.
\newcommand{\SincRad}{\operatorname{sinc}_{\mathrm{rad}}}
\newcommand{\SincPi}{\operatorname{sinc}_{\pi}}

% Binary arithmetic and Thue--Morse notation.
\newcommand{\BinaryDigitSum}{\operatorname{wt}_{2}}
\newcommand{\ThueMorseSignSymbol}{\varepsilon}
\newcommand{\ThueMorseSign}[1]{\varepsilon_{#1}}
\newcommand{\PadicValuation}[1]{\nu_{#1}}
\newcommand{\TwoAdicValuation}{\nu_{2}}

% q-series notation.  Every base is an explicit argument.
\newcommand{\QPochhammer}[3]{\left(#1;#2\right)_{#3}}
\newcommand{\GaussianBinomial}[3]{%
  \genfrac{[}{]}{0pt}{}{#1}{#2}_{#3}%
}
\newcommand{\QInteger}[2]{\left[#1\right]_{#2}}

% Named special functions.
\newcommand{\Polylogarithm}{\operatorname{Li}}
\newcommand{\LambertW}{W}
\newcommand{\GammaFunction}{\Gamma}
\newcommand{\RiemannZeta}{\zeta}

% Asymptotics.  BigO and LittleO are operator tokens for existing formula
% grammar; prefer the At forms so that the limiting regime is printed.
\newcommand{\BigO}{\mathcal{O}}
\newcommand{\LittleO}{o}
\newcommand{\BigOAt}[2]{\mathcal{O}_{#2}\!\left(#1\right)}
\newcommand{\LittleOAt}[2]{o_{#2}\!\left(#1\right)}

\definecolor{deepolive}{RGB}{67,79,45}
\definecolor{softolive}{RGB}{236,239,226}
\definecolor{darkteal}{RGB}{31,83,84}
\definecolor{warmgray}{RGB}{92,87,78}
\definecolor{softcoral}{RGB}{246,225,219}
\definecolor{linkblue}{RGB}{31,76,122}
\hypersetup{colorlinks=true,linkcolor=deepolive,citecolor=darkteal,urlcolor=linkblue,
 hypertexnames=false,
 pdftitle={Gaussian Coefficient Calculus: q-Binomial Kernels, Newton Bases, Bell Polynomials, Parameter Jets, Inversion, and Asymptotics},
 pdfauthor={Consolidated companion volume to Combinatorial Coefficient Calculus},
 pdfsubject={Gaussian coefficients, q-Stirling and q-Eulerian arrays, Bell polynomials, Jackson calculus, Lagrange inversion, root-of-unity jets, geometric filters, asymptotic regimes}}
\allowdisplaybreaks[3]
\raggedbottom
\emergencystretch=2em
\setlist{itemsep=2pt,topsep=4pt,parsep=1pt}
\setcounter{tocdepth}{1}
\setcounter{secnumdepth}{2}
\pagestyle{fancy}
\fancyhf{}
\fancyhead[L]{\small\nouppercase{\leftmark}}
\fancyhead[R]{\small\thepage}
\renewcommand{\headrulewidth}{0.4pt}
\renewcommand{\chaptermark}[1]{\markboth{\thechapter.\ #1}{}}
\fancypagestyle{plain}{\fancyhf{}\fancyfoot[C]{\thepage}\renewcommand{\headrulewidth}{0pt}}
\newtheorem{theorem}{Theorem}[chapter]
\newtheorem{lemma}[theorem]{Lemma}
\newtheorem{proposition}[theorem]{Proposition}
\newtheorem{corollary}[theorem]{Corollary}
\theoremstyle{definition}
\newtheorem{definition}[theorem]{Definition}
\newtheorem{example}[theorem]{Example}
\theoremstyle{remark}
\newtheorem{remark}[theorem]{Remark}
\newtheorem{warning}[theorem]{Warning}
\newtcolorbox{masterbox}[1][]{enhanced,breakable,colback=softolive,colframe=deepolive,
 boxrule=0.7pt,arc=1mm,left=2mm,right=2mm,top=1.5mm,bottom=1.5mm,
 title={Master principle},fonttitle=\bfseries,#1}
\newtcolorbox{auditbox}[1][]{enhanced,breakable,colback=softcoral,colframe=warmgray,
 boxrule=0.6pt,arc=1mm,left=2mm,right=2mm,top=1.5mm,bottom=1.5mm,
 title={Normalization and domain},fonttitle=\bfseries,#1}
% ---- primary notation (the spine source's conventions) ----
\newcommand{\N}{\mathbb N}
\newcommand{\Z}{\mathbb Z}
\newcommand{\Q}{\mathbb Q}
\newcommand{\C}{\mathbb C}
\newcommand{\R}{\mathbb R}
\newcommand{\F}{\mathbb F}
\newcommand{\GB}[3][q]{\begin{bmatrix}#2\\#3\end{bmatrix}_{#1}}
\newcommand{\qi}[2][q]{[#2]_{#1}}
\newcommand{\qf}[2][q]{[#2]_{#1}!}
\newcommand{\Y}[1]{\mathcal B_{#1}}
\newcommand{\B}[2]{\mathcal B_{#1,#2}}
\newcommand{\qB}[2]{\mathfrak B^{[q]}_{#1,#2}}
\newcommand{\coeff}[2]{[#1^{#2}]}
\newcommand{\fall}[2]{#1^{\underline{#2}}}
\newcommand{\rising}[2]{#1^{\overline{#2}}}
\DeclareMathOperator{\Li}{Li}
\DeclareMathOperator{\Var}{Var}
\DeclareMathOperator{\ord}{ord}
\DeclareMathOperator{\Res}{Res}
\DeclareMathOperator{\diag}{diag}
\DeclareMathOperator{\Id}{Id}
\newcommand{\dd}{\,\mathrm d}
\newcommand{\E}{\mathbb E}
\newcommand{\qbname}{\texorpdfstring{$q$}{q}}
\newcommand{\logq}{\vartheta}
% ---- aliases for the second source's macros (all render in the primary notation) ----
\newcommand{\qb}[2]{\GB{#1}{#2}}
\newcommand{\qbb}[3]{\GB[#3]{#1}{#2}}
\newcommand{\poch}[3]{(#1;#2)_{#3}}
\newcommand{\coef}[2]{[#1^{#2}]}
\newcommand{\Ecal}{\mathcal E}
\newcommand{\Pcal}{\mathcal P}
\newcommand{\Bhat}{\widehat B}
\newcommand{\Bq}{\mathfrak B^{[q]}}
\newcommand{\ee}{\mathrm e}
\newcommand{\ii}{\mathrm i}
\newcommand{\stir}[2]{\genfrac{\{}{\}}{0pt}{}{#1}{#2}}
\newcommand{\sone}[2]{s(#1,#2)}
% ---- aliases for the third source's macros ----
\newcommand{\G}[2]{\GB{#1}{#2}}
\newcommand{\Gb}[3]{\GB[#3]{#1}{#2}}
\newcommand{\qnum}[1]{[#1]_q}
\newcommand{\qfac}[1]{[#1]_q!}
\newcommand{\qp}[2]{(#1;q)_{#2}}
\newcommand{\ct}{\operatorname{CT}}
\newcommand{\Dq}{D_q}
\newcommand{\des}{\operatorname{des}}
\newcommand{\maj}{\operatorname{maj}}
\newcommand{\inv}{\operatorname{inv}}
\newcommand{\Sg}[2]{\left\{\begin{matrix}#1\\#2\end{matrix}\right\}}
\newcommand{\id}{\mathrm{id}}
\newcommand{\one}{\mathbf 1}
\newcommand{\Qn}{Q_n}
\newcommand{\Bcal}{\mathcal B}
\newcommand{\Yc}{\mathcal B}
\newcommand{\Poch}{\mathcal P}
\title{\Huge\bfseries Gaussian Coefficient Calculus\\[4mm]
 \Large \(q\)-Binomial Kernels, Newton Bases, Bell Polynomials,\\
 Parameter Jets, Inversion, and Asymptotics\\[7mm]
 \large A consolidated companion volume to \emph{Combinatorial Coefficient Calculus}}
\author{Consolidated from three research reports with detailed proofs}
\date{September 4, 2026}

% ---- hand chunk 01-front.tex ----
\begin{document}
\frontmatter
\hypersetup{pageanchor=false}
\maketitle
\hypersetup{pageanchor=true}
\chapter*{Abstract}
\addcontentsline{toc}{chapter}{Abstract}
This volume extends the basis-change and coefficient-extraction framework of
\emph{Combinatorial Coefficient Calculus} to Gaussian, or \(q\)-binomial,
coefficients. It consolidates three independently written reports on the
same subject into one document with a single notation, one statement of each
shared result, and complete proofs. The development is organized around
mathematical operations rather than a list of identities. Finite geometric
products produce Gaussian kernels; logarithms of those products produce
geometric power sums; ordinary Bell polynomials recover the coefficients;
Newton interpolation on arbitrary nodes produces Gaussian, \(q\)-Stirling, and
colored-node connection arrays; and ordinary formal composition and Lagrange
inversion transport the resulting formulae to implicit series.

The shared core comprises the finite and reciprocal \(q\)-binomial theorems,
Vandermonde and Chu--Vandermonde sums, Gaussian inversion with its
subspace-lattice interpretation, the ramified product-to-Bell master formula,
divisor-sum coefficients of Gaussian polynomials, Jackson differentiation and
geometric interpolation, two explicitly distinguished \(q\)-Stirling
normalizations, a weighted Bell--Riordan calculus, Lagrange inversion with
finite-product input, all-order expansions at \(q=1\), a uniform deflation
procedure at every root of unity, and Gaussian extrapolation filters on
geometric grids. Around that core the volume develops the material that only
one of the sources contained: a general node calculus with an affine-geometric
master connection, colored and type-\(B\) \(q\)-Stirling arrays, a
\(q\)-Dobinski identity, \(q\)-Worpitzky and Stirling--Eulerian transforms, a
path form of the divided-difference chain rule, weighted Gaussian transforms
and the failure of ordinary translation, the area-type \(q\)-Catalan
functional equation, Good-determinant entries for coupled inverses, Dilcher
divisor identities with their harmonic limit, positive linearization, weighted
and differentiated finite sums, multinomial factorial-ratio criteria, three
separated asymptotic regimes including a uniform double-scaling expansion, and
the exact recovery of geometric-uniform convolution moments, which is the
Fabius--Rvachev application.

Every theorem is proved or reduced explicitly to a proved result. Classical
identities are identified as such; no priority is claimed for the synthesis.
The three exact verification programs shipped with the sources are retained
unchanged and are described in the closing chapter.

\chapter*{Relation to the canonical manuscript and to the sources}
\addcontentsline{toc}{chapter}{Relation to the canonical manuscript and to the sources}
The canonical manuscript \cite{ccc} develops two central mechanisms: changing
to a basis adapted to an operator, and using Bell polynomials to recover
coefficients from component data. It separates formal from analytic
statements and keeps Bernoulli numbers distinct from Bell polynomials. This
volume retains that organization and those distinctions.

Its first point of attachment is the manuscript's Newton expansion: the
additive nodes \(0,1,2,\ldots\) are replaced either by geometric nodes
\(a,aq,aq^2,\ldots\), by \(q\)-integer nodes \([0]_q,[1]_q,[2]_q,\ldots\), or
by an arbitrary node list, for which the connection coefficients are
complete and elementary symmetric functions of the nodes. Its second point of
attachment is the exponential formula: a finite geometric alphabet has power
sums \([N]_{q^r}\), so the ordinary Bell polynomials already used in the
manuscript encode \(q\)-binomial identities without any \(q\)-Bell
polynomial. Its third point of attachment is ordinary composition and
reversion: a change to \(q\)-factorial coefficient normalization changes the
matrices, not the meaning of substitution.

\paragraph{The three sources.}
This volume replaces three reports filed together on September 4, 2026:
a book-length treatment organized by operations (Gaussian kernels, Bell
polynomials, basis changes, parameter jets, inversion, and geometric
filters); an article emphasizing the ramified master kernel, weighted
transforms, the area-type \(q\)-Catalan equation, coupled inverses, and a
normalization audit; and an article emphasizing the general node calculus,
\(q\)-Eulerian polynomials, the divided-difference chain rule, asymptotic
regimes, and the geometric-uniform moment application. Where two or three of
them proved the same theorem, one statement and one proof are kept, chosen
for completeness, and the others' variant formulations are recorded as
remarks or corollaries when they add a formula. Where their normalizations
differed, the definitions in \cref{q1:qcc:setting} fix the convention, and
the provenance chapter records the dictionary between the sources' symbols.
Passages that were written by one source and kept verbatim keep their
internal cross-references; the labels of shared results resolve to the single
retained statement.

This is a standalone mathematical extension, not a replacement of the
canonical manuscript and not a claim that its results have been formalized
in Lean. It loads the corpus notation file \path{docs/fabius-notation.tex}
for consistency with the other volumes, and defines the short macros it uses
itself. The typography follows the manuscript's Libertinus, olive, and coral
conventions.

\begin{auditbox}[title={What ``an extension'' means here}]
The finite \(q\)-binomial theorem, Gaussian inversion, \(q\)-Chu--Vandermonde,
Carlitz-type \(q\)-Stirling arrays, \(q\)-Worpitzky, divisor identities, and
\(q\)-Lucas are established mathematics. They are rederived to connect their
hypotheses and normalizations. The main additional value is the common
coefficient machinery: explicit all-order root-of-unity jets, full
geometric-filter error formulae, the affine-geometric node connection, the
path chain rule, and the uniform double-scaling expansion. No exhaustive
novelty claim is made.
\end{auditbox}
\begingroup
\small
\tableofcontents
\endgroup

\mainmatter

% ---- Q1 lines 153--226 ----
\chapter{The coefficient-calculus setting}
\label{q1:qcc:setting}
\section{Rings, indices, and specialization}
Unless stated otherwise, \(q\) is an indeterminate and calculations with
quotients take place over \(\Q(q)\), extended by the finitely many auxiliary
parameters that occur. A formal series in \(z\) is an element of \(K[[z]]\) for a
commutative \(\Q\)-algebra \(K\). Composition \(A(B(z))\) is defined when
\(B(0)=0\); inversion under multiplication requires an invertible constant term.
These statements have no convergence content.

All integer indices are nonnegative unless explicitly qualified. Empty products
are one, sums over empty ranges are zero, and \(0^0=1\) in coefficient formulae.
Gaussian coefficients outside \(0\le k\le n\) are zero. Define
\begin{equation}\label{q1:qcc:notation}
 (a;q)_N=\prod_{j=0}^{N-1}(1-aq^j),\qquad
 \qi n=\sum_{j=0}^{n-1}q^j,\qquad
 \qf n=\prod_{j=1}^n\qi j.
\end{equation}
In particular, \([0]_q=0\) and \([0]_q!=1\). The polynomial interpretation of
\([n]_q\) makes \([n]_1=n\) immediate. The quotient
\((1-q^n)/(1-q)\) is a useful representation, not the definition at \(q=1\).
Our \(q\)-shifted factorial convention agrees with \cite{dlmf172}.

When a rational representation equals a polynomial, specialization means
specializing that polynomial, after cancellation. In particular, the expression
\((q;q)_n/((q;q)_k(q;q)_{n-k})\) must not be evaluated as an uncancelled
\(0/0\) at a root of unity. Later we give explicit polynomial and local-series
methods that do not make this mistake.

\section{Bell and Bernoulli normalizations}
The complete and partial exponential Bell polynomials are characterized by
\begin{align}
 \exp\left(\sum_{r\ge1}x_r\frac{z^r}{r!}\right)
 &=\sum_{n\ge0}\Y n(x_1,\ldots,x_n)\frac{z^n}{n!},\label{q1:qcc:bell-def}\\
 \frac1{k!}\left(\sum_{r\ge1}x_r\frac{z^r}{r!}\right)^k
 &=\sum_{n\ge k}\B nk(x_1,\ldots)\frac{z^n}{n!}.\label{q1:qcc:partial-def}
\end{align}
Thus \(\Y0=1\), and
\begin{equation}\label{q1:qcc:bell-finite}
 \Y n(x_1,\ldots,x_n)
 =n!\!\sum_{\substack{m_1,\ldots,m_n\ge0\\\sum r m_r=n}}
 \prod_{r=1}^n\frac1{m_r!}\left(\frac{x_r}{r!}\right)^{m_r}.
\end{equation}
For \(\B nk\), impose also \(\sum m_r=k\). Expanding each exponential in
\eqref{q1:qcc:bell-def} proves \eqref{q1:qcc:bell-finite}. In particular, a formula in
Bell polynomials is a finite-sum formula, not a hidden recurrence.

Bernoulli numbers are denoted \(\beta_r\), with
\begin{equation}\label{q1:qcc:bernoulli-def}
 \frac{t}{e^t-1}=\sum_{r\ge0}\beta_r\frac{t^r}{r!},
 \qquad \beta_1=-\frac12.
\end{equation}
We use \(s(n,k)\) for signed ordinary Stirling numbers of the first kind and
\(S(n,k)\) for ordinary Stirling numbers of the second kind. Their \(q\)-versions
will be defined independently, with the normalization stated in the definition.

\section{Three distinct operations}
There are three operations that should not be conflated. One may specialize the
variables of a symmetric polynomial to a geometric progression. One may replace
an ordinary derivative by the Jackson derivative. Or one may merely record
ordinary series coefficients using \([n]_q!\) rather than \(n!\). All three lead
to Gaussian coefficients, but their product and composition rules are not
interchangeable.

\begin{masterbox}
The robust route is
\[
 \text{finite alphabet}\ \longrightarrow\ \text{power sums}
 \ \longrightarrow\ \text{formal logarithm}
 \ \longrightarrow\ \text{Bell coefficients}.
\]
Operator-adapted interpolation and formal reversion can then be attached to this
route without guessing which ordinary factorials should acquire a subscript \(q\).
\end{masterbox}
% ---- hand chunk 02-kernel-head.tex ----

\section{The partition-profile kernel and additional notation}
\label{q2:sec:kernel}
Several passages below, written originally in a slightly different
notation, use the following abbreviations. We write
\[
 D_n=(q;q)_n=(1-q)^n[n]_q!,\qquad
 G_{n,k}(q)=\GB nk
\]
when the Gaussian polynomial is to be viewed as a function of its base, and,
for a scalar \(Q\), \(G_{n,k}(Q)\) for its value at \(q=Q\). A finite
product of shifted factorials is often expanded through the following
partition-profile kernel, which is the complete Bell polynomial in a
different normalization.


% ---- Q2 lines 156--164 ----
For $n\ge0$ let
\[
 \Pcal(n)=\{(m_1,\ldots,m_n)\in\N^n:\ \sum_{j=1}^n jm_j=n\}.
\]
For $n=0$ this consists of the empty tuple. Define
\begin{equation}\label{q2:eq:E-kernel}
 \Ecal_n(L)=\sum_{\bm m\in\Pcal(n)}
       \prod_{j=1}^n\frac{(L_j/j)^{m_j}}{m_j!}.
\end{equation}
% ---- hand chunk 02b-kernel-glue.tex ----

In terms of the complete exponential Bell polynomial \(\Y n\) of
\eqref{q1:qcc:bell-def},

% ---- Q2 lines 170--187 ----
Consequently
\begin{equation}\label{q2:eq:E-Bell}
 \Ecal_n(L)=\frac1{n!}\Y n(0!L_1,1!L_2,\ldots,(n-1)!L_n),\qquad
 \exp\left(\sum_{j\ge1}\frac{L_j}jz^j\right)
       =\sum_{n\ge0}\Ecal_n(L)z^n.
\end{equation}
Both statements follow by multiplying the individual exponential series. In
particular, \eqref{q2:eq:E-kernel} is a definition by a finite sum, not by a recurrence.
The first values provide a useful normalization check:
\begin{align}\label{q2:eq:E-small}
 \Ecal_0&=1,&\Ecal_1&=L_1,&
 \Ecal_2&=\frac{L_1^2+L_2}{2},\\
 \Ecal_3&=\frac{L_1^3+3L_1L_2+2L_3}{6},&
 \Ecal_4&=\frac{L_1^4+6L_1^2L_2+3L_2^2+8L_1L_3+6L_4}{24}.
 \nonumber
\end{align}
The $q$-dependence will usually enter through the arguments $L_j$, not through a
new Bell-polynomial definition.
% ---- Q1 lines 228--286 ----
\chapter{Gaussian polynomials and finite products}
\label{q1:qcc:gaussian}
\section{A division-free definition}
\begin{definition}
For \(0\le k\le n\), set
\begin{equation}\label{q1:qcc:subset}
 \GB nk=\sum_{0\le i_1<\cdots<i_k<n}
 q^{i_1+\cdots+i_k-\binom{k}{2}}.
\end{equation}
The \(k=0\) summand is the empty subset and contributes one.
\end{definition}
Every exponent is nonnegative, because \(i_j\ge j-1\). Consequently this is a
polynomial in \(\Z_{\ge0}[q]\), valid at every complex \(q\).

\begin{theorem}[Pascal, factorial, and symmetry formulae]\label{q1:qcc:pascal}
For \(n\ge1\), the Gaussian polynomials satisfy the first two recurrences
below; the remaining identities hold for \(0\le k\le n\):
\begin{align}
 \GB nk&=\GB{n-1}k+q^{n-k}\GB{n-1}{k-1},\label{q1:qcc:pascal1}\\
 \GB nk&=q^k\GB{n-1}k+\GB{n-1}{k-1},\label{q1:qcc:pascal2}\\
 \GB nk&=\frac{(q;q)_n}{(q;q)_k(q;q)_{n-k}}
       =\frac{\qf n}{\qf k\qf{n-k}},\label{q1:qcc:quotient}\\
 \GB nk&=\GB n{n-k},\qquad
 \GB[q^{-1}]nk=q^{-k(n-k)}\GB nk.\label{q1:qcc:reciprocity}
\end{align}
The quotient identities hold in \(\Q(q)\). For \(0\le k\le n\), the polynomial
has degree \(k(n-k)\), constant coefficient one, and leading coefficient one.
\end{theorem}
\begin{proof}
Separate the subsets in \eqref{q1:qcc:subset} according as they contain \(n-1\).
Removing that element decreases the exponent by \(n-k\), proving
\eqref{q1:qcc:pascal1}. The factorial quotient has the same boundary values and
satisfies the same recurrence: after putting the two terms over a common
denominator the remaining numerator is
\((1-q^{n-k})+q^{n-k}(1-q^k)=1-q^n\).
Induction therefore proves \eqref{q1:qcc:quotient}. Its symmetry gives
\eqref{q1:qcc:pascal2} by replacing \(k\) with \(n-k\).

The smallest subset \(\{0,\ldots,k-1\}\) has exponent zero, and the largest
\(\{n-k,\ldots,n-1\}\) has exponent \(k(n-k)\). Each is unique.
Reflecting indices by \(i_j\mapsto n-1-i_{k+1-j}\) replaces an exponent \(d\)
by \(k(n-k)-d\), proving reciprocity and palindromicity.
\end{proof}

A subset gives a partition inside a \(k\)-by-\((n-k)\) rectangle by taking
\(\lambda_j=i_{k+1-j}-(k-j)\). These integers decrease weakly, lie between
zero and \(n-k\), and sum to the exponent in \eqref{q1:qcc:subset}. The construction
is reversible. This proves the rectangular-partition interpretation directly.
At \(q=1\), every subset has weight one, so \(\GB[1]nk=\binom nk\).

\begin{example}
\[
 \GB42=1+q+2q^2+q^3+q^4,\qquad
 \GB63=1+q+2q^2+3q^3+3q^4+3q^5+3q^6+2q^7+q^8+q^9.
\]
The coefficient sums are respectively \(6\) and \(20\); the symmetry centers
are respectively \(2\) and \(9/2\).
\end{example}

% ---- hand chunk 03-models-head.tex ----

\section{Rectangle partitions, words, and flags}
The subset sum in \cref{q1:qcc:subset} has two further finite models, and
the multinomial generalization has a finite-field model. They are used
later for positivity arguments and for the \(q\)-Eulerian transforms.


% ---- Q3 lines 230--262 ----
\begin{theorem}[Three polynomial models]
\label{q3:thm:subset}
For $0\le k\le n$,
\begin{align}
 \G nk
 &=\sum_{0\le i_1<\cdots<i_k<n}
       q^{i_1+\cdots+i_k-\binom k2},
 \label{q3:eq:subset}\\
 &=\sum_{0\le b_1\le\cdots\le b_k\le n-k}q^{b_1+\cdots+b_k}.
 \label{q3:eq:rectangle}
\end{align}
Equivalently, it enumerates partitions fitting in a $k\times(n-k)$
rectangle by their size. It also enumerates binary words with $k$ ones
and $n-k$ zeros by the number of pairs consisting of a zero to the left
of a one.
\end{theorem}
\begin{proof}
Let the first sum be $H_{n,k}$. Subsets not containing $n-1$ contribute
$H_{n-1,k}$; deleting $n-1$ from the others decreases the exponent by
$(n-1)-(k-1)=n-k$. Thus $H$ satisfies \cref{q1:qcc:pascal2} with the same
boundary values as $G$. The substitution $b_j=i_j-(j-1)$ gives the second
sum. Reversing the $b_j$ gives a partition inside the rectangle.
In the associated binary word, the $j$th one has exactly
$i_j-(j-1)$ zeros before it. Summing gives the stated word statistic.
\end{proof}
Reversal of the word changes that statistic to $k(n-k)$ minus itself,
which explains reciprocity without division. For example,
\begin{equation}
 \G42=1+q+2q^2+q^3+q^4.
 \label{q3:eq:G42}
\end{equation}
The two objects of weight $2$ are the rectangle partitions $(2,0)$ and
$(1,1)$.
% ---- Q3 lines 264--293 ----
\section{Multinomials, finite fields, and flags}
For $n_1+\cdots+n_r=n$, write
\begin{equation}
 \Gb n{n_1,\ldots,n_r}q
 =\frac{\qfac n}{\prod_{i=1}^r\qfac{n_i}}
 =\prod_{i=1}^r\G{n-\sum_{j<i}n_j}{n_i}.
 \label{q3:eq:qmultinomial}
\end{equation}
The right side proves polynomiality. Assigning the positions of letter
$1$, then letter $2$, and so on shows that it enumerates words with these
multiplicities by the usual inversion statistic; reversal also gives the
complementary convention.

\begin{proposition}[Subspace and flag counts]
\label{q3:prop:flags}
If $Q$ is a prime power, $G_{n,k}(Q)$ counts the $k$-dimensional
subspaces of $\mathbb F_Q^n$. The multinomial in
\cref{q3:eq:qmultinomial}, evaluated at $q=Q$, counts flags whose successive
dimension increments are $n_1,\ldots,n_r$.
\end{proposition}
\begin{proof}
There are $(Q^n-1)(Q^n-Q)\cdots(Q^n-Q^{k-1})$ ordered independent
$k$-tuples in $\mathbb F_Q^n$. Each $k$-subspace has
$(Q^k-1)(Q^k-Q)\cdots(Q^k-Q^{k-1})$ ordered bases. Dividing and canceling
the powers of $Q$ gives \cref{q1:qcc:quotient}. Choose each next subspace
in the quotient by the preceding one to obtain the flag product.
\end{proof}
The indeterminate $q$ is not itself a field size. The finite-field
interpretation applies only at prime powers; the polynomial identities
hold independently of it.
% ---- Q1 lines 287--315 ----
\section{The finite and reciprocal binomial theorems}
\begin{theorem}\label{q1:qcc:finite-binomial}
For every \(N\ge0\),
\begin{equation}\label{q1:qcc:finite}
 (z;q)_N=\sum_{k=0}^N(-1)^kq^{\binom{k}{2}}\GB Nk z^k.
\end{equation}
For \(N\ge1\), as a formal series in \(z\),
\begin{equation}\label{q1:qcc:reciprocal}
 \frac1{(z;q)_N}=\sum_{r\ge0}\GB{N+r-1}r z^r.
\end{equation}
When \(N=0\), the reciprocal is one.
\end{theorem}
\begin{proof}
To obtain \(z^k\) in the finite product, choose the \(-zq^i\) term in precisely
\(k\) factors. Its weight is \((-1)^kq^{\sum i}\); comparison with
\eqref{q1:qcc:subset} proves \eqref{q1:qcc:finite}.

Write \(H_N(z)=1/(z;q)_N\). Then
\((1-q^{N-1}z)H_N(z)=H_{N-1}(z)\).
If \(c_{N,r}=[z^r]H_N\), it follows that
\(c_{N,r}=c_{N-1,r}+q^{N-1}c_{N,r-1}\).
The candidate \(\GB{N+r-1}r\) has this recurrence by
\eqref{q1:qcc:pascal1}, with \(c_{1,r}=1\) and \(c_{N,0}=1\).
Induction on \(N+r\) proves the reciprocal identity.
\end{proof}

These classical identities are the finite and negative-power backbone of the
extension; compare the standard conventions in \cite{dlmf172,dlmf175}.

% ---- hand chunk 04-negative-head.tex ----

\section{A negative upper index}
\label{q2:sec:negative}
The reciprocal identity \eqref{q1:qcc:reciprocal} suggests extending the
upper index downward. Only the upper index is extended here; the lower index
stays nonnegative.


% ---- Q2 lines 269--280 ----
For $M\ge1$ and $r\ge0$ define a negative upper entry by the finite product
\[
 \qb{-M}r:=\prod_{j=0}^{r-1}\frac{1-q^{-M-j}}{1-q^{j+1}}.
\]
Factoring each numerator gives the explicit conversion
\begin{equation}\label{q2:eq:negative-upper}
 \qb{-M}r=(-1)^r q^{-Mr-\binom r2}\qb{M+r-1}r.
\end{equation}
Only the upper index has been extended here. This avoids the competing
conventions that arise when both indices are negative; a broader treatment is
available in \cite{formichella}.

% ---- Q1 lines 316--349 ----
\section{A formal infinite-binomial identity without a topology error}
\begin{theorem}\label{q1:qcc:infinite-binomial}
There is a unique \(R_a(z)\in\Q(q,a)[[z]]\) with \(R_a(0)=1\) satisfying
\begin{equation}\label{q1:qcc:binomial-functional}
 (1-z)R_a(z)=(1-az)R_a(qz).
\end{equation}
It is
\begin{equation}\label{q1:qcc:formal-binomial-series}
 R_a(z)=\sum_{r\ge0}\frac{(a;q)_r}{(q;q)_r}z^r.
\end{equation}
For complex \(|q|<1\) and \(|z|<1\), it also equals
\begin{equation}\label{q1:qcc:analytic-binomial}
 R_a(z)=\frac{(az;q)_\infty}{(z;q)_\infty}.
\end{equation}
\end{theorem}
\begin{proof}
Comparison of coefficients in \eqref{q1:qcc:binomial-functional} gives
\((1-q^r)c_r=(1-aq^{r-1})c_{r-1}\), with \(c_0=1\).
This proves both uniqueness and \eqref{q1:qcc:formal-binomial-series}.
For \(|q|<1\), the infinite products converge locally uniformly in their
arguments, since the tails are bounded by convergent geometric sums. The
denominator has no zero for \(|z|<1\). Their ratio satisfies the same functional
equation, is analytic at zero, and has constant term one. Its Taylor coefficients
are therefore the coefficients already computed. Unless the series terminates, the coefficient ratio tends
to one, so the series converges absolutely for \(|z|<1\); termination
makes convergence immediate.
\end{proof}
\begin{auditbox}
With \(q\) merely an indeterminate, the product
\(\prod_{j\ge0}(1-zq^j)\) is not defined by \(z\)-adic convergence: its linear
coefficient changes at every step. Equation \eqref{q1:qcc:formal-binomial-series}
or the functional equation is the formal definition. Equation
\eqref{q1:qcc:analytic-binomial} is an additional analytic realization.
\end{auditbox}
% ---- Q1 lines 351--383 ----
\chapter{Convolution, inversion, and linearization}
\label{q1:qcc:convolution}
\section{Two-block and many-block Vandermonde identities}
\begin{theorem}\label{q1:qcc:vandermonde}
For \(m,n,r\ge0\),
\begin{equation}\label{q1:qcc:vandermonde-formula}
 \GB{n+m}r=
 \sum_{k=\max(0,r-m)}^{\min(n,r)}
 q^{(n-k)(r-k)}\GB nk\GB m{r-k}.
\end{equation}
More generally, if \(n_1+\cdots+n_s=N\), then
\begin{equation}\label{q1:qcc:block-vandermonde}
 \GB Nr=\sum_{\substack{0\le k_i\le n_i\\\sum k_i=r}}
 q^{\sum_{i<j}(n_i-k_i)k_j}\prod_{i=1}^s\GB{n_i}{k_i}.
\end{equation}
\end{theorem}
\begin{proof}
Split \((-z;q)_{n+m}=(-z;q)_n(-zq^n;q)_m\) and use
\eqref{q1:qcc:finite} with \(-z\) in place of \(z\).
After dividing the coefficient of \(z^r\) by \(q^{\binom r2}\), the exponent is
\[
 \binom k2+\binom{r-k}2+n(r-k)-\binom r2=(n-k)(r-k).
\]
The same computation for consecutive blocks of lengths \(n_i\) gives
\(\sum_i\binom{k_i}2+\sum_{i<j}n_i k_j-\binom r2
 =\sum_{i<j}(n_i-k_i)k_j\), proving the general case.
\end{proof}
In particular, symmetry and the substitution \(k\mapsto n-k\) give
\begin{equation}\label{q1:qcc:central}
 \sum_{k=0}^n q^{k^2}\GB nk^2=\GB{2n}n.
\end{equation}
The exponent \(k^2\) is essential; removing it does not give a Gaussian analogue
of the ordinary central-binomial identity.
% ---- Q2 lines 327--349 ----
\begin{corollary}[Three useful companions]\label{q2:cor:companions}
For $M,N\ge1$ and $r\ge0$,
\begin{align}
 \qb{M+N+r-1}r
 &=\sum_{k=0}^r q^{M(r-k)}\qb{M+k-1}k\qb{N+r-k-1}{r-k},
 \label{q2:eq:negative-vander}\\
 \delta_{r0}
 &=\sum_{k=0}^{\min(N,r)}(-1)^kq^{\binom k2}
     \qb Nk\qb{N+r-k-1}{r-k}.\label{q2:eq:orth-product}
\end{align}
For $0\le k\le n$,
\begin{equation}\label{q2:eq:hockey}
 \sum_{j=k}^n q^{j-k}\qb jk=\qb{n+1}{k+1}.
\end{equation}
\end{corollary}
\begin{proof}
The first formula compares coefficients in
$1/\poch zq{M+N}=1/(\poch zqM\poch{zq^M}qN)$.
The second compares coefficients in
$\poch zqN/\poch zqN=1$.
For the third, rearrange the Pascal identity as
$q^{j-k}\qb jk=\qb{j+1}{k+1}-\qb j{k+1}$ and telescope.
\end{proof}
% ---- Q1 lines 385--420 ----
\section{Gaussian inversion and the subspace lattice}
\begin{theorem}[Gaussian binomial inversion]\label{q1:qcc:inversion}
For arbitrary sequences over \(\Q(q)\),
\begin{equation}\label{q1:qcc:inverse-pair}
 b_n=\sum_{k=0}^n\GB nk a_k
 \quad\Longleftrightarrow\quad
 a_n=\sum_{k=0}^n(-1)^{n-k}q^{\binom{n-k}{2}}\GB nk b_k.
\end{equation}
Both identities also make sense over \(\Z[q]\) when the sequences do.
\end{theorem}
\begin{proof}
The coefficient of \(a_j\) after substitution on the right is
\[
 \sum_{k=j}^n(-1)^{n-k}q^{\binom{n-k}{2}}\GB nk\GB kj.
\]
Cancel factorials to obtain
\(\GB nk\GB kj=\GB nj\GB{n-j}{k-j}\).
The remaining sum is \((1;q)_{n-j}\), after reversing its index and using
symmetry. It is zero unless \(n=j\), when it is one. The two triangular
matrices therefore multiply to the identity. Their unit diagonal makes either
one-sided inverse also the inverse on every finite truncation.
\end{proof}

The same inverse has a geometric interpretation. For a prime power \(Q\), the
number of \(k\)-dimensional subspaces of \(\F_Q^n\) is \(\GB[Q]nk\).
Indeed, count ordered linearly independent \(k\)-tuples and divide by the number
of ordered bases of \(\F_Q^k\):
\[
 \frac{\prod_{j=0}^{k-1}(Q^n-Q^j)}{\prod_{j=0}^{k-1}(Q^k-Q^j)}=\GB[Q]nk.
\]
An interval of rank difference \(r\) in the subspace lattice is itself a
subspace lattice of dimension \(r\). Its incidence-algebra M\"obius function is
\((-1)^rQ^{\binom r2}\), because its defining convolution with the constant-one
function is exactly \((1;Q)_r=0\) for \(r>0\). This gives a direct counterpart
of the canonical manuscript's use of incidence algebras for ordinary inversion.

% ---- Q3 lines 371--384 ----
\begin{corollary}[Gaussian annihilator]
\label{q3:cor:annihilator}
For an indeterminate $x$,
\begin{equation}
 \sum_{j=0}^N(-1)^jq^{\binom j2}\G Nj x^j=\qp xN.
 \label{q3:eq:annihilator}
\end{equation}
In particular the sum at $x=1$ is zero for $N>0$. If $q\ne0$ and
$x=q^{-r}$ with $0\le r<N$, it is again zero.
\end{corollary}
\begin{proof}
This is \cref{q1:qcc:finite}; in the two specializations one factor of the
product is zero.
\end{proof}
% ---- Q1 lines 421--505 ----
\section{A positive linearization formula}
\begin{theorem}\label{q1:qcc:linearization}
For \(0\le k,\ell\le n\),
\begin{equation}\label{q1:qcc:linearization-formula}
 \GB nk\GB n\ell
 =\sum_{r=\max(k,\ell)}^{\min(n,k+\ell)}
 q^{(r-k)(r-\ell)}\GB nr\GB rk\GB k{k+\ell-r}.
\end{equation}
\end{theorem}
\begin{proof}
First let \(q=Q\) be a prime power and count ordered pairs \((V,W)\) of
subspaces of dimensions \(k,\ell\). Set \(U=V+W\), \(\dim U=r\), and
\(J=V\cap W\), \(\dim J=j=k+\ell-r\). Choose \(U\), then \(V\subseteq U\),
then \(J\subseteq V\), giving the three Gaussian factors.
In \(U/J\), the space \(V/J\) has dimension \(k-j=r-\ell\).
The possible \(W/J\) are complements of this space, of dimension
\(\ell-j=r-k\). Relative to one fixed complement, every other complement is
the graph of a unique linear map into \(V/J\). There are
\(Q^{(r-k)(r-\ell)}\) such maps. This proves the formula for every prime power.
Both sides are polynomials in \(q\); infinitely many equal values prove the
polynomial identity.
\end{proof}
For example,
\begin{equation}\label{q1:qcc:qinteger-square}
 \qi n^2=\qi n+q(1+q)\GB n2.
\end{equation}
This is a Gaussian factorial-moment identity. A general power analogue follows
from \(q\)-Stirling basis change in \cref{q1:qcc:qstirling}.

\section{Pochhammer convolution and terminating hypergeometric sums}
\begin{theorem}\label{q1:qcc:poch-convolution}
For \(N\ge0\),
\begin{equation}\label{q1:qcc:poch-convolution-formula}
 (ab;q)_N=\sum_{j=0}^N\GB Nj a^j(b;q)_j(a;q)_{N-j}.
\end{equation}
Consequently, wherever the displayed denominators are nonzero,
\begin{align}
 \sum_{j=0}^N
 \frac{(q^{-N};q)_j(a;q)_j}{(q;q)_j(c;q)_j}q^j
 &=a^N\frac{(c/a;q)_N}{(c;q)_N},\label{q1:qcc:chu1}\\
 \sum_{j=0}^N
 \frac{(q^{-N};q)_j(a;q)_j}{(q;q)_j(c;q)_j}
 \left(\frac{cq^N}{a}\right)^j
 &=\frac{(c/a;q)_N}{(c;q)_N}.\label{q1:qcc:chu2}
\end{align}
The last two formulae are the terminating \(q\)-Chu--Vandermonde identities.
\end{theorem}
\begin{proof}
The functional equation \eqref{q1:qcc:binomial-functional} implies
\(R_{ab}(z)=R_a(z)R_b(az)\): both sides have constant term one and satisfy
\((1-z)R(z)=(1-abz)R(qz)\).
Comparison of \(z^N\) in \eqref{q1:qcc:formal-binomial-series}, followed by
multiplication by \((q;q)_N\), proves \eqref{q1:qcc:poch-convolution-formula}.
This proof is entirely formal.

Divide that formula by \((a;q)_N\). Direct reversal of finite products gives
\begin{equation}\label{q1:qcc:conversion-chu}
 \GB Nj\frac{(a;q)_{N-j}}{(a;q)_N}a^j
 =\frac{(q^{-N};q)_j}{(q;q)_j(q^{1-N}/a;q)_j}q^j.
\end{equation}
For completeness, use
\[
 (q^{-N};q)_j=(-1)^j q^{-Nj+\binom j2}\frac{(q;q)_N}{(q;q)_{N-j}}
\]
and
\[
 (aq^{N-j};q)_j=(-a)^j q^{Nj-j(j+1)/2}(q^{1-N}/a;q)_j
\]
in the left side. The powers reduce to \(q^j\), proving
\eqref{q1:qcc:conversion-chu}.
Now set the original \(a=q^{1-N}/c\) and rename the original \(b\) as \(a\).
The product reversal
\((x;q)_N=(-x)^Nq^{\binom N2}(q^{1-N}/x;q)_N\)
turns \((ab;q)_N/(a;q)_N\) into the right side of \eqref{q1:qcc:chu1}.

Finally replace \(q,a,c\) in \eqref{q1:qcc:chu1} by
\(q^{-1},a^{-1},c^{-1}\). Reversing each finite product changes the summand's
extra power into \((cq^N/a)^j\); the prefactor \(a^{-N}\) on the right cancels
the reversal factor \(a^N\). This yields \eqref{q1:qcc:chu2}.
All calculations first take place in a rational-function field; any admissible
specialization then follows.
\end{proof}
The identities are classical; \cite{dlmf176} records their hypergeometric
normalizations. The proof above explains why they are coefficient-convolution
identities before they are hypergeometric summations.
% ---- hand chunk 05-transform-head.tex ----

\section{Divided-power coordinates and weighted Gaussian transforms}
\label{q2:sec:transforms}
Gaussian inversion, \cref{q1:qcc:inversion}, is the \(c=1\) case of a
one-parameter family of transforms. The family is best seen through the two
\(q\)-exponentials, which also govern the \(q\)-factorial coordinates used
for composition in \cref{q1:qcc:composition}.


% ---- Q2 lines 596--625 ----
Define
\begin{equation}\label{q2:eq:qexp}
 e_q(z)=\sum_{n\ge0}\frac{z^n}{\qf n},\qquad
 E_q(z)=\sum_{n\ge0}q^{\binom n2}\frac{z^n}{\qf n}.
\end{equation}
Using \cref{q1:qcc:infinite-binomial} with $a=0$ and the finite theorem followed by its
coefficientwise limit gives
\begin{equation}\label{q2:eq:qexp-products}
 e_q(z)=\frac1{\poch{(1-q)z}q\infty},\qquad
 E_q(z)=\poch{-(1-q)z}q\infty,
 \qquad e_q(z)E_q(-z)=1.
\end{equation}
For completeness, the coefficientwise limit used for the second equality is
$q^{\binom n2}/D_n$: in the finite theorem,
$\qb Nn=\prod_{j=0}^{n-1}(1-q^{N-j})/D_n$ tends $q$-adically to $1/D_n$.
Thus \eqref{q2:eq:qexp-products} also holds purely formally over $\Q(q)$.
The functions in \eqref{q2:eq:qexp} use $q$-factorials; a convention with $D_n$ in
the denominator rescales their argument by $1-q$.

For sequences $a,b$ put
\begin{equation}\label{q2:eq:qconv}
 (a\star_q b)_n=\sum_{k=0}^n\qb nk a_kb_{n-k},\qquad
 \mathscr G_q(a;z)=\sum_{n\ge0}a_n\frac{z^n}{\qf n}.
\end{equation}
Multiplication of formal series immediately proves
\begin{equation}\label{q2:eq:qconv-series}
 \mathscr G_q(a\star_q b;z)=\mathscr G_q(a;z)\mathscr G_q(b;z).
\end{equation}
Hence $\star_q$ is associative and commutative and has identity
$\epsilon_n=\delta_{n0}$. Alternatively, associativity follows directly from
% ---- Q2 lines 629--671 ----
\begin{theorem}[Weighted Gaussian inversion]\label{q2:thm:transform}
Let $c$ be a scalar. The two transforms
\begin{align}
 b_n&=\sum_{k=0}^n\qb nk c^{n-k}a_k,\label{q2:eq:transform-forward}\\
 a_n&=\sum_{k=0}^n(-1)^{n-k}q^{\binom{n-k}2}\qb nk c^{n-k}b_k
 \label{q2:eq:transform-inverse}
\end{align}
are mutually inverse. These are polynomial identities, valid over every
commutative ring after assigning a value to $q$ and $c$.
\end{theorem}
\begin{proof}
Over $\Q(q)$, the forward transform multiplies the generating series by
$e_q(cz)$ and the inverse multiplies by $E_q(-cz)$. Their product is $1$.
For a denominator-free proof, the coefficient of $a_j$ after both transforms is
\[
 c^{n-j}\qb nj
 \sum_{r=0}^{n-j}(-1)^r q^{\binom r2}\qb{n-j}r
 =c^{n-j}\qb nj\poch1q{n-j}.
\]
This is $1$ when $n=j$ and $0$ otherwise. Interchanging the transforms gives the
same sum, proving both inverse compositions. Since the entire argument is a
polynomial identity with integer coefficients, arbitrary specialization follows.
\end{proof}

\section{Why ordinary translation fails}
Let $T_c$ denote \eqref{q2:eq:transform-forward}. Define the homogeneous polynomial
\[
 H_r(c,d;q)=\sum_{j=0}^r\qb rj c^j d^{r-j}.
\]
Associativity of convolution gives the exact composition law
\begin{equation}\label{q2:eq:transform-composition}
 (T_cT_d a)_n=\sum_{k=0}^n\qb nk H_{n-k}(c,d;q)a_k.
\end{equation}
For example,
\[
 H_2(c,d;q)=c^2+(1+q)cd+d^2,
\]
which is not $(c+d)^2$ for generic commuting $c,d$. Thus
$T_cT_d\ne T_{c+d}$ in general, and $T_{-c}$ is not the inverse of $T_c$.
The missing factor in the naive inverse is precisely $q^{\binom{n-k}2}$.
This contrasts with the ordinary weighted-translation theorem in the canonical manuscript.
The correct replacement is either \eqref{q2:eq:transform-composition}, or the
multiplicative cocycle of \cref{q1:qcc:poch-convolution} when shifted-factorial kernels are used.
% ---- Q2 lines 943--957 ----
There is, however, a genuine Gaussian binomial theorem for noncommuting variables.
\begin{proposition}[Quantum-plane binomial theorem]\label{q2:prop:quantum-plane}
In an associative algebra with a central parameter $q$ and $YX=qXY$,
\begin{equation}\label{q2:eq:quantum-plane}
 (X+Y)^n=\sum_{k=0}^n\qb nk X^{n-k}Y^k.
\end{equation}
\end{proposition}
\begin{proof}
Right-multiply by $X+Y$ and use $Y^kX=q^kXY^k$.
The new coefficient is $q^k\qb nk+\qb n{k-1}=\qb{n+1}k$.
The initial case is $n=0$.
\end{proof}
This theorem is consistent with the failure of ordinary scalar translation in
\cref{q2:sec:transforms}: the Gaussian cross term is produced by a commutation
relation, not by the ordinary expansion of commuting $c+d$.
% ---- Q1 lines 507--617 ----
\chapter{Symmetric functions and the Bell master formula}
\label{q1:qcc:bell-master}
\section{Elementary, complete, and power-sum coordinates}
For a finite alphabet \(X=(x_1,\ldots,x_N)\), define
\[
 e_r(X)=\sum_{i_1<\cdots<i_r}x_{i_1}\cdots x_{i_r},\qquad
 h_r(X)=\sum_{i_1\le\cdots\le i_r}x_{i_1}\cdots x_{i_r},
 \qquad p_r(X)=\sum_i x_i^r.
\]
Both degree-zero functions are one. Direct multiplication gives
\begin{align}
 \sum_{r\ge0}e_r(X)z^r&=\prod_i(1+x_i z),\label{q1:qcc:e-gf}\\
 \sum_{r\ge0}h_r(X)z^r&=\prod_i(1-x_i z)^{-1}.\label{q1:qcc:h-gf}
\end{align}
Taking formal logarithms gives, for \(r\ge0\),
\begin{align}
 e_r(X)&=\frac1{r!}\Y r\bigl(p_1,-1!p_2,2!p_3,\ldots,
                         (-1)^{r-1}(r-1)!p_r\bigr),\label{q1:qcc:e-bell}\\
 h_r(X)&=\frac1{r!}\Y r\bigl(p_1,1!p_2,2!p_3,\ldots,(r-1)!p_r\bigr).
 \label{q1:qcc:h-bell}
\end{align}
These identities follow from \eqref{q1:qcc:bell-def}; no additional combinatorial
normalization is involved. Also
\(\sum_{j=0}^r(-1)^j e_j(X)h_{r-j}(X)=\delta_{r0}\), because the two
generating functions with opposite signs multiply to one.

\begin{theorem}[Principal specialization]\label{q1:qcc:principal}
For \(X_N=(1,q,\ldots,q^{N-1})\),
\begin{equation}\label{q1:qcc:principal-formula}
 p_r(X_N)=\qi[q^r]N,\qquad
 e_k(X_N)=q^{\binom k2}\GB Nk,\qquad
 h_k(X_N)=\GB{N+k-1}k\quad(N\ge1).
\end{equation}
\end{theorem}
\begin{proof}
The first formula is a geometric sum. The other two follow by comparing
\eqref{q1:qcc:e-gf} and \eqref{q1:qcc:h-gf} with \eqref{q1:qcc:finite} and
\eqref{q1:qcc:reciprocal}.
\end{proof}
Thus a Gaussian coefficient is an ordinary symmetric polynomial at geometric
nodes. In particular,
\begin{align}
 q^{\binom k2}\GB Nk
 &=\frac1{k!}\Y k\left(\left((-1)^{r-1}(r-1)!\qi[q^r]N\right)_{r=1}^k\right),
 \label{q1:qcc:gaussian-e-bell}\\
 \GB{N+k-1}k
 &=\frac1{k!}\Y k\left(\left((r-1)!\qi[q^r]N\right)_{r=1}^k\right).
 \label{q1:qcc:gaussian-h-bell}
\end{align}
For instance,
\[
 \GB{N+1}2=\frac{\qi N^2+\qi[q^2]N}{2},\qquad
 q\GB N2=\frac{\qi N^2-\qi[q^2]N}{2}.
\]
The apparent halves cancel because these are symmetric-function identities over
\(\Z[q]\).

\section{Arbitrary powers of geometric products}
\begin{theorem}[Geometric-product coefficient formula]\label{q1:qcc:master-theorem}
Let \(N_\ell\ge0\), \(d_\ell\ge1\), and let \(a_\ell,\alpha_\ell\) be
commuting parameters in a \(\Q\)-algebra. Define formally
\begin{equation}\label{q1:qcc:master-product}
 F(z)=\prod_{\ell=1}^L(a_\ell z;q^{d_\ell})_{N_\ell}^{-\alpha_\ell},
 \qquad
 P_r=\sum_{\ell=1}^L\alpha_\ell a_\ell^r\qi[q^{d_\ell r}]{N_\ell}.
\end{equation}
Powers with nonintegral exponents mean \(\exp(-\alpha_\ell\log(\cdot))\)
with constant term one. Then
\begin{equation}\label{q1:qcc:master-coefficient}
 [z^n]F(z)=\frac1{n!}\Y n\bigl(P_1,1!P_2,\ldots,(n-1)!P_n\bigr)
 =\sum_{\substack{m_1,\ldots,m_n\ge0\\\sum r m_r=n}}
 \prod_{r=1}^n\frac{P_r^{m_r}}{r^{m_r}m_r!}.
\end{equation}
\end{theorem}
\begin{proof}
Each factor has a formal logarithm, and
\[
 \log F(z)
 =\sum_{\ell=1}^L\alpha_\ell\sum_{j=0}^{N_\ell-1}
       \sum_{r\ge1}\frac{(a_\ell q^{d_\ell j}z)^r}{r}
 =\sum_{r\ge1}P_r\frac{z^r}{r}.
\]
At every fixed degree this rearranges only finitely many contributions from the
finite alphabets. Apply \eqref{q1:qcc:bell-def} with
\(x_r=(r-1)!P_r\), then \eqref{q1:qcc:bell-finite}.
\end{proof}
The first four coefficients are
\begin{align*}
 f_1&=P_1,\qquad f_2=\frac{P_1^2+P_2}{2},\\
 f_3&=\frac{P_1^3+3P_1P_2+2P_3}{6},\\
 f_4&=\frac{P_1^4+6P_1^2P_2+3P_2^2+8P_1P_3+6P_4}{24}.
\end{align*}
Negative integral \(\alpha_\ell\) encode numerator products, positive integral
ones encode repeated reciprocal products, and arbitrary parameters encode
binomial powers. Several bases \(q^{d_\ell}\) require no new machinery.

\begin{corollary}[A two-base convolution]\label{q1:qcc:two-base}
For \(M,N\ge1\),
\begin{multline}\label{q1:qcc:two-base-formula}
 \sum_{j=0}^n a^j b^{n-j}
 \GB[q]{M+j-1}j\GB[q^d]{N+n-j-1}{n-j}\\
 =\frac1{n!}\Y n\left(
 \left((r-1)!\bigl(a^r\qi[q^r]M+b^r\qi[q^{dr}]N\bigr)\right)_{r=1}^n
 \right).
\end{multline}
\end{corollary}
\begin{proof}
Expand the two reciprocal products using \eqref{q1:qcc:reciprocal} and multiply.
Alternatively, apply \cref{q1:qcc:master-theorem} directly to their product. Equating
the two coefficients proves the identity.
\end{proof}
% ---- hand chunk 06-ramified-head.tex ----

\section{The ramified kernel}
\label{q2:sec:ramified}
\Cref{q1:qcc:master-theorem} allows each factor its own base \(q^{d_\ell}\).
A factor may also depend on a power \(z^{d}\) of the series variable. The
following statement, in the partition-profile notation of
\cref{q2:sec:kernel}, records the resulting ramification rule, which cannot
be omitted; it also covers infinite products through their logarithms.


% ---- Q2 lines 434--476 ----
For a finite $N\ge0$ define
\begin{equation}\label{q2:eq:S}
 S_N(t)=\sum_{j=0}^{N-1}t^j.
\end{equation}
This polynomial, rather than $(1-t^N)/(1-t)$, is the preferred expression at
$t=1$. When $N=\infty$, write $S_\infty(t)=(1-t)^{-1}$ with the formal or
analytic interpretation of \cref{q1:qcc:infinite-binomial}.

\begin{theorem}[Ramified shifted-factorial coefficient calculus]\label{q2:thm:master}
Let $d_r,h_r$ be positive integers, $N_r\in\N\cup\{\infty\}$, and let
$a_r,\alpha_r$ be commuting parameters, $1\le r\le s$. Set
\begin{equation}\label{q2:eq:master-P}
 P(z)=\prod_{r=1}^s\poch{a_r z^{d_r}}{q^{h_r}}{N_r}^{\alpha_r},
\end{equation}
and define, by finite sums for each $j\ge1$,
\begin{equation}\label{q2:eq:master-L}
 L_j=-\sum_{\substack{1\le r\le s\\d_r\mid j}}
       d_r\alpha_r a_r^{j/d_r}
             S_{N_r}\!\left(q^{h_rj/d_r}\right).
\end{equation}
Then for every $n\ge0$ and every scalar $\beta$,
\begin{equation}\label{q2:eq:master-coef}
 [z^n]P(z)^\beta
 =\Ecal_n(\beta L)
 =\sum_{\bm m\in\Pcal(n)}\prod_{j=1}^n
                   \frac{(\beta L_j/j)^{m_j}}{m_j!}.
\end{equation}
Here infinite factors are defined by their logarithms over $\Q(q)$; finite
factors need no such qualification.
\end{theorem}
\begin{proof}
For one finite factor,
\[
 \log\poch{az^d}{q^h}N
  =-\sum_{\ell\ge1}\frac{a^\ell}{\ell}
               S_N(q^{h\ell})z^{d\ell}.
\]
The same formula defines an infinite factor with $S_\infty$. In degree $j$, a
factor contributes only if $d\mid j$; writing $\ell=j/d$ makes its contribution
$-d a^{j/d}S_N(q^{hj/d})/j$. Summing the exponents $\alpha_r$ yields
$\log P=\sum_{j\ge1}L_jz^j/j$. Equation \eqref{q2:eq:E-Bell} proves
\eqref{q2:eq:master-coef}. All coefficient extractions use only $L_1,\ldots,L_n$.
\end{proof}
% ---- Q2 lines 501--535 ----
\section{Composition, arbitrary powers, and parameter derivatives}
\begin{corollary}[Composition with a unit product]\label{q2:cor:compose-P}
If $H(u)=\sum_{k\ge0}h_ku^k$ and $P$ is as in \cref{q2:thm:master}, then
\begin{equation}\label{q2:eq:compose-P}
 [z^n]H(P(z)-1)
 =\sum_{k=0}^n h_k\sum_{j=0}^k(-1)^{k-j}\binom kj\Ecal_n(jL).
\end{equation}
If $M(\bm m)=\sum_jm_j$, then
\begin{equation}\label{q2:eq:beta-derivative}
 [z^n]P(z)^\beta(\log P(z))^r
 =\sum_{\bm m\in\Pcal(n)}
   (M(\bm m))_{\underline r}\,\beta^{M(\bm m)-r}
   \prod_{j=1}^n\frac{(L_j/j)^{m_j}}{m_j!},
\end{equation}
where terms with $M(\bm m)<r$ are zero, so no negative power of $\beta$ is used.
\end{corollary}
\begin{proof}
Expand $(P-1)^k=\sum_{j=0}^k(-1)^{k-j}\binom kjP^j$.
Since $P-1$ has positive $z$-order, only $k\le n$ can contribute.
For the second formula differentiate \eqref{q2:eq:master-coef} $r$ times in $\beta$;
its degree in $\beta$ is at most $n$.
\end{proof}

For example, put
\[
 P(z)=\frac{\poch{-z^2}{q^2}M^\gamma}{\poch{az}qN^\alpha}.
\]
The logarithmic inputs are explicitly
\begin{equation}\label{q2:eq:mixed-example-L}
 L_j=\alpha a^jS_N(q^j)
   -\mathbf1_{2\mid j}\,2\gamma(-1)^{j/2}S_M(q^j).
\end{equation}
Thus all powers, logarithmic factors, and compositions of this mixed even/ordinary
product are evaluated by one finite kernel. This example also makes clear why the
ramification factor $d_r$ in \eqref{q2:eq:master-L} cannot be omitted.
% ---- hand chunk 07-pochpower.tex ----

\section{Reciprocal powers by ordinary binomials}
For a single reciprocal factor, the Bell form of
\cref{q1:qcc:master-theorem} has an elementary companion. Write
\(\lambda\) for the exponent.

% ---- Q3 lines 1068--1089 ----
For a formal parameter $\lambda$, \cref{q1:qcc:master-theorem} gives
\begin{equation}
 \coef zn(z;q)_N^{-\lambda}
 =\sum_{\sum j m_j=n}\prod_{j=1}^n
       \frac{\lambda^{m_j}[N]_{q^j}^{m_j}}{j^{m_j}m_j!}.
 \label{q3:eq:poch-power}
\end{equation}
It is polynomial in $\lambda$ of degree at most $n$. With $\lambda$ a
positive integer, one may instead expand each factor ordinarily:
\begin{equation}
 \coef zn(z;q)_N^{-\lambda}
 =\sum_{d_0+\cdots+d_{N-1}=n}
       q^{\sum_{j=0}^{N-1}jd_j}
       \prod_{j=0}^{N-1}\binom{\lambda+d_j-1}{d_j}.
 \label{q3:eq:poch-power-factor}
\end{equation}
The proof is multiplication of $(1-q^jz)^{-\lambda}$; hence
\cref{q3:eq:poch-power,q3:eq:poch-power-factor} form an explicit identity
between a partition sum and an ordinary-binomial weighted sum.
More general mixed ratios require only the signed power sums $L_j$ of
\cref{q1:qcc:master-theorem}. In particular, negative $\alpha_\ell$ introduce numerator
factors without changing the coefficient mechanism.
% ---- Q1 lines 619--653 ----
\section{A determinant and the recovery of logarithmic data}
\begin{proposition}\label{q1:qcc:determinant}
If \(F(z)=\exp(\sum_{r\ge1}P_rz^r/r)=\sum f_nz^n\), then
\begin{equation}\label{q1:qcc:det-formula}
 f_n=\frac1{n!}
 \det\begin{pmatrix}
 P_1&-1&0&\cdots&0\\
 P_2&P_1&-2&\cdots&0\\
 P_3&P_2&P_1&\ddots&\vdots\\
 \vdots&\vdots&\ddots&\ddots&-(n-1)\\
 P_n&P_{n-1}&\cdots&P_2&P_1
 \end{pmatrix}.
\end{equation}
Conversely, put \(x_j=j!f_j\). Then
\begin{equation}\label{q1:qcc:inverse-bell-log}
 P_n=\frac1{(n-1)!}\sum_{k=1}^n(-1)^{k-1}(k-1)!
                   \B nk(x_1,\ldots,x_{n-k+1}).
\end{equation}
\end{proposition}
\begin{proof}
The equation \(F'=(\sum_{r\ge1}P_rz^{r-1})F\) gives
\(nf_n=\sum_{r=1}^nP_rf_{n-r}\). If \(D_n\) is the determinant above, expansion
along its last row gives
\[
 D_n=\sum_{r=1}^n\frac{(n-1)!}{(n-r)!}P_rD_{n-r},\qquad D_0=1.
\]
The superdiagonal products supply the factorial ratio, and the Laplace signs
cancel their signs. Therefore \(D_n/n!\) satisfies the same recurrence and
initial value as \(f_n\).
For the inverse formula, expand
\(\log F=\sum_{k\ge1}(-1)^{k-1}(F-1)^k/k\), use
\eqref{q1:qcc:partial-def}, and compare \(z^n\) with \(P_n/n\).
\end{proof}
This is precisely the moment--cumulant mechanism in the canonical manuscript, with
geometric power sums serving as the logarithmic data.
% ---- Q1 lines 655--742 ----
\chapter{Partial fractions and divisor identities}
\label{q1:qcc:divisor}
\section{An interpolation formula for complete symmetric functions}
\begin{lemma}\label{q1:qcc:partial-fraction}
If \(x_1,\ldots,x_N\) are distinct and nonzero, then
\begin{equation}\label{q1:qcc:partial-fraction-formula}
 \prod_{i=1}^N\frac1{1-x_i z}
 =\sum_{i=1}^N\frac{x_i^{N-1}}{\prod_{j\ne i}(x_i-x_j)}\frac1{1-x_i z}.
\end{equation}
Consequently,
\begin{equation}\label{q1:qcc:h-divided}
 h_r(x_1,\ldots,x_N)
 =\sum_{i=1}^N\frac{x_i^{r+N-1}}{\prod_{j\ne i}(x_i-x_j)}.
\end{equation}
\end{lemma}
\begin{proof}
Multiply the proposed rational identity by \(\prod_j(1-x_jz)\).
Both sides are polynomials of degree at most \(N-1\); at \(z=1/x_i\), only
one summand on the right survives, and it equals one. Agreement at \(N\)
distinct points proves the identity. Expanding each geometric denominator and
comparing coefficients proves \eqref{q1:qcc:h-divided}.
\end{proof}
The divided-difference expression on the right has removable singularities when
nodes coincide; the left side supplies the polynomial continuation. This is a
simple instance of why a rational summation formula can encode a polynomial.

\section{A finite divisor-sum master identity}
\begin{theorem}[Dilcher-type identity and Bell form]\label{q1:qcc:dilcher}
For \(N\ge1\), \(r\ge0\), and indeterminate \(q\), let
\(u_j=q^j/(1-q^j)\). Then
\begin{align}
 &\sum_{j=1}^N(-1)^{j-1}
  \frac{q^{\binom j2+rj}}{(1-q^j)^r}\GB Nj
 =h_r(u_1,\ldots,u_N)\label{q1:qcc:dilcher-formula}\\
 &\hspace{13mm}
 =\sum_{1\le i_1\le\cdots\le i_r\le N}
      \frac{q^{i_1+\cdots+i_r}}{(1-q^{i_1})\cdots(1-q^{i_r})}
 \nonumber\\
 &\hspace{13mm}
 =\frac1{r!}\Y r\left(
       \left((s-1)!\sum_{j=1}^N\frac{q^{js}}{(1-q^j)^s}\right)_{s=1}^r
       \right).
 \label{q1:qcc:dilcher-bell}
\end{align}
For \(r=0\), every right-hand side is interpreted as one.
\end{theorem}
\begin{proof}
Apply \cref{q1:qcc:partial-fraction} to \(q,q^2,\ldots,q^N\). For the node
\(q^j\), the residue coefficient is
\[
 \frac1{\prod_{i\ne j}(1-q^{i-j})}
 =\frac{(-1)^{j-1}q^{\binom j2}}{(q;q)_{j-1}(q;q)_{N-j}}.
\]
Multiplication by \((q;q)_N\) therefore yields
\begin{equation}\label{q1:qcc:divisor-rational}
 \frac{(q;q)_N}{(zq;q)_N}
 =\sum_{j=1}^N(-1)^{j-1}q^{\binom j2}(1-q^j)\GB Nj\frac1{1-zq^j}.
\end{equation}
Set \(z=1+t\). The left side becomes
\(\prod_{j=1}^N(1-u_jt)^{-1}\).
The coefficient of \(t^r\) on the right is exactly the left side of
\eqref{q1:qcc:dilcher-formula}. The complete-symmetric and Bell descriptions now
follow from \eqref{q1:qcc:h-gf} and \eqref{q1:qcc:h-bell}.
\end{proof}
This is the finite identity associated with Dilcher's divisor-function
identities \cite{dilcher}. The partial-fraction proof simultaneously explains the
alternating single sum, the weakly ordered multiple sum, and the Bell polynomial.
For \(r=1\), it reduces to
\[
 \sum_{j=1}^N(-1)^{j-1}\frac{q^{j(j+1)/2}}{1-q^j}\GB Nj
 =\sum_{j=1}^N\frac{q^j}{1-q^j}.
\]
For \(r=2\), the right side is \(\frac12((\sum u_j)^2+\sum u_j^2)\),
and for \(r=3\), it is
\(\frac16((\sum u_j)^3+3(\sum u_j)(\sum u_j^2)+2\sum u_j^3)\).

\section{The ordinary harmonic limit}
Multiply \eqref{q1:qcc:dilcher-formula} by \((1-q)^r\) and let \(q\to1\).
Since \((1-q)u_j\to1/j\), the limit is
\begin{equation}\label{q1:qcc:harmonic-limit}
 \sum_{j=1}^N(-1)^{j-1}\binom Nj\frac1{j^r}
 =h_r\left(1,\frac12,\ldots,\frac1N\right)
 =\frac1{r!}\Y r\bigl(H_N^{(1)},1!H_N^{(2)},\ldots,(r-1)!H_N^{(r)}\bigr),
\end{equation}
where \(H_N^{(s)}=\sum_{j=1}^N j^{-s}\). This limit is termwise because all sums
are finite. Thus the ordinary binomial--harmonic identities of coefficient
calculus are literal degenerations of the Gaussian identities, not merely
analogous-looking formulae.
% ---- hand chunk 08-newton-head.tex ----

\chapter{Newton bases on arbitrary nodes}
\label{q3:sec:newton}
Gaussian coefficients are Newton connection coefficients for geometric
nodes, in the same sense in which ordinary Stirling numbers are the
connection coefficients for the nodes \(0,1,2,\ldots\). This chapter develops
the node calculus in general, then specializes it to geometric,
affine-geometric, and \(q\)-integer nodes. The affine-geometric case is the
engine behind the closed \(q\)-Stirling formulae of
\cref{q1:qcc:qstirling}.


% ---- Q3 lines 546--702 ----
\section{The general node calculus}
Let $a_0,a_1,\ldots$ be elements of a commutative ring and set
\begin{equation}
 P_0^{\boldsymbol a}(x)=1,\qquad
 P_k^{\boldsymbol a}(x)=\prod_{j=0}^{k-1}(x-a_j).
 \label{q3:eq:node-basis}
\end{equation}
Since these polynomials are monic of successive degrees, there are unique
arrays $S_{\boldsymbol a}$ and $s_{\boldsymbol a}$ such that
\begin{equation}
 x^n=\sum_{k=0}^nS_{\boldsymbol a}(n,k)P_k^{\boldsymbol a}(x),\qquad
 P_n^{\boldsymbol a}(x)=\sum_{k=0}^ns_{\boldsymbol a}(n,k)x^k.
 \label{q3:eq:node-connections}
\end{equation}
This common framework for binomial, Stirling, and Gaussian coefficients
is classical; Voigt's short paper \cite{voigt} is particularly close in
spirit to the basis-change program of the canonical manuscript.

For a finite list $A=(a_0,\ldots,a_k)$, the elementary and complete
homogeneous polynomials are defined by
\[
 e_m(A)=\sum_{i_1<\cdots<i_m}a_{i_1}\cdots a_{i_m},\qquad
 h_m(A)=\sum_{i_1\le\cdots\le i_m}a_{i_1}\cdots a_{i_m}.
\]
Both have value one at $m=0$; $e_m=0$ above the list length.

\begin{theorem}[Universal finite node formulae]
\label{q3:thm:node}
For $0\le k\le n$,
\begin{align}
 S_{\boldsymbol a}(n,k)&=h_{n-k}(a_0,\ldots,a_k),
 \label{q3:eq:node-h}\\
 s_{\boldsymbol a}(n,k)&=(-1)^{n-k}e_{n-k}(a_0,\ldots,a_{n-1}).
 \label{q3:eq:node-e}
\end{align}
The arrays are inverse lower-triangular matrices. In particular,
\begin{equation}
 \sum_{j=k}^ns_{\boldsymbol a}(n,j)S_{\boldsymbol a}(j,k)=\delta_{n,k}.
 \label{q3:eq:node-inverse}
\end{equation}
Their recurrences are
\begin{align}
 S_{\boldsymbol a}(n+1,k)&=S_{\boldsymbol a}(n,k-1)+a_kS_{\boldsymbol a}(n,k),
 \label{q3:eq:node-Srec}\\
 s_{\boldsymbol a}(n+1,k)&=s_{\boldsymbol a}(n,k-1)-a_ns_{\boldsymbol a}(n,k).
 \label{q3:eq:node-srec}
\end{align}
\end{theorem}
\begin{proof}
The identity $xP_k=P_{k+1}+a_kP_k$ gives
\cref{q3:eq:node-Srec}. Separating monomials in $h_m(a_0,\ldots,a_k)$
according as they contain $a_k$ gives the same recurrence for
\cref{q3:eq:node-h}, including the boundary values. Formula
\cref{q3:eq:node-e} follows by expanding the product, and multiplication by
$x-a_n$ gives \cref{q3:eq:node-srec}. Substituting one basis expansion into
the other proves matrix inversion by uniqueness of monomial coefficients.
\end{proof}
All these identities are valid even when nodes coincide. Only the next
representation requires distinct nodes.

\begin{theorem}[Divided differences and a nonrecursive spectral sum]
\label{q3:thm:divdiff}
Over a field, assume $a_0,\ldots,a_k$ are distinct. Then
\begin{equation}
 S_{\boldsymbol a}(n,k)
 =\sum_{j=0}^k\frac{a_j^n}{\prod_{0\le i\le k,\ i\ne j}(a_j-a_i)}.
 \label{q3:eq:node-spectral}
\end{equation}
More generally, the coefficient of $P_k^{\boldsymbol a}$ in the Newton
interpolant to values $f(a_0),\ldots,f(a_M)$ is
\begin{equation}
 f[a_0,\ldots,a_k]
 =\sum_{j=0}^k\frac{f(a_j)}{\prod_{i\ne j}(a_j-a_i)}.
 \label{q3:eq:divided-difference}
\end{equation}
Thus its degree-$M$ interpolating polynomial is
$\sum_{k=0}^Mf[a_0,\ldots,a_k]P_k^{\boldsymbol a}(x)$.
\end{theorem}
\begin{proof}
The Lagrange interpolant on the first $k+1$ nodes is
\[
 L_k(x)=\sum_{j=0}^kf(a_j)
 \frac{\prod_{i\ne j}(x-a_i)}{\prod_{i\ne j}(a_j-a_i)}.
\]
Its leading coefficient is the sum in \cref{q3:eq:divided-difference}.
In a Newton expansion, every term with index greater than $k$ vanishes
on these nodes, and every term with index less than $k$ has smaller
degree. Hence that leading coefficient is precisely the $k$th Newton
coefficient. Apply this to $f(x)=x^n$ for \cref{q3:eq:node-spectral}.
\end{proof}

\section{Geometric nodes and the Gaussian basis pair}
\begin{theorem}[The Gaussian Newton pair]
\label{q3:thm:gaussian-newton}
Put $P_k(x)=\prod_{j=0}^{k-1}(x-q^j)$. Then
\begin{align}
 x^n&=\sum_{k=0}^n\G nk P_k(x),
 \label{q3:eq:gaussian-newton}\\
 P_n(x)&=\sum_{k=0}^n(-1)^{n-k}q^{\binom{n-k}2}\G nk x^k.
 \label{q3:eq:gaussian-newton-inverse}
\end{align}
\end{theorem}
\begin{proof}
By \cref{q3:eq:node-h}, the first coefficient is
$h_{n-k}(1,q,\ldots,q^k)$, which is $G_{n,k}(q)$ by
\cref{q1:qcc:reciprocal}. The second is the elementary symmetric
coefficient in the same geometric list, and \cref{q1:qcc:finite} gives
exactly the stated phase. Equivalently, the first recurrence is
\cref{q1:qcc:pascal1} with $a_k=q^k$.
\end{proof}
This is the precise sense in which Gaussian coefficients are
Stirling-type coefficients: the nodes $0,1,2,\ldots$ produce ordinary
Stirling numbers, while $1,q,q^2,\ldots$ produce Gaussian polynomials.
At $q=1$ these geometric nodes all equal one, so
\cref{q3:eq:gaussian-newton} becomes the ordinary binomial expansion
$x^n=\sum_k\binom nk(x-1)^k$, not the ordinary Stirling expansion.

For generic nonzero $q$, the explicit geometric-grid formula is
\begin{equation}
 f[1,q,\ldots,q^k]
 =\frac{q^{-\binom k2}}{(q;q)_k}
   \sum_{j=0}^k(-1)^jq^{\binom{k-j}2}\G kj f(q^j).
 \label{q3:eq:geometric-dd}
\end{equation}
Indeed,
\begin{equation}
 \prod_{i\ne j}(q^j-q^i)
 =(-1)^jq^{jk-j(j+1)/2}(q;q)_j(q;q)_{k-j},
 \label{q3:eq:geometric-denominator}
\end{equation}
obtained by separating $i<j$ from $i>j$. Substitution in
\cref{q3:eq:divided-difference} proves \cref{q3:eq:geometric-dd}.
This is a finite interpolation formula; an infinite Newton expansion
would need separate convergence assumptions.

\section{An affine-geometric master connection}
\begin{theorem}[Affine-geometric nodes]
\label{q3:thm:affine}
For nodes $a_j=A+Bq^j$, the connection coefficients are
\begin{equation}
 S_{A+Bq^\bullet}(n,k)
 =\sum_{d=k}^n\binom nd A^{n-d}B^{d-k}\G dk.
 \label{q3:eq:affine-geometric}
\end{equation}
This is a polynomial identity in $A,B,q$, including $B=0$.
\end{theorem}
\begin{proof}
First take $B$ invertible and write $x=A+By$. Then
$P_k^{\boldsymbol a}(x)=B^kP_k(y)$. Expand $(A+By)^n$ ordinarily,
then expand each $y^d$ by \cref{q3:eq:gaussian-newton}. The coefficient of
$B^kP_k(y)$ is \cref{q3:eq:affine-geometric}. Both sides are polynomials,
so the identity descends from $\Q(A,B,q)$ to $\Z[A,B,q]$ and can be
specialized at $B=0$.
\end{proof}
This theorem is a compact engine for several families that otherwise
look unrelated. The next section makes its $q$-arithmetic specializations
explicit.
% ---- Q2 lines 690--721 ----
\section{The \texorpdfstring{$q$}{q}-integer lattice: the deformation of ordinary Newton expansion}
Define the monic polynomials
\begin{equation}\label{q2:eq:Fk}
 F_k(x)=\prod_{i=0}^{k-1}(x-\qi i),\qquad F_0(x)=1.
\end{equation}
Their normalized versions
\[
 N_k(x;q)=\frac{F_k(x)}{q^{\binom k2}\qf k}
\]
satisfy $N_k(\qi m;q)=\qb mk$ for $m\ge0$. To see this, use
$\qi m-\qi i=q^i\qi{m-i}$ and multiply over $i$.
This basis belongs to the broader setting of quantum integer-valued polynomials
\cite{harmanhopkins}; its normalization is stated here explicitly.

\begin{theorem}[Finite $q$-Newton expansion]\label{q2:thm:qnewton}
For every polynomial $p$ of degree at most $d$ over $R$,
\begin{align}
 p(x)&=\sum_{k=0}^d c_k N_k(x;q),\label{q2:eq:qnewton}\\
 c_k&=\sum_{j=0}^k(-1)^{k-j}q^{\binom{k-j}2}\qb kj p(\qi j).
 \label{q2:eq:qnewton-coef}
\end{align}
At $q=1$ this becomes
$p(x)=\sum_{k=0}^d\Delta^kp(0)\binom xk$.
\end{theorem}
\begin{proof}
The distinct degrees and unit leading coefficients of $N_0,\ldots,N_d$ make
them a basis. Evaluate its expansion at $x=\qi m$ for $0\le m\le d$:
$p(\qi m)=\sum_{k=0}^m\qb mk c_k$. Gaussian inversion gives
\eqref{q2:eq:qnewton-coef}. At $q=1$, the factors in $N_k$ become
$(x-i)$, their denominator becomes $k!$, and the sum for $c_k$ is the ordinary
finite-difference formula. All denominators in this specialization are nonzero.
\end{proof}
% ---- Q1 lines 744--878 ----
\chapter{Jackson differentiation and geometric Newton interpolation}
\label{q1:qcc:jackson}
\section{The operator and its commutation rule}
Let \(T_qf(x)=f(qx)\), and let \(X\) denote multiplication by \(x\). Define
\begin{equation}\label{q1:qcc:jackson-def}
 D_qf(x)=\frac{f(x)-f(qx)}{(1-q)x}.
\end{equation}
For polynomials, this has no singularity at \(x=0\), because the numerator has
zero constant term. It acts by
\(D_qx^m=[m]_qx^{m-1}\), and hence
\begin{equation}\label{q1:qcc:q-weyl}
 D_qX=1+qXD_q,\qquad D_qT_q=qT_qD_q.
\end{equation}
These identities follow on every monomial and hence on every polynomial. This is
the convention for the Jackson derivative used throughout; compare
\cite{dlmf172}.

\begin{theorem}[The shifted product rule]\label{q1:qcc:leibniz}
For \(n\ge0\),
\begin{equation}\label{q1:qcc:q-leibniz}
 D_q^n(fg)(x)=\sum_{k=0}^n\GB nk
 (D_q^k f)(q^{n-k}x)(D_q^{n-k}g)(x).
\end{equation}
\end{theorem}
\begin{proof}
The first-order formula is
\(D_q(fg)=(D_qf)g+(T_qf)D_qg\), obtained by adding and subtracting
\(f(qx)g(x)\) in the numerator. Suppose the formula is true at order \(n\).
Differentiating a summand gives a contribution with \(k\) unchanged and a
contribution with \(k\) increased by one. In the latter, differentiation of
\((D_q^kf)(q^{n-k}x)\) produces a factor \(q^{n-k}\).
Thus the new coefficient at index \(k\) is
\(\GB nk+q^{n-k+1}\GB n{k-1}=\GB{n+1}k\), by
\eqref{q1:qcc:pascal1}. The shifts are exactly those in the stated formula.
\end{proof}
\begin{auditbox}
The argument shift \(q^{n-k}x\) is part of the theorem. Replacing the ordinary
binomial coefficient by a Gaussian coefficient in the usual Leibniz rule, while
leaving all factors evaluated at \(x\), is false.
\end{auditbox}

\section{A geometric Taylor basis}
Define
\begin{equation}\label{q1:qcc:q-newton-basis}
 (x-a)_{q}^{[k]}=\prod_{j=0}^{k-1}(x-aq^j).
\end{equation}
This is a polynomial basis in \(x\), not a power of \(x-a\).
\begin{theorem}[Jackson--Newton expansion]\label{q1:qcc:q-taylor}
For every polynomial \(p\) of degree at most \(d\),
\begin{equation}\label{q1:qcc:q-taylor-formula}
 p(x)=\sum_{k=0}^d\frac{(D_q^kp)(a)}{\qf k}(x-a)_q^{[k]}.
\end{equation}
The identity first holds over \(\Q(q,a)\), and specializes wherever its
coefficients are defined. At \(q=1\) it becomes ordinary Taylor expansion at
\(a\).
\end{theorem}
\begin{proof}
The polynomial \(D_q(x-a)_q^{[k]}\) has degree \(k-1\) and leading coefficient
\([k]_q\). At each \(x=aq^j\), \(0\le j\le k-2\), its two numerator terms
vanish. Therefore
\[
 D_q(x-a)_q^{[k]}=[k]_q(x-a)_q^{[k-1]}.
\]
This argument at generic \(a,q\) proves a polynomial identity and hence also
covers \(a=0\). Expand \(p\) in this monic basis, apply \(D_q^k\), and set
\(x=a\). Terms with index less than \(k\) vanish under differentiation; terms
with larger index retain a factor \(x-a\); the index \(k\) term yields
\([k]_q!\) times its coefficient. This proves the formula.
\end{proof}
For example,
\[
 x^2=a^2+(1+q)a(x-a)+(x-a)(x-aq),
\]
and
\[
 x^3=a^3+[3]_qa^2(x-a)+[3]_qa(x-a)(x-aq)
                  +(x-a)(x-aq)(x-aq^2).
\]
At a sample node \(x=aq^n\), factors with \(k>n\) vanish, and direct
factorization gives the discrete Gaussian transform
\begin{equation}\label{q1:qcc:sample-transform}
 p(aq^n)=\sum_{k=0}^n\GB nk\bigl(a(q-1)\bigr)^k
                  q^{\binom k2}(D_q^kp)(a).
\end{equation}
Indeed,
\(\prod_{j<k}(aq^n-aq^j)
 =[a(q-1)]^kq^{\binom k2}[k]_q!\GB nk\).
This identifies Gaussian inversion with the extraction of geometric Newton
coefficients.

\section{Explicit stencils and divided differences}
\begin{theorem}\label{q1:qcc:stencil}
For \(n\ge0\),
\begin{equation}\label{q1:qcc:q-stencil}
 D_q^nf(x)=\frac1{(1-q)^nx^n}
 \sum_{j=0}^n(-1)^j q^{\binom j2-(n-1)j}\GB nj f(q^jx).
\end{equation}
At distinct nonzero nodes \(a,aq,\ldots,aq^n\),
\begin{equation}\label{q1:qcc:divided-difference}
 f[a,aq,\ldots,aq^n]=\frac{(D_q^nf)(a)}{\qf n}.
\end{equation}
The left side denotes the leading coefficient of the degree-at-most-\(n\)
interpolating polynomial through the given data.
\end{theorem}
\begin{proof}
The commutation relation \(T_qX^{-1}=q^{-1}X^{-1}T_q\) gives
\[
 D_q^n=(1-q)^{-n}X^{-n}\prod_{r=0}^{n-1}(1-q^{-r}T_q).
\]
Expand the finite product with base \(q^{-1}\), and use reciprocity
\eqref{q1:qcc:reciprocity}. Its \(T_q^j\) coefficient becomes
\((-1)^jq^{\binom j2-(n-1)j}\GB nj\), proving
\eqref{q1:qcc:q-stencil}.
Let \(p\) interpolate \(f\) at these nodes. Formula \eqref{q1:qcc:q-stencil}
shows that \((D_q^nf)(a)=(D_q^np)(a)\). The latter equals \([n]_q!\) times
the leading coefficient of \(p\), either by monomial differentiation or by
\cref{q1:qcc:q-taylor}. This proves \eqref{q1:qcc:divided-difference}.
\end{proof}

\begin{corollary}[An exact Gaussian coefficient extractor]\label{q1:qcc:extractor}
If \(p\) has degree at most \(n\), then
\begin{equation}\label{q1:qcc:extractor-formula}
 [x^n]p(x)=\frac1{(q;q)_n}
 \sum_{j=0}^n(-1)^j q^{\binom j2-(n-1)j}\GB nj p(q^j).
\end{equation}
The numerator annihilates \(x^r\) for every \(0\le r<n\).
\end{corollary}
\begin{proof}
Apply \eqref{q1:qcc:q-stencil} at \(x=1\) and divide by \([n]_q!\).
Alternatively, the numerator on \(x^r\) is
\((q^{r-n+1};q)_n\), by \eqref{q1:qcc:finite}. One factor is zero for
\(0\le r<n\), while for \(r=n\) the product is \((q;q)_n\).
\end{proof}
This is the geometric counterpart of the additive finite-difference stencil in
the source. Its cancelling objects are monomials evaluated at geometric nodes.
% ---- hand chunk 10a-multileibniz-head.tex ----

\section{The iterated product rule}
Repeated application of \cref{q1:qcc:leibniz} to the first factor and the
remaining product gives, for \(r\) factors,

% ---- Q3 lines 1181--1188 ----
\begin{equation}
 \Dq^n\!\left(\prod_{i=1}^rf_i\right)(x)
 =\sum_{n_1+\cdots+n_r=n}
    \frac{\qfac n}{\prod_i\qfac{n_i}}
    \prod_{i=1}^r(\Dq^{n_i}f_i)
              \left(q^{n_{i+1}+\cdots+n_r}x\right).
 \label{q3:eq:q-multi-leibniz}
\end{equation}
% ---- hand chunk 10b-multileibniz-tail.tex ----

The successive Gaussian coefficients telescope to the displayed Gaussian
multinomial, and the argument shifts accumulate from the right: the
\(i\)th factor is evaluated at \(q^{n_{i+1}+\cdots+n_r}x\). The same
coefficient occurs in the noncommutative binomial formula of
\cref{q2:prop:quantum-plane}.


% ---- Q3 lines 1216--1300 ----
\section{An exact Jackson chain rule through divided differences}
\label{q3:sec:chain}
The preceding distinctions become especially important for composition.
There is no general identity
$\Dq(f\circ g)=(\Dq f)\circ g\cdot\Dq g$.
For $f(u)=u^2$ and $g(x)=x^2$, its two sides would be
$\qnum4x^3$ and $\qnum2^2x^3$, which differ for generic $q$.
A correct finite chain rule uses divided differences. The general
construction belongs to the divided-difference chain-rule literature
\cite{floaterlyche}; the following proof also explains its relation to
coefficient calculus.

\begin{theorem}[Path form of the divided-difference chain rule]
Let $x_0,\ldots,x_n$ be distinct and put $y_i=g(x_i)$. Suppose first that
the $y_i$ are distinct. For $n\ge1$,
\begin{equation}
 \boxed{\displaystyle
 (f\circ g)[x_0,\ldots,x_n]
 =\sum_{k=1}^n\ \sum_{0=i_0<i_1<\cdots<i_k=n}
   f[y_{i_0},\ldots,y_{i_k}]
   \prod_{r=1}^k g[x_{i_{r-1}},\ldots,x_{i_r}].}
 \label{q3:eq:dd-chain}
\end{equation}
Whenever confluent limits exist, the same formula extends to coinciding
$y_i$ by replacing divided differences by their confluent values.
For $x_i=aq^i$ and $a\ne0$, multiplication of the right side by
$\qfac n$ gives $(\Dq^n(f\circ g))(a)$.
\end{theorem}
\begin{proof}
We give the finite matrix argument in detail. Let $T$ be the
$(n+1)\times(n+1)$ upper bidiagonal matrix whose diagonal is
$x_0,\ldots,x_n$ and whose first superdiagonal consists of ones. Direct
back substitution in $(zI-T)R=I$ shows that
\[
 R_{ij}=\frac1{\prod_{h=i}^j(z-x_h)}\quad(i\le j).
\]
For a polynomial $h$, applying polynomial functional calculus, or
expanding this rational expression by partial fractions, therefore gives
\begin{equation}
 h(T)_{ij}=h[x_i,\ldots,x_j].
 \label{q3:eq:matrix-dd}
\end{equation}
For completeness, the partial-fraction coefficient of $1/(z-x_r)$ is
$\prod_{i\le h\le j,\,h\ne r}(x_r-x_h)^{-1}$; multiplying by $h(x_r)$
and summing is exactly the explicit divided-difference formula.

More generally, let $A$ be upper triangular with diagonal $a_i$ and
strict upper part $N$. The finite expansion
\[
 (zI-A)^{-1}
 =\sum_{k=0}^n\bigl((zI-\operatorname{diag}A)^{-1}N\bigr)^k
                (zI-\operatorname{diag}A)^{-1}
\]
shows that its $(0,n)$ entry is a sum over increasing paths
$0=i_0<\cdots<i_k=n$. The contribution of a path is
\[
 \frac{A_{i_0i_1}\cdots A_{i_{k-1}i_k}}
      {(z-a_{i_0})\cdots(z-a_{i_k})}.
\]
The same partial-fraction calculation gives
\[
 f(A)_{0n}=\sum_{\text{paths}}% ed.: restored \text; the source line carried a tab character here
       f[a_{i_0},\ldots,a_{i_k}]
       A_{i_0i_1}\cdots A_{i_{k-1}i_k}.
\]
Take $A=g(T)$, use \cref{q3:eq:matrix-dd}, and use
$f(g(T))=(f\circ g)(T)$. This proves \cref{q3:eq:dd-chain} for
polynomials. For arbitrary data with distinct $y_i$, interpolate $g$
on the $x_i$ and $f$ on the $y_i$; every term is unchanged. Passage to
confluent limits proves the stated extension when the necessary
smoothness or analyticity is available. Finally apply
\cref{q1:qcc:divided-difference} to the geometric nodes.
\end{proof}

For $n=2$ the two paths give the particularly transparent identity
\begin{align}
 (f\circ g)[x_0,x_1,x_2]
 &=f[y_0,y_2]g[x_0,x_1,x_2] \notag\\
 &\quad+f[y_0,y_1,y_2]g[x_0,x_1]g[x_1,x_2].
 \label{q3:eq:dd-chain-two}
\end{align}
As $q\to1$, geometric nodes coalesce and divided differences become
ordinary derivatives divided by factorials. Thus this is a finite-node
precursor of ordinary Fa\`a di Bruno, rather than an unjustified
replacement of all factorials by $q$-factorials.
% ---- Q1 lines 880--1024 ----
\chapter{Two \qbname-Stirling arrays and normal ordering}
\label{q1:qcc:qstirling}
\section{The monic \qbname-integer factorial basis}
Define
\begin{equation}\label{q1:qcc:qstirling-basis}
 P_k(x)=\prod_{j=0}^{k-1}(x-[j]_q),\qquad
 P_n(x)=\sum_{k=0}^n s_q(n,k)x^k,\qquad
 x^n=\sum_{k=0}^n S_q(n,k)P_k(x).
\end{equation}
These are changes between two monic polynomial bases. They define \(s_q,S_q\)
unambiguously over \(\Z[q]\), even when specialized nodes coincide.

\begin{theorem}\label{q1:qcc:qstirling-basic}
The arrays are mutually inverse, and
\begin{align}
 s_q(n+1,k)&=s_q(n,k-1)-[n]_qs_q(n,k),\label{q1:qcc:sq-rec}\\
 S_q(n+1,k)&=S_q(n,k-1)+[k]_qS_q(n,k),\label{q1:qcc:Sq-rec}\\
 s_q(n,k)&=(-1)^{n-k}e_{n-k}([0]_q,\ldots,[n-1]_q),\label{q1:qcc:sq-e}\\
 S_q(n,k)&=h_{n-k}([0]_q,\ldots,[k]_q).\label{q1:qcc:Sq-h}
\end{align}
Both arrays have entry one at \((0,0)\) and vanish outside the triangular range.
\end{theorem}
\begin{proof}
Multiplication of \(P_n(x)\) by \(x-[n]_q\) gives
\eqref{q1:qcc:sq-rec}. The identity \(xP_k=P_{k+1}+[k]_qP_k\) gives
\eqref{q1:qcc:Sq-rec}. The coefficient formula \eqref{q1:qcc:sq-e} follows by
choosing roots from the product defining \(P_n\).
For \eqref{q1:qcc:Sq-h}, use
\[
 h_r([0]_q,\ldots,[k]_q)
 =h_r([0]_q,\ldots,[k-1]_q)+[k]_qh_{r-1}([0]_q,\ldots,[k]_q).
\]
With \(r=n-k\), this is precisely the recurrence for \(S_q\), with the same
boundary values. Finally, substituting either basis change into the other and
using uniqueness of coefficients proves matrix inversion.
\end{proof}
At \(q=1\), these are the ordinary signed first-kind and second-kind Stirling
numbers. At \(q=0\), all positive-index nodes equal one, so for \(n\ge1\),
\[
 S_0(n,k)=\binom{n-1}{k-1},\qquad
 s_0(n,k)=(-1)^{n-k}\binom{n-1}{k-1}\quad(1\le k\le n).
\]
This follows from \(P_n(x)=x(x-1)^{n-1}\) and from the complete-symmetric
formula. It illustrates the usefulness of a polynomial definition at degenerate
bases.

\section{A Gaussian closed sum, with every power of \qbname{} accounted for}
\begin{theorem}\label{q1:qcc:qstirling-explicit}
Over \(\Q(q)\), for \(0\le k\le n\),
\begin{equation}\label{q1:qcc:qstirling-sum}
 S_q(n,k)=\frac{q^{-\binom k2}}{[k]_q!}
 \sum_{j=0}^k(-1)^{k-j}q^{\binom{k-j}{2}}\GB kj[j]_q^n.
\end{equation}
The expression has the polynomial continuation specified by
\eqref{q1:qcc:qstirling-basis}.
\end{theorem}
\begin{proof}
The coefficient of the degree-\(k\) Newton basis polynomial in the interpolation
of \(x^n\) at \([0]_q,\ldots,[k]_q\) is
\[
 \sum_{j=0}^k\frac{[j]_q^n}{\prod_{0\le i\le k,\ i\ne j}([j]_q-[i]_q)}.
\]
It equals the coefficient in the full Newton expansion because terms of degree
larger than \(k\) vanish at all these nodes. For \(i<j\),
\([j]_q-[i]_q=q^i[j-i]_q\), whereas for \(i>j\) it is
\(-q^j[i-j]_q\). Hence the denominator is
\[
 (-1)^{k-j}q^{jk-j(j+1)/2}[j]_q![k-j]_q!.
\]
Its inverse, multiplied by \([k]_q!\), has exponent
\(\binom{j+1}{2}-jk=\binom{k-j}{2}-\binom k2\).
This proves \eqref{q1:qcc:qstirling-sum}.
\end{proof}
There is no guesswork in the factor \(q^{-\binom k2}\): it is the normalization
of the monic \(q\)-integer Newton basis. Omitting it produces a different array.

\section{Bell formulae for every diagonal}
For \(k\ge0\) and \(r\ge1\), let
\begin{equation}\label{q1:qcc:qinteger-powersum}
 U_r(k)=\sum_{j=0}^k[j]_q^r
 =\frac1{(1-q)^r}\sum_{a=0}^r(-1)^a\binom ra[k+1]_{q^a}.
\end{equation}
The \(a=0\) term means \(k+1\). The equality follows by expanding
\((1-q^j)^r\) binomially and summing the geometric powers \(q^{aj}\).
Equations \eqref{q1:qcc:sq-e}--\eqref{q1:qcc:Sq-h} and the Bell formulae give
\begin{align}
 S_q(n,k)&=\frac1{(n-k)!}\Y{n-k}
   \left(\left((r-1)!U_r(k)\right)_{r=1}^{n-k}\right),\label{q1:qcc:Sq-bell}\\
 s_q(n,k)&=\frac{(-1)^{n-k}}{(n-k)!}\Y{n-k}
   \left(\left((-1)^{r-1}(r-1)!U_r(n-1)\right)_{r=1}^{n-k}\right).
 \label{q1:qcc:sq-bell}
\end{align}
These are recurrence-free finite formulae involving only binomial coefficients,
\(q\)-integers, factorials, and the finite sum \eqref{q1:qcc:bell-finite}.
For example,
\[
 S_q(n,n-1)=\sum_{j=1}^{n-1}[j]_q
          =\sum_{a=0}^{n-2}(n-1-a)q^a,
\]
and
\[
 S_q(n,n-2)=\frac12\left(U_1(n-2)^2+U_2(n-2)\right).
\]
The latter uses the alphabet ending at \([n-2]_q\), not at \([n-1]_q\).

\section{The normal-ordering normalization}
Define
\begin{equation}\label{q1:qcc:normalized-stirling}
 \widetilde S_q(n,k)=q^{\binom k2}S_q(n,k).
\end{equation}
Then
\begin{equation}\label{q1:qcc:normalized-stirling-rec}
 \widetilde S_q(n+1,k)=[k]_q\widetilde S_q(n,k)
                         +q^{k-1}\widetilde S_q(n,k-1).
\end{equation}
Both this and \eqref{q1:qcc:Sq-rec} occur in \(q\)-Stirling literature; the
restricted-growth-word viewpoint and related identities are discussed in
\cite{cai}. The defining recurrence must accompany any comparison of conventions.

\begin{theorem}[Normal ordering and factorial moments]\label{q1:qcc:normal-ordering}
On polynomials,
\begin{equation}\label{q1:qcc:normal-ordering-formula}
 (XD_q)^n=\sum_{k=0}^n\widetilde S_q(n,k)X^kD_q^k.
\end{equation}
For integers \(m,n\ge0\),
\begin{equation}\label{q1:qcc:factorial-moments}
 [m]_q^n=\sum_{k=0}^{\min(m,n)}
          \widetilde S_q(n,k)[k]_q!\GB mk.
\end{equation}
\end{theorem}
\begin{proof}
Induction from \eqref{q1:qcc:q-weyl} gives
\(D_qX^k=[k]_qX^{k-1}+q^kX^kD_q\).
Left multiplication of the right side of \eqref{q1:qcc:normal-ordering-formula}
by \(XD_q\) therefore produces recurrence
\eqref{q1:qcc:normalized-stirling-rec}, with the correct initial value.
This proves normal ordering.
Apply both sides to \(x^m\). The left side is \([m]_q^nx^m\), and
\(X^kD_q^kx^m=[k]_q!\GB mk x^m\); terms with \(k>m\) vanish.
Comparison proves \eqref{q1:qcc:factorial-moments}.
\end{proof}
For \(n=2\), normal ordering reads
\((XD_q)^2=XD_q+qX^2D_q^2\). Thus
\(\widetilde S_q(2,2)=q\), while \(S_q(2,2)=1\). This tiny example is an
effective audit of the two normalizations.
% ---- hand chunk 11-affine-forms.tex ----

\section{Affine-geometric closed forms}
The nodes \([j]_q\) are affine-geometric: \([j]_q=A+Bq^j\) with
\(A=(1-q)^{-1}\) and \(B=-(1-q)^{-1}\). \Cref{q3:thm:affine} therefore
gives finite mixed ordinary/Gaussian sums for both arrays, in addition to the
symmetric-function forms of \cref{q1:qcc:qstirling-basic} and the spectral
sum of \cref{q1:qcc:qstirling-explicit}.

\begin{corollary}[Affine-geometric closed forms]\label{q3:cor:qstirling-gaussian}
For \(0\le k\le n\),
\begin{align}
 S_q(n,k)&=(1-q)^{k-n}
       \sum_{d=k}^n(-1)^{d-k}\binom nd\GB dk,
 \label{q3:eq:qstirling-gaussian}\\
 s_q(n,k)&=(1-q)^{k-n}
 \sum_{d=k}^n(-1)^{d-k}\binom dk
        q^{\binom{n-d}2}\GB nd.
 \label{q3:eq:qstirling-first-gaussian}
\end{align}
All apparent singularities at \(q=1\) are removable: both sides are
polynomials in \(q\).
\end{corollary}
\begin{proof}
Substituting \(A=(1-q)^{-1}\), \(B=-(1-q)^{-1}\) in
\eqref{q3:eq:affine-geometric} gives
\(A^{n-d}B^{d-k}=(-1)^{d-k}(1-q)^{k-n}\), which is
\eqref{q3:eq:qstirling-gaussian}. For the inverse array put
\(y=1-(1-q)x\), so that \(x-[j]_q=-(y-q^j)/(1-q)\) and
\(P_n(x)=(-1)^n(1-q)^{-n}\prod_{j=0}^{n-1}(y-q^j)\) in the notation of
\eqref{q1:qcc:qstirling-basis}. Insert
\eqref{q3:eq:gaussian-newton-inverse} for the product, expand
\(y^d=\sum_k\binom dk(-(1-q))^kx^k\) ordinarily, and compare the
coefficients of \(x^k\): the sign is
\((-1)^{n}(-1)^{n-d}(-1)^k=(-1)^{d-k}\) and the power of \(1-q\) is
\(k-n\). This is \eqref{q3:eq:qstirling-first-gaussian}. Polynomiality of
both arrays, from \eqref{q1:qcc:qstirling-basis}, proves the cancellation
claims.
\end{proof}

% ---- Q3 lines 772--816 ----
The first few rows are
\begin{equation}
\begin{array}{c|lllll}
 n\backslash k&0&1&2&3&4\\\hline
 0&1\\
 1&0&1\\
 2&0&1&1\\
 3&0&1&2+q&1\\
 4&0&1&3+3q+q^2&3+2q+q^2&1
\end{array}
\label{q3:eq:qstirling-table}
\end{equation}
In particular $S_q(n,1)=1$ for $n\ge1$. The limit $q\to1$ recovers
ordinary Stirling numbers, while $q=0$ gives
$S_0(n,k)=\binom{n-1}{k-1}$ for $1\le k\le n$: all positive-index
nodes have become one, and \cref{q1:qcc:Sq-h} counts weak
compositions of $n-k$ into $k$ positions.

\section{Colored and type-\texorpdfstring{$B$}{B} node specializations}
\begin{corollary}[Nodes $\qnum{r+mj}$]
\label{q3:cor:colored}
Fix integers $m\ge1$ and $r\ge0$, and define $S_q^{(m,r)}(n,k)$ using
the nodes $a_j=[r+mj]_q$. Then
\begin{equation}
 S_q^{(m,r)}(n,k)
 =(1-q)^{k-n}\sum_{d=k}^n
       \binom nd(-q^r)^{d-k}\Gb dk{q^m}.
 \label{q3:eq:colored}
\end{equation}
It is a polynomial with nonnegative integer coefficients. For $(m,r)=(2,1)$
its defining recurrence is
\begin{equation}
 S_q^{(2,1)}(n+1,k)=S_q^{(2,1)}(n,k-1)+[2k+1]_qS_q^{(2,1)}(n,k),
 \label{q3:eq:typeB}
\end{equation}
the type-$B$ convention discussed in \cite{dingzeng}.
\end{corollary}
\begin{proof}
Write $[r+mj]_q=(1-q)^{-1}-q^r(1-q)^{-1}(q^m)^j$ and use
\cref{q3:thm:affine}. Nonnegativity follows from the complete homogeneous
formula \cref{q3:eq:node-h}; \cref{q3:eq:typeB} is \cref{q3:eq:node-Srec}.
\end{proof}
No signed-permutation interpretation is needed to prove
\cref{q3:eq:colored}. It gives a finite Gaussian formula for the entire
node family, including its boundary values $S_q^{(m,r)}(n,0)=[r]_q^n$.
% ---- hand chunk 12-dobinski.tex ----

\section{A \qbname-Dobinski identity}
The factorial-moment identity \eqref{q1:qcc:factorial-moments} has a
generating-function form in which the normal-ordering normalization
\(\widetilde S_q\) is forced by the \(q\)-Weyl relation, not chosen by
analogy with the ordinary Dobinski formula.

\begin{theorem}[\(q\)-Dobinski]\label{q3:thm:qdobinski}
Let \(T_{n,q}(z)=\sum_{k=0}^n\widetilde S_q(n,k)z^k\) and let
\(e_q(z)=\sum_{j\ge0}z^j/[j]_q!\) be the \(q\)-exponential of
\eqref{q2:eq:qexp}. Then
\begin{equation}\label{q3:eq:qdobinski}
 T_{n,q}(z)=\frac1{e_q(z)}\sum_{j\ge0}[j]_q^n\frac{z^j}{[j]_q!}.
\end{equation}
This holds formally in \(z\) over \(\Q(q)\); for \(0<q<1\) it also holds
analytically when \(|(1-q)z|<1\).
\end{theorem}
\begin{proof}
Multiply \eqref{q1:qcc:factorial-moments} by \(z^m/[m]_q!\) and sum over
\(m\ge0\):
\[
 \sum_{m\ge0}[m]_q^n\frac{z^m}{[m]_q!}
 =\sum_{k=0}^n\widetilde S_q(n,k)[k]_q!
   \sum_{m\ge k}\frac{z^m}{[k]_q!\,[m-k]_q!}
 =\sum_{k=0}^n\widetilde S_q(n,k)z^k\,e_q(z),
\]
because \(\GB mk/[m]_q!=1/([k]_q!\,[m-k]_q!)\) and the \(k\)-sum is finite,
so the interchange is an identity of formal series. Dividing by the unit
\(e_q(z)\) gives \eqref{q3:eq:qdobinski}. Analytically, \([j]_q\) is bounded
by \((1-q)^{-1}\) for \(0<q<1\), so both series converge absolutely in the
stated disk.
\end{proof}
At \(q=1\) the identity is Dobinski's formula
\(\sum_kS(n,k)z^k=e^{-z}\sum_jj^nz^j/j!\). For \(n=2\) it reads
\(z+qz^2=e_q(z)^{-1}\sum_j[j]_q^2z^j/[j]_q!\), which is the normal-ordering
audit \(\widetilde S_q(2,2)=q\) once more.

% ---- Q3 lines 868--977 ----
\chapter{\texorpdfstring{$q$}{q}-Eulerian polynomials and the Stirling--Eulerian bridge}
\label{q3:sec:eulerian}
For a permutation $\pi=(\pi_1,\ldots,\pi_n)$, let
$\operatorname{Des}(\pi)=\{i:\pi_i>\pi_{i+1}\}$,
$\des(\pi)=|\operatorname{Des}(\pi)|$, and
$\maj(\pi)=\sum_{i\in\operatorname{Des}(\pi)}i$. Define
\begin{equation}
 A_n(t,q)=\sum_{\pi\in\mathfrak S_n}t^{\des(\pi)}q^{\maj(\pi)}
       =\sum_{j=0}^{n-1}A_{n,j}(q)t^j \qquad(n\ge1).
 \label{q3:eq:qeulerian}
\end{equation}
The major index, rather than an unspecified inversion statistic, is
important in the identities below.

\begin{theorem}[$q$-Worpitzky and a finite Eulerian sum]
\label{q3:thm:worpitzky}
For $n\ge1$,
\begin{equation}
 \sum_{m\ge0}[m+1]_q^n t^m=\frac{A_n(t,q)}{(t;q)_{n+1}}.
 \label{q3:eq:qeulerian-gf}
\end{equation}
Consequently, for $m\ge0$,
\begin{align}
 [m+1]_q^n&=\sum_{j=0}^{n-1}A_{n,j}(q)\G{m+n-j}n,
 \label{q3:eq:qworpitzky}\\
 A_{n,j}(q)&=\sum_{r=0}^j(-1)^rq^{\binom r2}\G{n+1}r[j-r+1]_q^n.
 \label{q3:eq:qeulerian-explicit}
\end{align}
\end{theorem}
\begin{proof}
Expand $[m+1]_q^n$ as the weight sum of vectors
$(a_1,\ldots,a_n)\in\{0,\ldots,m\}^n$, with weight $q^{\sum a_i}$.
Sort the labels by decreasing $a_i$, resolving equal entries by
increasing label. This gives a unique permutation $\pi$ such that
$a_{\pi_1}\ge\cdots\ge a_{\pi_n}$, with a strict inequality at every
descent of $\pi$. Put $d=\des(\pi)$ and subtract from the $i$th sorted
entry the number of descents at positions at least $i$. The resulting
sequence is
$ m-d\ge b_1\ge\cdots\ge b_n\ge0$, and its weight is shifted by
$\maj(\pi)$. Conversely, this construction reconstructs exactly the
vectors producing $\pi$.
Write $e_0=m-d-b_1$, $e_i=b_i-b_{i+1}$ for $1\le i<n$, and $e_n=b_n$.
Then $m=d+\sum_{i=0}^ne_i$ and $\sum b_i=\sum_{i=1}^ni e_i$.
The generating function for these gaps is
\[
 t^dq^{\maj(\pi)}\prod_{i=0}^n\frac1{1-tq^i}.
\]
Sum over permutations to get \cref{q3:eq:qeulerian-gf}. Extract $t^m$
using \cref{q1:qcc:reciprocal} for \cref{q3:eq:qworpitzky}; multiply by
$(t;q)_{n+1}$ and use \cref{q1:qcc:finite} for
\cref{q3:eq:qeulerian-explicit}.
\end{proof}

\begin{theorem}[Gaussian Stirling--Eulerian transforms]
\label{q3:thm:stirling-eulerian}
For $1\le k\le n$,
\begin{equation}
 q^{\binom k2}\qfac k S_q(n,k)
 =\sum_{\ell=1}^k q^{k(k-\ell)}A_{n,\ell-1}(q)
       \G{n-\ell}{k-\ell}.
 \label{q3:eq:stirling-eulerian}
\end{equation}
Conversely,
\begin{equation}
 A_n(t,q)=\sum_{k=1}^n
 q^{\binom k2}\qfac kS_q(n,k)
       t^{k-1}(tq^{k+1};q)_{n-k}.
 \label{q3:eq:eulerian-stirling}
\end{equation}
\end{theorem}
\begin{proof}
Set $U_k=q^{\binom k2}\qfac kS_q(n,k)$. The Newton expansion at
$x=[m+1]_q$ gives
$[m+1]_q^n=\sum_{k=1}^nU_kG_{m+1,k}(q)$.
Since
\[
 \sum_{m\ge0}G_{m+1,k}(q)t^m=\frac{t^{k-1}}{(t;q)_{k+1}},
\]
comparison with \cref{q3:eq:qeulerian-gf} proves
\cref{q3:eq:eulerian-stirling}.
To invert it, we first establish the finite identity
\begin{equation}
 \sum_{r=0}^N\G Nr q^{r(r-1)}u^r(uq^r;q)_{N-r}=1.
 \label{q3:eq:eulerian-kernel}
\end{equation}
The coefficient of $u^j$ on the left equals
\[
 \G Nj q^{\binom j2}
   \sum_{r=0}^j(-1)^{j-r}q^{\binom r2}\G jr,
\]
which is $\delta_{j,0}$ by \cref{q3:cor:annihilator}. Set
$N=n-\ell$ and $u=tq^{\ell+1}$ in \cref{q3:eq:eulerian-kernel} and
multiply by $t^{\ell-1}$. This expresses $t^{\ell-1}$ as
\[
 \sum_{k=\ell}^n q^{k(k-\ell)}\G{n-\ell}{k-\ell}
          t^{k-1}(tq^{k+1};q)_{n-k}.
\]
Insert this expansion for each monomial of $A_n(t,q)$ in
\cref{q3:eq:eulerian-stirling}. These basis polynomials have distinct
lowest powers, so their coefficients are unique. The resulting
coefficient of the $k$th basis polynomial is \cref{q3:eq:stirling-eulerian}.
\end{proof}
Identity \cref{q3:eq:stirling-eulerian} is classical; its normalization and
history are recorded in \cite{dingzeng}. The proof here exhibits it as
another triangular coefficient transform, directly parallel to the
ordinary Stirling--Eulerian relations in the canonical manuscript.
For example $A_3(t,q)=1+(2q+2q^2)t+q^3t^2$, and at $(n,k)=(3,2)$
both sides of \cref{q3:eq:stirling-eulerian} equal
$2q+3q^2+q^3$.

% ---- Q1 lines 1026--1108 ----
\chapter{Weighted Bell polynomials and Riordan composition}
\label{q1:qcc:composition}
\section{Gaussian convolution as a change of coefficient coordinates}
Set \(w_n=[n]_q!\), and record a sequence by
\[
 A(z)=\sum_{n\ge0}a_n\frac{z^n}{w_n}.
\]
Ordinary multiplication of series is then the Gaussian convolution
\begin{equation}\label{q1:qcc:q-convolution}
 (a\star_q b)_n=\sum_{k=0}^n\GB nk a_kb_{n-k}.
\end{equation}
Associativity and commutativity follow from ordinary series multiplication.
More explicitly, for \(s\) factors the convolution coefficient is
\begin{equation}\label{q1:qcc:q-multinomial}
 \sum_{i_1+\cdots+i_s=n}
 \frac{[n]_q!}{[i_1]_q!\cdots[i_s]_q!}
 a^{(1)}_{i_1}\cdots a^{(s)}_{i_s}.
\end{equation}
The factorial ratio is the Gaussian multinomial. It is polynomial because it is
the product
\(\GB n{i_1}\GB{n-i_1}{i_2}\cdots\).
Thus associativity extends polynomially to every specialization of \(q\), even
when the divided-power representation itself is not available.

\begin{proposition}[Explicit convolution inverse]\label{q1:qcc:convolution-inverse}
If \(a_0\) is invertible, the inverse \(b\) for \(\star_q\) satisfies
\(b_0=a_0^{-1}\), and for \(n>0\),
\begin{equation}\label{q1:qcc:convolution-inverse-formula}
 b_n=\frac1{a_0}\sum_{r=1}^n\frac{(-1)^r}{a_0^r}
 \sum_{\substack{i_1+\cdots+i_r=n\\i_j\ge1}}
 \frac{[n]_q!}{[i_1]_q!\cdots[i_r]_q!}\prod_{j=1}^r a_{i_j}.
\end{equation}
\end{proposition}
\begin{proof}
Write \(A(z)=a_0(1+U(z))\), where \(U(0)=0\). The formal geometric expansion
\(A(z)^{-1}=a_0^{-1}\sum_{r\ge0}(-U(z))^r\) is coefficientwise finite.
Multiplying its degree-\(n\) coefficient by \(w_n\) gives the displayed formula.
\end{proof}
In particular, with
\[
 e_q(z)=\sum_{n\ge0}\frac{z^n}{[n]_q!},\qquad
 E_q(z)=\sum_{n\ge0}q^{\binom n2}\frac{z^n}{[n]_q!},
\]
Gaussian inversion is the coefficient identity \(e_q(z)E_q(-z)=1\).

\section{Composition coefficients without an ambiguous ``\qbname-Bell'' convention}
Let
\(B(z)=\sum_{j\ge1}b_jz^j/w_j\). Define
\begin{equation}\label{q1:qcc:weighted-bell-def}
 \qB nk(b)=\frac{w_n}{w_k}[z^n]B(z)^k.
\end{equation}
This is a specified weighted coefficient family; it is not a claim that all
objects called \(q\)-Bell polynomials in the literature have this meaning.

\begin{theorem}\label{q1:qcc:weighted-bell}
For \(0\le k\le n\),
\begin{equation}\label{q1:qcc:weighted-bell-finite}
 \qB nk(b)=\frac{w_nk!}{w_k}
 \sum_{\substack{m_1,\ldots,m_n\ge0\\\sum m_j=k,\ \sum j m_j=n}}
 \prod_{j=1}^n\frac1{m_j!}\left(\frac{b_j}{w_j}\right)^{m_j}.
\end{equation}
Equivalently,
\begin{equation}\label{q1:qcc:weighted-bell-transport}
 \qB nk(b)=\frac{w_nk!}{w_kn!}\B nk
       \left(\frac{1!b_1}{w_1},\frac{2!b_2}{w_2},\ldots\right).
\end{equation}
If \(A(z)=\sum_{k\ge0}a_kz^k/w_k\), then
\begin{equation}\label{q1:qcc:composition-formula}
 w_n[z^n]A(B(z))=\sum_{k=0}^n a_k\qB nk(b).
\end{equation}
\end{theorem}
\begin{proof}
Expand \(B(z)^k\) by the ordinary multinomial theorem. A multiplicity vector
\((m_j)\) contributes \(k!\prod_j(b_j/w_j)^{m_j}/m_j!\).
The two constraints count the factors and their total degree. This proves
\eqref{q1:qcc:weighted-bell-finite}. Compare it with the defining finite sum for
\(\B nk\) to obtain \eqref{q1:qcc:weighted-bell-transport}. Finally, expand
\(A(B)\) in powers of \(B\) and use \eqref{q1:qcc:weighted-bell-def}.
\end{proof}
At \(q=1\), \(w_n=n!\), and \(\mathfrak B^{[q]}_{n,k}\) reduces to the
ordinary exponential Bell polynomial \(\B nk\). At \(q=0\), \(w_n=1\),
and it becomes the ordinary coefficient \([z^n]B(z)^k\). This cleanly separates
coefficient normalization from the choice of a deformation of set partitions.
% ---- Q2 lines 857--869 ----
The first nontrivial examples expose a frequent normalization error:
\begin{align}
 \qB32&=\frac{2\qi3}{\qi2}b_1b_2,\label{q2:eq:Bq32}\\
 \qB42&=\frac{2\qi4}{\qi2}b_1b_3
             +\frac{\qi4\qi3}{(\qi2)^2}b_2^2,\label{q2:eq:Bq42}\\
 \qB43&=\frac{3\qi4}{\qi2}b_1^2b_2.
 \label{q2:eq:Bq43}
\end{align}
Also $\qB n1=b_n$ and $\qB nn=b_1^n$. The rational coefficient
$2\qi3/\qi2$ is not $\qi3$. It need not be a polynomial in $q$, and it may have a
pole at a root of unity. This is expected: the coordinate system itself uses
inverses of $q$-factorials.

% ---- Q1 lines 1110--1139 ----
\section{Weighted Riordan arrays}
For \(A(0)\ne0\) and \(B(0)=0\) with invertible linear coefficient, put
\begin{equation}\label{q1:qcc:riordan-def}
 \mathcal M_q(A,B)_{n,k}=\frac{w_n}{w_k}[z^n]A(z)B(z)^k.
\end{equation}
\begin{theorem}\label{q1:qcc:riordan-group}
These lower triangular matrices obey
\begin{align}
 \mathcal M_q(A,B)\mathcal M_q(C,D)
 &=\mathcal M_q\bigl(A\,C(B),D(B)\bigr),\label{q1:qcc:riordan-product}\\
 \mathcal M_q(A,B)^{-1}
 &=\mathcal M_q\left(\frac1{A(\overline B)},\overline B\right),
 \label{q1:qcc:riordan-inverse}
\end{align}
where \(\overline B\) is the ordinary compositional inverse of \(B\).
\end{theorem}
\begin{proof}
The \((n,k)\) entry of the product equals
\[
 \frac{w_n}{w_k}\sum_j[z^n]A(z)B(z)^j\,[t^j]C(t)D(t)^k.
\]
For fixed \(n\), the sum is finite because \(B(0)=0\). Summing before taking
\([z^n]\) produces \(A(z)C(B(z))D(B(z))^k\), proving the product law.
Inserting \(D=\overline B\) and \(C=1/A(\overline B)\) gives
\(\mathcal M_q(1,z)\), the identity matrix. The same substitution in the other
order gives the other inverse identity.
\end{proof}
This is diagonal transport of ordinary Riordan calculus. It extends the canonical manuscript's
Riordan framework while making clear which operation is unchanged: composition
is still ordinary substitution of formal power series.
% ---- Q3 lines 1433--1443 ----
As a useful check,
\[
 \mathcal M_q(e_q(z),z)_{n,k}=\G nk,
\]
and its inverse is $\mathcal M_q(E_q(-z),z)$. This recovers Gaussian inversion,
including the factor $(-1)^{n-k}q^{\binom{n-k}2}$.
If polynomials $P_n(x)$ are defined by
$g(z)e_q(xz)=\sum P_n(x)z^n/[n]_q!$, then
$\Dq P_n(x)=\qnum n P_{n-1}(x)$: apply $\Dq$ in $x$ to the generating
function and compare coefficients. This supplies the corresponding
$q$-Appell lowering property with no ambiguity about normalization.
% ---- Q1 lines 1141--1152 ----
\section{Why Jackson composition is a different problem}
The first-order chain rule for the Jackson derivative is
\begin{equation}\label{q1:qcc:divided-chain}
 D_q(F\circ G)(x)=F[G(x),G(qx)]\,D_qG(x),
\end{equation}
where \(F[u,v]=(F(u)-F(v))/(u-v)\), with its polynomial continuation on the
diagonal. This follows simply by factoring the numerator in
\eqref{q1:qcc:jackson-def}. It is not usually \((D_qF)(G(x))D_qG(x)\).
For example, with \(F(u)=u^2\) and \(G(x)=x^2\), the actual answer is
\([4]_qx^3\), whereas that guess gives \([2]_q^2x^3\).
The weighted Bell family in this chapter describes ordinary composition and
must not be silently substituted into a higher Jackson chain rule.
% ---- hand chunk 16-chain-xref.tex ----

The exact higher-order rule is the path formula of \cref{q3:sec:chain}: the
\(n\)th Jackson derivative of a composition is a sum over increasing paths of
products of divided differences, not a substitution of \([j]_q!\) for
\(j!\) in Fa\`a di Bruno's formula.


% ---- Q1 lines 1154--1259 ----
\chapter{Lagrange inversion with geometric product data}
\label{q1:qcc:lagrange}
\section{The ordinary inversion theorem over a \qbname-dependent field}
\begin{theorem}[Formal Lagrange--B\"urmann]\label{q1:qcc:lagrange-theorem}
Let \(\Phi(0)=1\). There is a unique \(y(z)\in zK[[z]]\) satisfying
\(y=z\Phi(y)\). For \(n\ge1\),
\begin{equation}\label{q1:qcc:lagrange-formula}
 [z^n]H(y(z))=\frac1n[t^{n-1}]H'(t)\Phi(t)^n.
\end{equation}
In particular, for \(s\ge1\) and \(n\ge s\),
\begin{equation}\label{q1:qcc:lagrange-powers}
 [z^n]y(z)^s=\frac{s}{n}[t^{n-s}]\Phi(t)^n.
\end{equation}
\end{theorem}
\begin{proof}
The equation determines the coefficient of degree \(n\) of \(y\) from lower
degrees, starting with \([z]y=1\). This proves existence and uniqueness.
The inverse substitution is \(z=t/\Phi(t)\).
For a Laurent series, residue means the coefficient of degree \(-1\).
The residue of a formal derivative is zero, and the substitution rule
\(\Res_z A(u(z))u'(z)\,dz=\Res_u A(u)\,du\) follows on monomials:
for powers other than \(-1\), the integrand is a derivative; for the power
\(-1\), the logarithmic derivative has residue one.
Thus
\begin{align*}
 [z^n]H(y(z))
 &=\Res_t H(t)\left(\frac{\Phi(t)^n}{t^{n+1}}
                   -\frac{\Phi(t)^{n-1}\Phi'(t)}{t^n}\right)\dd t\\
 &=\frac1n\Res_t H'(t)\frac{\Phi(t)^n}{t^n}\dd t.
\end{align*}
The second equality is obtained by taking the zero residue of the derivative
of \(H(t)\Phi(t)^nt^{-n}\). It is \eqref{q1:qcc:lagrange-formula}; taking
\(H(t)=t^s\) proves \eqref{q1:qcc:lagrange-powers}.
\end{proof}
This is the ordinary theorem used in the source, now over a field containing
\(q\); see also the proof-oriented survey \cite{gessel}.

\section{A recurrence-free inverse for a product of shifted factorials}
\begin{theorem}\label{q1:qcc:inverse-product}
Suppose \(\Phi(t)\) is the product \eqref{q1:qcc:master-product}, with logarithmic
power sums \(P_r\). The solution of \(y=z\Phi(y)\) has coefficients
\begin{equation}\label{q1:qcc:inverse-product-formula}
 [z^n]y(z)=\frac1{n!}\Y{n-1}
       \left(\left(n(r-1)!P_r\right)_{r=1}^{n-1}\right)\qquad(n\ge1).
\end{equation}
More generally,
\begin{equation}\label{q1:qcc:inverse-product-powers}
 [z^n]y(z)^s=\frac{s}{n(n-s)!}\Y{n-s}
       \left(\left(n(r-1)!P_r\right)_{r=1}^{n-s}\right),\qquad n\ge s\ge1.
\end{equation}
\end{theorem}
\begin{proof}
The logarithm of \(\Phi(t)^n\) has power sums \(nP_r\).
Apply \cref{q1:qcc:master-theorem} inside \eqref{q1:qcc:lagrange-powers}.
The case \(s=1\) has denominator \(n(n-1)!=n!\).
\end{proof}
This solves, coefficient by coefficient and without auxiliary recurrence, the
inverse problem for
\[
 f(x)=x\prod_{\ell=1}^L(a_\ell x;q^{d_\ell})_{N_\ell}^{\alpha_\ell}.
\]
For \(f(x)=x(x;q)_M^\alpha\), the formula becomes
\begin{equation}\label{q1:qcc:simple-inverse}
 [z^n]f^{-1}(z)=\frac1{n!}\Y{n-1}
   \left(\left(\alpha n(r-1)![M]_{q^r}\right)_{r=1}^{n-1}\right).
\end{equation}
An alternative finite sum, obtained by expanding each factor separately, is
\begin{equation}\label{q1:qcc:simple-inverse-multisum}
 [z^n]f^{-1}(z)=\frac1n
 \sum_{\substack{r_0,\ldots,r_{M-1}\ge0\\\sum r_j=n-1}}
 \prod_{j=0}^{M-1}\frac{\rising{\alpha n}{r_j}}{r_j!}
 q^{\sum j r_j}.
\end{equation}
Here the rising factorial makes the formula polynomial in \(\alpha\), including
values at which gamma-quotient representations would be inconvenient.

\begin{example}[A two-color inverse]\label{q1:qcc:inverse-example}
The inverse of \(x(1-x)(1-qx)\) begins
\begin{align*}
 y(z)={}&z+(1+q)z^2+(2+3q+2q^2)z^3\\
       &+(5+10q+10q^2+5q^3)z^4+O(z^5).
\end{align*}
At \(q=1\), the coefficient formula is
\(n^{-1}\binom{3n-2}{n-1}\), giving \(1,2,7,30,\ldots\).
\end{example}
The coefficients for positive integral \(\alpha\) count rooted plane trees
whose children are split into \(\alpha M\) ordered blocks; a child in a block of
color \(j\) has weight \(q^j\). Each vertex contributes one factor
\(\Phi(t)\), and the root contributes \(z\), proving the functional equation.
Every tree with \(n\) vertices has \(n-1\) edges. Reversing all color labels
\(j\mapsto M-1-j\) therefore proves
\begin{equation}\label{q1:qcc:inverse-reciprocity}
 [z^n]y(z;q^{-1})=q^{-(M-1)(n-1)}[z^n]y(z;q).
\end{equation}
The same identity follows algebraically by reversing the finite product and using
homogeneity in \eqref{q1:qcc:simple-inverse}. The tree argument additionally proves
nonnegativity and integrality in this parameter range.

\begin{auditbox}[title={Ordinary inversion is not obtained by replacing \(n\) with \([n]_q\)}]
The factor \(1/n\) in Lagrange inversion is an ordinary integer factor. For
\(M=1\), the equation is \(y=z/(1-y)\), independent of \(q\), and its
coefficients are Catalan numbers. Replacing \(1/n\) by \(1/[n]_q\) would make
the answer depend on \(q\), contradicting the defining equation. Genuine
\(q\)-shifted inversion theories change additional algebraic structure; they
are not the operation performed here.
\end{auditbox}
% ---- Q3 lines 1514--1575 ----
\section{Optional-slot trees and a Gaussian divisibility identity}
A second choice of input reveals polynomial positivity that is not
obvious from Lagrange's factor $1/n$.

\begin{theorem}[Gaussian slot-tree coefficients]\label{q3:thm:slot-tree}
Fix an integer $N\ge1$ and define $T_N(z)$ by
\begin{equation}
 T_N(z)=z\prod_{j=0}^{N-1}(1+q^jT_N(z)).
 \label{q3:eq:slot-tree-eq}
\end{equation}
For $n\ge1$, its coefficient $C_{N,n}(q)=\coef zn T_N(z)$ satisfies
\begin{equation}
 C_{N,n}(q)=\frac1n
   \sum_{\substack{k_1+\cdots+k_n=n-1\\0\le k_i\le N}}
        q^{\sum_{i=1}^n\binom{k_i}2}
        \prod_{i=1}^n\G N{k_i}.
 \label{q3:eq:slot-gaussian}
\end{equation}
Moreover,
\begin{align}
 C_{N,n}(q)&\in\Z_{\ge0}[q],
 \label{q3:eq:slot-positive}\\
 q^{(N-1)(n-1)}C_{N,n}(q^{-1})&=C_{N,n}(q),
 \label{q3:eq:slot-reciprocal}\\
 C_{N,n}(1)&=\frac1n\binom{Nn}{n-1}
 =\frac1{(N-1)n+1}\binom{Nn}{n}.
 \label{q3:eq:slot-fuss}
\end{align}
Thus the numerator of \cref{q3:eq:slot-gaussian} is divisible by $n$
coefficientwise in $\Z[q]$.
\end{theorem}
\begin{proof}
Use Lagrange inversion with $\Phi(w)=(-w;q)_N$, expand each of the $n$
factors by the finite $q$-binomial theorem, and extract degree $n-1$.
This proves \cref{q3:eq:slot-gaussian}.
For integrality and positivity, consider rooted trees in which each
vertex has $N$ labeled child slots $0,\ldots,N-1$, each either empty or
occupied by one subtree. Give an edge in slot $j$ weight $q^j$, and give
each vertex weight $z$. Decomposing at the root gives
\cref{q3:eq:slot-tree-eq}. There are finitely many such trees of each size;
their weight sum is a polynomial with nonnegative integer coefficients.
Uniqueness of the formal solution identifies it with $T_N$.

Reverse every slot label $j\mapsto N-1-j$. A tree of $n$ vertices has
$n-1$ edges, so its exponent $w$ changes to $(N-1)(n-1)-w$. This
weight-reversing bijection proves reciprocity. Finally put $q=1$ in
Lagrange's formula: $\Phi(w)^n=(1+w)^{Nn}$. The equality of the two
ordinary binomial expressions is a direct factorial cancellation.
\end{proof}

For $N=2$ an especially simple form is
\begin{equation}
 C_{2,n}(q)=\frac1n\sum_{j=0}^{n-1}
                \binom nj\binom n{j+1}q^j.
 \label{q3:eq:narayana-slot}
\end{equation}
For example $C_{2,3}=1+3q+q^2$. The formula follows by extracting
$w^{n-1}$ from $(1+w)^n(1+qw)^n$.
These are Narayana polynomials under the displayed normalization.
They are not the same $q$-Catalan family as a Gaussian quotient such as
$\G{2n}n/[n+1]_q$. The name ``$q$-Catalan'' by itself does not specify
a weight statistic or a unique deformation.
% ---- Q2 lines 1047--1056 ----
For a reciprocal example, take
$w=z/\poch{aw}qN^r$ with integers $r,N\ge1$. Then
\begin{equation}\label{q2:eq:reciprocal-reversion}
 [z^n]w^m=\frac mn a^{n-m}
 \sum_{\substack{k_1+\cdots+k_{rn}=n-m\\k_i\ge0}}
       \prod_{i=1}^{rn}\qb{N+k_i-1}{k_i}.
\end{equation}
Here \cref{q1:qcc:finite-binomial} replaces the finite theorem. Formula
\eqref{q1:qcc:inverse-product-formula} continues to apply when $r$ is not an integer, whereas
\eqref{q2:eq:reciprocal-reversion} uses $rn$ separate factors and therefore does not.
% ---- Q2 lines 1058--1116 ----
\section{A genuinely shifted equation: the area-type \texorpdfstring{$q$}{q}-Catalan series}\label{q2:sec:qcatalan}
The equation in \cref{q3:thm:slot-tree} contains $q$-dependent coefficients but no dilation
of its unknown series. A true $q$-functional equation behaves differently. The
existence of several inequivalent $q$-Catalan families is well documented;
\cite{cigler} develops a uniform viewpoint. We give a complete derivation of the
particular normalization needed here.

\begin{theorem}[A quotient solution and a finite coefficient formula]\label{q2:thm:qcat}
Over $\Q(q)$ define
\begin{equation}\label{q2:eq:U-Airy}
 U(z)=\sum_{n\ge0}u_nz^n,\qquad
 u_n=\frac{(-1)^nq^{n(n-1)}}{D_n}.
\end{equation}
Then
\begin{equation}\label{q2:eq:C-Airy}
 C(z)=\frac{U(qz)}{U(z)}
\end{equation}
is the unique formal solution with constant coefficient $1$ of
\begin{equation}\label{q2:eq:area-eqn}
 C(z)=1+zC(z)C(qz).
\end{equation}
For each $n\ge0$, its coefficients have the nonrecursive finite expression
\begin{equation}\label{q2:eq:area-finite}
 [z^n]C(z)=\sum_{k=0}^n q^k u_k
   \sum_{\bm m\in\Pcal(n-k)}
       (-1)^{M(\bm m)}M(\bm m)!
           \prod_{j=1}^{n-k}\frac{u_j^{m_j}}{m_j!},
 \qquad M(\bm m)=\sum_jm_j.
\end{equation}
\end{theorem}
\begin{proof}
The coefficients in \eqref{q2:eq:U-Airy} satisfy, for $n\ge1$,
$(1-q^n)u_n=-q^{2n-2}u_{n-1}$, as follows immediately by cancelling $D_n$.
Consequently $U(z)-U(qz)=-zU(q^2z)$. Division by the unit $U(z)$ gives
\[
 \frac{U(qz)}{U(z)}=1+z\frac{U(q^2z)}{U(qz)}\frac{U(qz)}{U(z)},
\]
which proves \eqref{q2:eq:area-eqn}. Its coefficient of degree $n\ge1$ involves
only lower-degree coefficients on the right, so the solution is unique.
Finally expand $1/U=\sum_{r\ge0}(-1)^r(U-1)^r$ by the ordinary multinomial
theorem and multiply by $U(qz)$. This is exactly \eqref{q2:eq:area-finite}.
\end{proof}

The first terms are
\begin{equation}\label{q2:eq:area-first}
 C(z)=1+z+(1+q)z^2+(1+2q+q^2+q^3)z^3+O(z^4).
\end{equation}
Equation \eqref{q2:eq:area-eqn}, or its binary-tree interpretation, shows that every
coefficient belongs to $\Z_{\ge0}[q]$. Indeed, attach two ordered subtrees to a root
and give the second subtree an extra factor $q$ for each of its vertices. The
resulting weight rule is precisely $zC(z)C(qz)$. This proves that singularities
in the rational summands of \eqref{q2:eq:area-finite} cancel, including at $q=1$.
At $q=1$ the equation is $C=1+zC^2$, so the usual Catalan sequence is recovered.

The polynomials in \eqref{q2:eq:area-first}, the slot family in \eqref{q3:eq:narayana-slot},
and the Gaussian quotient $\qb{2n}n/\qi{n+1}$ are not interchangeable.
For example, their usual Catalan index $n=2$ gives respectively $1+q$,
$1+q$, and $1+q^2$; at $n=3$ the first two already differ. The implicit equation
and normalization, not the name ``$q$-Catalan'', must select the intended family.
% ---- Q1 lines 1261--1308 ----
\section{A two-variable inverse with its Jacobian factor exposed}
The source also develops multivariate Good inversion. The following complete
specialization shows how geometric-product Bell data enter a coupled system,
while retaining the correction that a naive product of one-variable formulae
would omit.

\begin{theorem}[Cross-coupled inversion]\label{q1:qcc:coupled}
Let \(A(0)=B(0)=1\), and let
\[
 u=z_1A(v),\qquad v=z_2B(u).
\]
For \(m,n>0\), \(r,s\ge0\), \(m\ge r\), and \(n\ge s\),
\begin{equation}\label{q1:qcc:coupled-formula}
 [z_1^mz_2^n]u^rv^s
 =\frac{ms+nr-rs}{mn}\,
 [t^{n-s}]A(t)^m\,[t^{m-r}]B(t)^n.
\end{equation}
Each coefficient on the right has the Bell form
\eqref{q1:qcc:master-coefficient} when \(A,B\) are finite geometric products.
\end{theorem}
\begin{proof}
Substitute the first equation into the second:
\(v=z_2B(z_1A(v))\). Apply the one-variable formula
\eqref{q1:qcc:lagrange-formula} in \(z_2\), with \(z_1\) a parameter, to
\(u^rv^s=z_1^rA(v)^rv^s\). The derivative is
\(s t^{s-1}A(t)^r+r t^sA(t)^{r-1}A'(t)\).
Furthermore,
\[
 [z_1^{m-r}]B(z_1A(t))^n
 =A(t)^{m-r}[w^{m-r}]B(w)^n.
\]
The first derivative term therefore contributes
\(\frac{s}{n}[t^{n-s}]A(t)^m\) times the \(B\)-coefficient. The second
contributes
\[
 \frac{r}{n}[t^{n-s-1}]A(t)^{m-1}A'(t)
 =\frac{r(n-s)}{mn}[t^{n-s}]A(t)^m.
\]
Adding gives \((ms+nr-rs)/(mn)\). This argument also covers \(r=0\) or
\(s=0\), since the corresponding derivative term is zero.
\end{proof}
For comparison, the determinant in the two-variable Good formula is
\[
 1-\frac{t_2A'(t_2)}{A(t_2)}\frac{t_1B'(t_1)}{B(t_1)}.
\]
Its correction subtracts \((n-s)(m-r)/(mn)\) from one, which is exactly the
prefactor in \eqref{q1:qcc:coupled-formula}. This gives an independent check on the
indices and on the otherwise easily omitted Jacobian contribution.
% ---- Q2 lines 1177--1190 ----
\section{The Good determinant behind the formula}
For a general system $w_i=z_i\Phi_i(\bm w)$ with unit $\Phi_i$, the ordinary
multivariate Lagrange--Good formula reads
\begin{equation}\label{q2:eq:Good}
 [\bm z^{\bm n}]H(\bm w)
 =[\bm t^{\bm n}]H(\bm t)\prod_i\Phi_i(\bm t)^{n_i}
   \det\left(\delta_{ij}-t_j\partial_{t_j}\log\Phi_i(\bm t)\right).
\end{equation}
It follows by the formal change of variables $z_i=t_i/\Phi_i(\bm t)$ in the
multivariate residue: the displayed matrix is exactly the Jacobian on the
logarithmic differentials $dt_j/t_j$. The residue invariance is the
multivariate counterpart of the substitution in \cref{q1:qcc:lagrange-theorem}, and the
formula is one of the canonical manuscript's central tools \cite{ccc}.

% ---- hand chunk 19-good-glue.tex ----

For the system of \cref{q1:qcc:coupled} the determinant is
\[
 1-uv\frac{A'(v)}{A(v)}\frac{B'(u)}{B(u)}.
\]
Extracting the two factors gives \(1-(m-r)(n-s)/(mn)\) times the product
of the two one-variable coefficients, which is the prefactor
\((ms+nr-rs)/(mn)\) of \eqref{q1:qcc:coupled-formula}. The proof given
there derives this special case without invoking a multivariate residue
theorem.


% ---- Q2 lines 1200--1214 ----
More generally, a factor
$\poch{a\bm t^{\bm\nu}}{q^h}N^\alpha$ has
\begin{equation}\label{q2:eq:Good-entry}
 t_j\partial_{t_j}\log\poch{a\bm t^{\bm\nu}}{q^h}N^\alpha
  =-\alpha\nu_j\sum_{\ell\ge1}
       a^\ell S_N(q^{h\ell})\bm t^{\ell\bm\nu}.
\end{equation}
Here $\bm\nu$ is a nonzero nonnegative integer vector. Thus, after imposing a
fixed multidegree bound, both the determinant and its product factors have
explicit finite expansions: expand the determinant over permutations and expand
the exponential of the logarithmic product over multiplicities indexed by
$(\bm\nu,\ell)$. Only those with $\ell\bm\nu\le\bm n$ can contribute.
This gives a direct route from the canonical manuscript's Good theorem to multivariate
$q$-shifted-factorial identities without changing its Jacobian or introducing
unjustified Jackson derivatives.
% ---- Q1 lines 1310--1398 ----
\chapter{The coefficients of a Gaussian polynomial}
\label{q1:qcc:gaussian-coefficients}
\section{A divisor-sum formula for each coefficient}
Write \(m=n-k\), \(D=k(n-k)\), and
\[
 G_{n,k}(q)=\GB nk=\sum_{r=0}^D c_{n,k}(r)q^r.
\]
This chapter extracts coefficients in the \emph{base} \(q\), rather than in an
auxiliary product variable. Since \(G_{n,k}(0)=1\), it has a formal logarithm
in \(\Q[[q]]\).

\begin{theorem}\label{q1:qcc:divisor-coefficients}
For \(r\ge1\), define the finite divisor sum
\begin{equation}\label{q1:qcc:arithmetic-kernel}
 A_r(n,k)=\sum_{\substack{d\mid r\\1\le d\le k}}d
          -\sum_{\substack{d\mid r\\m<d\le n}}d.
\end{equation}
Then
\begin{equation}\label{q1:qcc:gaussian-coeff-bell}
 c_{n,k}(r)=\frac1{r!}\Y r\left(
        \left((j-1)!A_j(n,k)\right)_{j=1}^r\right),
\end{equation}
or, with no named auxiliary polynomial,
\begin{equation}\label{q1:qcc:gaussian-coeff-finite}
 c_{n,k}(r)=\sum_{\substack{v_1,\ldots,v_r\ge0\\\sum jv_j=r}}
                 \prod_{j=1}^r\frac{A_j(n,k)^{v_j}}{j^{v_j}v_j!}.
\end{equation}
Also \(c_{n,k}(0)=1\), and the formula gives zero for \(r>D\).
\end{theorem}
\begin{proof}
The finite product quotient gives
\begin{align*}
 \log G_{n,k}(q)
 &=\sum_{d=1}^k\sum_{a\ge1}\frac{q^{da}}a
   -\sum_{d=m+1}^n\sum_{a\ge1}\frac{q^{da}}a\\
 &=\sum_{r\ge1}\frac{A_r(n,k)}r q^r.
\end{align*}
For fixed \(r\), only divisors \(d\mid r\) contribute, with coefficient
\(1/(r/d)=d/r\). The Bell exponential formula proves the result. The original
series is a polynomial of degree \(D\), proving the last assertion.
\end{proof}
Although individual terms in \eqref{q1:qcc:gaussian-coeff-finite} are rational and
can have signs, the answer is a nonnegative integer by the subset definition.
For example, when \((n,k)=(4,2)\), the first kernels are
\(A_1=1,A_2=3,A_3=-2,A_4=-1\). They yield
\[
 c(1)=1,\qquad c(2)=\frac{1+3}{2}=2,\qquad
 c(3)=\frac{1+9-4}{6}=1,
\]
and \(c(4)=1\), as expected.

\section{A second finite formula by inclusion--exclusion}
\begin{proposition}\label{q1:qcc:coefficient-inclusion}
For \(r\ge0\),
\begin{equation}\label{q1:qcc:coefficient-inclusion-formula}
 c_{n,k}(r)=\sum_{\varepsilon\in\{0,1\}^k}(-1)^{\sum\varepsilon_j}
 \sum_{\substack{v_1,\ldots,v_k\ge0\\
       \sum jv_j=r-\sum(m+j)\varepsilon_j}}1.
\end{equation}
An impossible inner constraint contributes zero. All inner variables may be
bounded by \(v_j\le\lfloor r/j\rfloor\).
\end{proposition}
\begin{proof}
Expand each numerator factor \(1-q^{m+j}\) by its two terms and each reciprocal
factor \((1-q^j)^{-1}\) as a geometric series. The sign records which numerator
terms were selected, and the displayed constraint is precisely the total
exponent. Coefficient extraction proves the formula.
\end{proof}
The two finite formulae organize the same coefficient differently. The first
collects repeated divisibility contributions into Bell blocks; the second uses
restricted partitions and inclusion--exclusion. Neither needs earlier values of
\(c_{n,k}(r)\).

\section{Finite Fourier extraction}
Let \(L>D\), and let \(\omega=e^{2\pi i/L}\). Orthogonality of roots of unity
implies, for \(0\le r\le D\),
\begin{equation}\label{q1:qcc:fourier-coefficient}
 c_{n,k}(r)=\frac1L\sum_{j=0}^{L-1}\omega^{-rj}G_{n,k}(\omega^j).
\end{equation}
Indeed, expand the polynomial and use
\(L^{-1}\sum_{j=0}^{L-1}\omega^{j(s-r)}=\delta_{sr}\) in this degree range.
Without the condition \(L>D\), the same expression gives the sum of coefficients
with exponent congruent to \(r\pmod L\). Later, \(q\)-Lucas supplies a direct
root-of-unity evaluation of every summand.

This identity is exact in a cyclotomic field. Floating-point use can suffer from
cancellation, so it is not automatically a preferable numerical algorithm. For
exact coefficient production, polynomial Pascal recursion or the divisor-sum
formula is usually simpler.
% ---- Q1 lines 1400--1545 ----
\chapter{Parameter derivatives and the expansion at one}
\label{q1:qcc:parameter-jets}
\section{Why the logarithmic parameter is the natural coordinate}
Put \(m=n-k\), \(D=km\), and write
\[
 G_{n,k}(q)=\GB nk=\sum_{j=0}^{D}c_{n,k}(j)q^j.
\]
The substitution \(q=e^t\) converts powers of \(q\) into exponentials. Thus
ordinary derivatives in \(t\) give power moments of the coefficient sequence,
whereas derivatives in \(q\) give falling-factorial moments. This is precisely
the distinction for which the canonical manuscript uses the Stirling basis changes.
Let \(\logq=q\,\mathrm d/\mathrm dq\), so that
\(\frac{\mathrm d}{\mathrm dt}F(e^t)=(\logq F)(e^t)\).

\begin{lemma}[The regular logarithm of an exponential difference]
\label{q1:qcc:regular-log}
As a formal series at \(u=0\),
\begin{equation}\label{q1:qcc:regular-log-eq}
 \log\frac{e^u-1}{u}
 =\frac u2+\sum_{r\ge2}\frac{\beta_r}{r}\frac{u^r}{r!}.
\end{equation}
All odd coefficients after the linear one vanish.
\end{lemma}
\begin{proof}
The quotient inside the logarithm has constant term one. Its logarithmic
 derivative is
\[
 \frac{e^u}{e^u-1}-\frac1u
 =1+\frac1u\left(\frac{u}{e^u-1}-1\right)
 =\frac12+\sum_{r\ge2}\beta_r\frac{u^{r-1}}{r!}.
\]
Integrate formally and use the zero constant term of the logarithm. The function
\(u/(e^u-1)+u/2\) is even, as direct substitution of \(-u\) shows. Hence
\(\beta_{2j+1}=0\) for \(j\ge1\).
\end{proof}

\begin{theorem}[All-order Gaussian cumulants]
\label{q1:qcc:cumulants}
Define
\begin{equation}\label{q1:qcc:kappa}
 \kappa_1(n,k)=\frac{km}{2},\qquad
 \kappa_r(n,k)=\frac{\beta_r}{r}
       \sum_{j=1}^{k}\bigl((m+j)^r-j^r\bigr)\quad(r\ge2).
\end{equation}
Then
\begin{align}
 \log\frac{G_{n,k}(e^t)}{\binom nk}
   &=\sum_{r\ge1}\kappa_r(n,k)\frac{t^r}{r!},
       \label{q1:qcc:log-cumulants}\\
 G_{n,k}(e^t)
   &=\binom nk\sum_{r\ge0}
      \Y r(\kappa_1,\ldots,\kappa_r)\frac{t^r}{r!}.
       \label{q1:qcc:cumulant-expansion}
\end{align}
These are both formal identities and convergent local Taylor expansions.
In particular,
\begin{equation}\label{q1:qcc:low-cumulants}
 \kappa_2=\frac{D(n+1)}{12},\qquad
 \kappa_3=0,\qquad
 \kappa_4=-\frac{D(n+1)\bigl(n(n+1)-D\bigr)}{120}.
\end{equation}
\end{theorem}
\begin{proof}
The cancelled product representation is
\[
 \frac{G_{n,k}(e^t)}{\binom nk}
 =\prod_{j=1}^{k}
 \frac{(e^{(m+j)t}-1)/((m+j)t)}{(e^{jt}-1)/(jt)}.
\]
Every factor now has constant term one. Apply
\cref{q1:qcc:regular-log} to each numerator and denominator. The linear
coefficient is \(\frac12\sum_{j=1}^k m=km/2\); the higher coefficients
are exactly \eqref{q1:qcc:kappa}. Exponentiation and the Bell generating
function give \eqref{q1:qcc:cumulant-expansion}.

For the second coefficient,
\(\sum_{j=1}^k((m+j)^2-j^2)=km(n+1)\). For the fourth,
expansion of the fourth powers and summation of powers of \(j\) gives
\[
 \sum_{j=1}^k((m+j)^4-j^4)
 =km(n+1)\bigl(n(n+1)-km\bigr).
\]
Together with \(\beta_2=1/6\), \(\beta_3=0\), and \(\beta_4=-1/30\),
these are \eqref{q1:qcc:low-cumulants}. The local analytic assertions follow
because the cancelled product is analytic and nonzero near \(t=0\).
\end{proof}

\subsection{Finite Bernoulli-polynomial form}
No power-sum recurrence is necessary in \eqref{q1:qcc:kappa}. Define
\(\beta_r(x)=\sum_{j=0}^r\binom rj\beta_jx^{r-j}\). Multiplying its
generating function \(te^{xt}/(e^t-1)\) by \(e^t-1\) shows
\(\beta_{r+1}(x+1)-\beta_{r+1}(x)=(r+1)x^r\). Consequently
\[
 P_r(M):=\sum_{j=1}^{M}j^r
   =\frac{\beta_{r+1}(M+1)-\beta_{r+1}(1)}{r+1}.
\]
The sum in \eqref{q1:qcc:kappa} is \(P_r(n)-P_r(m)-P_r(k)\).
This also makes its symmetry in \(k,m\) manifest.

\section{Exact coefficient moments and probability}
\begin{corollary}[Power and factorial moments]\label{q1:qcc:moment-formulas}
For \(r\ge0\),
\begin{align}
 \sum_{j=0}^{D}j^r c_{n,k}(j)
 &=\binom nk\Y r(\kappa_1,\ldots,\kappa_r),
       \label{q1:qcc:power-moments}\\
 G_{n,k}^{(r)}(1)
 &=\binom nk\sum_{a=0}^r s(r,a)
           \Y a(\kappa_1,\ldots,\kappa_a).
       \label{q1:qcc:ordinary-jets-one}
\end{align}
For \(r=0\), the first sum is the coefficient sum, with \(0^0=1\).
\end{corollary}
\begin{proof}
Differentiate the finite sum \(G_{n,k}(e^t)=\sum_jc_{n,k}(j)e^{jt}\).
This proves \eqref{q1:qcc:power-moments}. The identity
\(j^{\underline r}=\sum_a s(r,a)j^a\), which defines the signed
Stirling coefficients, gives \eqref{q1:qcc:ordinary-jets-one}.
\end{proof}

Choosing a partition uniformly from the partitions in the \(k\)-by-\(m\)
rectangle gives a random area \(X\) with
\(\Pr(X=j)=c_{n,k}(j)/\binom nk\). Its moment generating function is the
left side of \eqref{q1:qcc:cumulant-expansion}, divided by \(\binom nk\).
Thus the \(\kappa_r\) are its cumulants, and
\begin{equation}\label{q1:qcc:area-statistics}
 \E X=\frac D2,\qquad \Var(X)=\frac{D(n+1)}{12}.
\end{equation}
Complementing the rectangle sends \(X\) to \(D-X\); hence every odd
central moment is zero. Every even central moment is given by
\(\Y{2r}(0,\kappa_2,0,\kappa_4,\ldots,\kappa_{2r})\).
These assertions concern a finite distribution and require no limiting
probabilistic theorem.

\begin{example}
For \(G_{4,2}(q)=1+q+2q^2+q^3+q^4\), the coefficient sum is six,
\(\kappa_1=2\), \(\kappa_2=5/3\), and \(\kappa_4=-8/3\).
Therefore
\[
 G_{4,2}'(1)=6\cdot2=12,\qquad
 (\logq^2G_{4,2})(1)=6(2^2+5/3)=34,
\]
whereas
\(G_{4,2}''(1)=34-12=22\). Confusing the two kinds of differentiation
would already give an incorrect second derivative.
\end{example}
% ---- hand chunk 20-moments-lead.tex ----

Write \(\Delta_r=\sum_{j=1}^{k}\bigl((n-k+j)^r-j^r\bigr)\), so that
\(\kappa_r=\beta_r\Delta_r/r\) for \(r\ge2\) and \(\Delta_2=D(n+1)\). The
fourth central moment is then explicit:

% ---- Q3 lines 1760--1767 ----
\begin{equation}
 \mathbb E X=\frac D2,\qquad
 \operatorname{Var}(X)=\frac{D(n+1)}{12},\qquad
 \mathbb E(X-d/2)^4
     =3\left(\frac{D(n+1)}{12}\right)^2-\frac{\Delta_4}{120}.
 \label{q3:eq:area-moments}
\end{equation}
Here $\Delta_2=D(n+1)$, $B_2=1/6$, and $B_4=-1/30$.
% ---- hand chunk 20b-moments-glue.tex ----

Higher central moments are obtained by setting the first Bell argument to
zero, as noted above.

% ---- Q3 lines 1773--1785 ----
The first two ordinary derivatives illustrate why Euler derivatives
and ordinary derivatives must not be confused:
\begin{align}
 \left.\partial_q\G nk\right|_1&=\binom nk\frac D2,
 \label{q3:eq:gaussian-first-q-one}\\
 \left.\partial_q^2\G nk\right|_1
   &=\binom nk\left(\frac{D^2}{4}
             +\frac{D(n+1)}{12}-\frac D2\right).
 \label{q3:eq:gaussian-second-q-one}
\end{align}
The subtraction of $d/2$ converts the second raw moment into the second
falling-factorial moment. These exact identities require no limiting
probability theorem.
% ---- Q1 lines 1547--1565 ----
\section{Arbitrary parameter points and Eulerian rational functions}
For \(s\ge0\), define the rational function
\[
 \Li_{-s}(u)=\left(u\frac{\mathrm d}{\mathrm du}\right)^s\frac{u}{1-u}.
\]
For \(|u|<1\) this is \(\sum_{j\ge1}j^su^j\), but the rational
expression is the definition needed outside that disk. Set \(A_0(u)=1\).
Then
\begin{equation}\label{q1:qcc:li-eulerian}
 \Li_{-s}(u)=\frac{uA_s(u)}{(1-u)^{s+1}},\qquad
 A_{s+1}(u)=(1+su)A_s(u)+u(1-u)A_s'(u).
\end{equation}
The formula follows by applying \(u\,\mathrm d/\mathrm du\) to the
fraction. For \(s\ge1\), the coefficients of \(A_s\) are the ordinary
Eulerian numbers indexed by the number of descents: insertion of the largest
letter gives the same coefficient recurrence and initial value. Alternatively,
\eqref{q1:qcc:li-eulerian} is an independent algebraic definition of this
Eulerian normalization.

% ---- hand chunk 21-stirling-form.tex ----

For later use, write
\begin{equation}\label{q3:eq:Rr}
 R_r(x)=\left(x\frac{\dd}{\dd x}\right)^{r-1}\frac{x}{1-x}
       =\Li_{1-r}(x),\qquad r\ge1,
\end{equation}
so that \(R_{s+1}=\Li_{-s}\). The polylogarithm notation is optional: the
finite expression
\begin{equation}\label{q3:eq:Rr-stirling}
 R_r(x)=\sum_{\ell=1}^r(\ell-1)!\,S(r,\ell)\frac{x^\ell}{(1-x)^\ell}
\end{equation}
makes the rational nature explicit. To verify it, note that for \(|x|<1\)
the left side is \(\sum_{m\ge1}m^{r-1}x^m\), and
\(m^{r-1}=\sum_\ell S(r-1,\ell)\,m^{\underline\ell}\) together with
\(\sum_mm^{\underline\ell}x^m=\ell!\,x^\ell/(1-x)^{\ell+1}\) gives a sum of
the same shape; alternatively, differentiate the right side and use the
Stirling recurrence \(S(r+1,\ell)=\ell S(r,\ell)+S(r,\ell-1)\) to prove it
from \(R_1=x/(1-x)\). Rational continuation removes the restriction
\(|x|<1\). Thus both the Stirling and the Eulerian bases organize the same
rational functions; the Eulerian form is \eqref{q1:qcc:li-eulerian} with
\(s=r-1\).

% ---- Q1 lines 1566--1605 ----
\begin{theorem}[Logarithmic parameter derivatives]\label{q1:qcc:generic-jets}
At a nonzero point \(q\) for which none of \(q^1,\ldots,q^n\) equals
one, define
\begin{equation}\label{q1:qcc:L-r}
 L_r(q)=\sum_{j=1}^{k}
 \left(j^r\Li_{1-r}(q^j)-(m+j)^r\Li_{1-r}(q^{m+j})\right).
\end{equation}
Then, for \(r\ge0\),
\begin{align}
 \logq^rG_{n,k}(q)&=G_{n,k}(q)\Y r(L_1(q),\ldots,L_r(q)),
          \label{q1:qcc:theta-jets}\\
 G_{n,k}^{(r)}(q)&=q^{-r}G_{n,k}(q)
       \sum_{a=0}^r s(r,a)\Y a(L_1(q),\ldots,L_a(q)).
          \label{q1:qcc:q-jets}
\end{align}
All \(L_r\) are explicit rational functions, by
\eqref{q1:qcc:li-eulerian}.
\end{theorem}
\begin{proof}
For \(a\ge1\), repeated logarithmic differentiation gives
\[
 \logq^r\log(1-q^a)=-a^r\Li_{1-r}(q^a),\qquad r\ge1.
\]
Subtracting denominator contributions from numerator contributions gives
\(\logq^r\log G=L_r\). In the local coordinate \(q e^t\), the
logarithm of \(G(qe^t)/G(q)\) is
\(\sum_{r\ge1}L_r(q)t^r/r!\). Exponentiate and use Bell polynomials.
Finally, the operator identity
\begin{equation}\label{q1:qcc:stirling-operators}
 q^r\frac{\mathrm d^r}{\mathrm dq^r}
 =\logq(\logq-1)\cdots(\logq-r+1)
 =\sum_{a=0}^r s(r,a)\logq^a
\end{equation}
follows by applying both sides to each monomial. It proves the second formula.
\end{proof}

The displayed domain avoids unnecessary cancellation inside individual rational
terms. At roots of unity, the Gaussian polynomial is still regular, but some
\(L_r\) may have poles. The correct replacement is a deflated local
factorization, not termwise substitution into \eqref{q1:qcc:L-r}.
% ---- Q1 lines 1607--1822 ----
\chapter{Cyclotomic structure and all-order root-of-unity jets}
\label{q1:qcc:roots}
\section{Cyclotomic factorization and carries}
Let \(\Phi_d(q)\) denote the monic polynomial whose roots are the primitive
\(d\)-th roots of unity. The partition of roots by their exact order gives
\(q^a-1=\prod_{d\mid a}\Phi_d(q)\).

\begin{theorem}[Squarefree cyclotomic factorization]\label{q1:qcc:cyclotomic}
For \(0\le k\le n\), \(m=n-k\),
\begin{equation}\label{q1:qcc:cyclotomic-product}
 G_{n,k}(q)=\prod_{d=2}^{n}\Phi_d(q)^{\epsilon_d},\qquad
 \epsilon_d=\left\lfloor\frac nd\right\rfloor
      -\left\lfloor\frac kd\right\rfloor
      -\left\lfloor\frac md\right\rfloor\in\{0,1\}.
\end{equation}
Consequently every nonconstant Gaussian polynomial has distinct roots, all on
the unit circle. Its degree gives the identity
\begin{equation}\label{q1:qcc:totient-degree}
 \sum_{d=2}^{n}\varphi(d)\epsilon_d=k(n-k).
\end{equation}
\end{theorem}
\begin{proof}
Factor each \(1-q^a\) in the factorial quotient. There are
\(\lfloor n/d\rfloor\) multiples of \(d\) in its numerator factorial,
and the denominator subtracts the corresponding counts for \(k\) and \(m\).
The exponent of \(\Phi_1=q-1\) is \(n-k-m=0\); the signs also cancel.
Write \(k=ud+v\), \(m=wd+x\), with \(0\le v,x<d\). Then
\(\epsilon_d=\lfloor(v+x)/d\rfloor\), which is either zero or one.
The roots of different cyclotomic polynomials are disjoint, and each is simple.
Finally \(\deg\Phi_d=\varphi(d)\), and degrees add in the product.
\end{proof}
This carry interpretation is particularly useful: vanishing at a primitive
\(d\)-th root occurs precisely when the residues of \(k\) and \(n-k\)
produce a carry upon addition modulo \(d\).

\section{The finite \qbname-Lucas theorem}
The following classical root-of-unity formula is often called \(q\)-Lucas;
see, for example, \cite{formichella}. Its proof is short enough to
include, and the phase bookkeeping is important.

\begin{theorem}[\(q\)-Lucas]\label{q1:qcc:lucas}
Let \(\zeta\) be a primitive \(d\)-th root of unity, with \(d\ge1\).
Write \(n=ad+b\), \(k=rd+s\), where \(0\le b,s<d\). Then
\begin{equation}\label{q1:qcc:lucas-eq}
 G_{n,k}(\zeta)=\binom ar\,G_{b,s}(\zeta).
\end{equation}
The right side is zero when \(s>b\).
\end{theorem}
\begin{proof}
The finite product theorem and a division into full blocks give
\[
 \prod_{j=0}^{n-1}(1+\zeta^jz)
 =\bigl(1-(-z)^d\bigr)^a
       \prod_{j=0}^{b-1}(1+\zeta^jz).
\]
The second factor has degree smaller than \(d\). Hence its only possible
contribution to \(z^{rd+s}\) has degree \(s\), and coefficient comparison
therefore gives
\[
 \zeta^{\binom k2}G_{n,k}(\zeta)
 =(-1)^{r(d+1)}\binom ar
           \zeta^{\binom s2}G_{b,s}(\zeta).
\]
The phase identity
\(\zeta^{\binom{rd+s}{2}-\binom s2}=(-1)^{r(d+1)}\)
finishes the proof. For odd \(d\), the exponent is a multiple of \(d\).
For even \(d\), it is congruent to \(rd/2\) modulo \(d\), giving
\((-1)^r\). These are exactly the two values on the right.
The same argument covers \(d=1\).
\end{proof}

For example,
\[
 G_{2a,2r}(-1)=\binom ar,\qquad G_{2a,2r+1}(-1)=0,
\]
while \(G_{2a+1,2r}(-1)=G_{2a+1,2r+1}(-1)=\binom ar\).
If \(d\mid n\), \eqref{q1:qcc:lucas-eq} is also a fixed-subset count: a
permutation consisting of \(n/d\) disjoint \(d\)-cycles fixes a
\(k\)-element subset exactly when the subset is a union of whole cycles.
There are \(\binom{n/d}{k/d}\) such subsets if \(d\mid k\), and none
otherwise. This provides a concrete combinatorial interpretation for these
root evaluations.

\section{Deflating all apparent singularities}
A value formula such as \(q\)-Lucas does not determine derivatives at the
root. The next theorem supplies every derivative, including the case where the
Gaussian polynomial itself vanishes. The calculation uses the logarithmic
coordinate \(q=\zeta e^t\) and removes powers of \(t\) before taking a
logarithm.

For an integer \(a\ge1\), define
\begin{equation}\label{q1:qcc:root-constants}
 c_a(\zeta)=
 \begin{cases}
 -a,&d\mid a,\\
 1-\zeta^a,&d\nmid a.
 \end{cases}
\end{equation}
For \(r\ge1\), define
\begin{equation}\label{q1:qcc:eta}
 \eta_r(a;\zeta)=
 \begin{cases}
 a/2,&d\mid a,\ r=1,\\[1mm]
 \beta_r a^r/r,&d\mid a,\ r\ge2,\\[1mm]
 -a^r\Li_{1-r}(\zeta^a),&d\nmid a.
 \end{cases}
\end{equation}
All rational functions in the last line are evaluated away from their pole at
one. Put
\begin{align}
 \nu&=\left\lfloor\frac nd\right\rfloor
       -\left\lfloor\frac kd\right\rfloor
       -\left\lfloor\frac md\right\rfloor,\label{q1:qcc:root-order}\\
 C_\zeta&=\prod_{j=1}^{k}\frac{c_{m+j}(\zeta)}{c_j(\zeta)},\label{q1:qcc:root-C}\\
 \lambda_r&=\sum_{j=1}^{k}
       \bigl(\eta_r(m+j;\zeta)-\eta_r(j;\zeta)\bigr).
       \label{q1:qcc:root-lambda}
\end{align}
Here \(\nu=0\) when \(d=1\), and always \(C_\zeta\ne0\).

\begin{theorem}[Uniform all-order local formula]\label{q1:qcc:root-jets}
With the definitions above,
\begin{equation}\label{q1:qcc:root-factorization}
 \boxed{\quad
 G_{n,k}(\zeta e^t)
  =C_\zeta t^\nu
     \exp\left(\sum_{r\ge1}\lambda_r\frac{t^r}{r!}\right).
 \quad}
\end{equation}
In particular, for \(r\ge0\),
\begin{equation}\label{q1:qcc:root-t-coefficients}
 [t^{\nu+r}]G_{n,k}(\zeta e^t)
       =\frac{C_\zeta}{r!}\Y r(\lambda_1,\ldots,\lambda_r).
\end{equation}
Every ordinary Taylor derivative at \(\zeta\) is the finite sum
\begin{equation}\label{q1:qcc:root-ordinary-derivatives}
 G_{n,k}^{(R)}(\zeta)
 =C_\zeta\zeta^{-R}
   \sum_{j=\nu}^{R}s(R,j)\frac{j!}{(j-\nu)!}
               \Y{j-\nu}(\lambda_1,\ldots,\lambda_{j-\nu}).
\end{equation}
An empty sum is zero. Formula \eqref{q1:qcc:root-factorization} holds formally
in \(\Q(\zeta)[[t]]\) and analytically near \(t=0\).
\end{theorem}
\begin{proof}
For each factor in the product quotient, consider two cases. If \(d\mid a\),
then
\[
 1-\zeta^ae^{at}=-at\,\frac{e^{at}-1}{at}
 =c_a(\zeta)t
     \exp\left(\frac{at}{2}
          +\sum_{r\ge2}\frac{\beta_ra^r}{r}\frac{t^r}{r!}\right)
\]
by \cref{q1:qcc:regular-log}. If \(d\nmid a\), the factor has nonzero
constant term \(c_a(\zeta)\), and repeated differentiation yields
\[
 \log\frac{1-\zeta^ae^{at}}{1-\zeta^a}
 =-\sum_{r\ge1}a^r\Li_{1-r}(\zeta^a)\frac{t^r}{r!}.
\]
Multiplying numerator factors and dividing denominator factors produces the
constant \(C_\zeta\), the net power \(t^\nu\), and the sum of logarithms
\(\sum\lambda_rt^r/r!\). This proves \eqref{q1:qcc:root-factorization};
Bell expansion proves \eqref{q1:qcc:root-t-coefficients}.

The operator \(\logq\) becomes \(\mathrm d/\mathrm dt\). Therefore
\[
 (\logq^jG)(\zeta)=
 \begin{cases}
 0,&j<\nu,\\
 C_\zeta\dfrac{j!}{(j-\nu)!}
        \Y{j-\nu}(\lambda_1,\ldots,\lambda_{j-\nu}),&j\ge\nu.
 \end{cases}
\]
Apply \eqref{q1:qcc:stirling-operators} at \(q=\zeta\) to get
\eqref{q1:qcc:root-ordinary-derivatives}. Finally, after deflation all factors
are analytic units near zero, so their normalized logarithms converge locally.
No logarithm of a vanishing quantity was used.
\end{proof}

\begin{corollary}[First two derivatives]\label{q1:qcc:root-low-jets}
At a nonzero value, \(\nu=0\),
\[
 G(\zeta)=C_\zeta,\qquad
 G'(\zeta)=C_\zeta\zeta^{-1}\lambda_1,\qquad
 G''(\zeta)=C_\zeta\zeta^{-2}
             (\lambda_1^2+\lambda_2-\lambda_1).
\]
At a zero, necessarily simple, \(\nu=1\),
\[
 G'(\zeta)=C_\zeta\zeta^{-1},\qquad
 G''(\zeta)=C_\zeta\zeta^{-2}(2\lambda_1-1).
\]
\end{corollary}
\begin{proof}
Substitute \(R=0,1,2\) into
\eqref{q1:qcc:root-ordinary-derivatives}, using
\(s(2,1)=-1\), \(s(2,2)=1\), and
\(\Y2(x_1,x_2)=x_1^2+x_2\).
\end{proof}

\begin{example}[A zero at \(-1\)]
Take \(G_{4,1}(q)=1+q+q^2+q^3\). At \(\zeta=-1\),
\[
 \nu=1,\qquad C_{-1}=-2,\qquad
 \lambda_1=\frac32,\qquad\lambda_2=\frac{13}{12}.
\]
Thus \(G'(-1)=2\) and \(G''(-1)=-4\), in agreement with
\[
 G_{4,1}(-1+h)=2h-2h^2+h^3.
\]
\end{example}

At \(\zeta=1\), \(C_1=\binom nk\), \(\nu=0\), and
\(\lambda_r=\kappa_r(n,k)\): the expansion at one is recovered. Together,
\cref{q1:qcc:gaussian-coefficients,q1:qcc:generic-jets,q1:qcc:root-jets} give explicit
coefficients at every complex base.
% ---- hand chunk 22-local-head.tex ----

\section{A second local prescription through the linear factors}
\Cref{q1:qcc:root-jets} works in the logarithmic coordinate \(q=\zeta e^t\)
and deflates each factor \(1-\zeta^ae^{at}\). An equivalent prescription works
in the additive coordinate \(q=\zeta+h\) and uses the linear factors of the
cyclotomic factorization \eqref{q1:qcc:cyclotomic-product} directly. Write
\(\nu_d(n,k)\) for the exponent \(\epsilon_d\) in that factorization, and
\(\Phi_\ell\) for the \(\ell\)th cyclotomic polynomial. The Stirling-form
inputs of \eqref{q3:eq:Rr-stirling} express the logarithmic-coordinate
inputs \(\eta_r\) of \eqref{q1:qcc:eta} as finite rational sums,
\(-a^r\Li_{1-r}(\zeta^a)=-a^r\sum_{v=1}^{r}(v-1)!\,S(r,v)\,
\bigl(\zeta^a/(1-\zeta^a)\bigr)^v\), which is the form one of the sources
used; the two prescriptions agree because both compute the same Taylor
coefficients.


% ---- Q3 lines 1860--1897 ----
Fix $\zeta$ primitive of order $\ell\ge2$, let
$\nu=\nu_\ell(n,k)$, and put
\[
 H(q)=\frac{\G nk}{\Phi_\ell(q)^\nu}.
\]
When $\ell>n$, take $\nu=0$. Then $H(\zeta)\ne0$. The following is a
finite algebraic prescription for the full Taylor expansion at $\zeta$:
\begin{align}
 L_r&=(-1)^{r-1}(r-1)!
      \sum_{\substack{2\le d\le n\\d\ne\ell}}\nu_d(n,k)
      \sum_{\substack{\omega\text{ primitive}\\d\text{th root}}}
                         (\zeta-\omega)^{-r},
 \label{q3:eq:root-local-cumulants}\\
 \Gb nk{\zeta+h}
   &=\Phi_\ell(\zeta+h)^\nu H(\zeta)
        \sum_{m\ge0}\Yc_m(L_1,\ldots,L_m)\frac{h^m}{m!}.
 \label{q3:eq:root-local-expansion}
\end{align}
To prove this, factor $H$ into the linear factors $q-\omega$ of
\cref{q1:qcc:cyclotomic-product}. Its $r$th logarithmic derivative at
$\zeta$ is exactly \cref{q3:eq:root-local-cumulants}. Exponentiate its
Taylor series using Bell polynomials and multiply back the removed
factor. The series is analytic up to the nearest zero of $H$; it is
also an identity of formal Taylor series over the cyclotomic field.
Each prescribed coefficient is a finite sum, even though the full
Taylor series for $H$ written as an exponential is infinite.

If $\nu=1$, the first nonzero coefficient is particularly simple:
\begin{equation}
 \left.\partial_q\G nk\right|_{q=\zeta}
       =\Phi_\ell'(\zeta)H(\zeta).
 \label{q3:eq:root-simple-derivative}
\end{equation}
Higher coefficients are finite convolutions of the polynomial
$\Phi_\ell(\zeta+h)$ with the Bell coefficients in
\cref{q3:eq:root-local-expansion}. Thus root evaluation, multiplicity,
and all local derivatives fit the same coefficient calculus, provided
that a zero is removed before taking its logarithm.
% ---- Q1 lines 1824--1965 ----
\chapter{Weighted sums and differentiated identities}
\label{q1:qcc:weighted-sums}
\section{Polynomial weights in the summation index}
The finite product theorem can be differentiated in its auxiliary variable
without differentiating the parameter \(q\). This gives a second, independent
source of weighted Gaussian sums. Let
\[
 P_N(z)=(-z;q)_N
   =\sum_{k=0}^{N}q^{\binom k2}\GB Nk z^k,
 \qquad \Theta_z=z\frac{\mathrm d}{\mathrm dz}.
\]

\begin{theorem}[All power-weighted finite sums]\label{q1:qcc:weighted-Gaussian}
At a point where \(P_N(z)\ne0\), define
\[
 K_r(z)=-\sum_{j=0}^{N-1}\Li_{1-r}(-zq^j),\qquad r\ge1.
\]
Then, for \(r\ge0\),
\begin{equation}\label{q1:qcc:weighted-Gaussian-eq}
 \sum_{k=0}^{N}k^r q^{\binom k2}\GB Nk z^k
   =P_N(z)\Y r(K_1(z),\ldots,K_r(z)).
\end{equation}
The right side, after cancellation, is a polynomial and therefore extends to
the zeros of \(P_N\).
\end{theorem}
\begin{proof}
The left side is \(\Theta_z^rP_N(z)\). For every \(r\ge1\),
\(\Theta_z^r\log P_N(z)=K_r(z)\). Taylor expansion of
\(P_N(ze^t)/P_N(z)\) and the Bell generating function give the result.
The polynomial extension follows because the two rational expressions agree
away from finitely many zeros.
\end{proof}
For example,
\[
 \sum_{k=0}^{N}kq^{\binom k2}\GB Nk z^k
  =(-z;q)_N\sum_{j=0}^{N-1}\frac{zq^j}{1+zq^j}.
\]
For a general polynomial \(p(k)=\sum_{r=0}^{R}a_rk^r\), sum
\eqref{q1:qcc:weighted-Gaussian-eq} with weights \(a_r\). Alternatively,
expand \(p\) in falling factorials and use ordinary derivatives in \(z\).

\section{Alternating sums at a product zero}
At \(z=1\), the product \((z;q)_N\) vanishes for \(N\ge1\), so a
logarithmic formula should first be deflated. Set
\[
 H_N(z)=(zq;q)_{N-1},\qquad (z;q)_N=(1-z)H_N(z).
\]
Assume for the moment that \(q^j\ne1\) for \(1\le j<N\), and put
\begin{equation}\label{q1:qcc:alternating-kernels}
 J_s=-(s-1)!\sum_{j=1}^{N-1}
             \frac{q^{js}}{(1-q^j)^s}\quad(s\ge1).
\end{equation}

\begin{proposition}[Alternating factorial moments]\label{q1:qcc:alternating}
For \(r\ge1\),
\begin{equation}\label{q1:qcc:alternating-eq}
 \sum_{k=0}^{N}(-1)^kq^{\binom k2}\GB Nk\fall{k}{r}
   =-r(q;q)_{N-1}\Y{r-1}(J_1,\ldots,J_{r-1}).
\end{equation}
For \(r=0\), the sum is zero. Ordinary powers are recovered by
\(k^r=\sum_{a=0}^{r}S(r,a)\fall{k}{a}\).
\end{proposition}
\begin{proof}
The left side is the \(r\)-th derivative of \((z;q)_N\) at one.
Leibniz's rule and the factor \(1-z\) leave only
\(-rH_N^{(r-1)}(1)\). The ordinary derivatives of
\(\log H_N(z)\) at one are exactly \(J_s\), and
\(H_N(1)=(q;q)_{N-1}\). Apply the ordinary Bell derivative formula,
which follows from the Taylor expansion of \(H_N(1+t)/H_N(1)\).
The zeroth-order sum is \((1;q)_N=0\). The conversion of powers to
falling factorials is the defining ordinary second-kind Stirling expansion.
\end{proof}

When \(r>N\), the left side is zero, and the right side is also zero:
it represents a derivative of the degree-\(N-1\) polynomial \(H_N\)
of order \(r-1>N-1\). At excluded roots of unity, polynomial
specialization or the root deflation of \cref{q1:qcc:root-jets} replaces
termwise evaluation of \eqref{q1:qcc:alternating-kernels}.

\begin{auditbox}[title={Which cancellation is actually present?}]
The Gaussian alternating kernel does not simply annihilate every polynomial
weight of degree less than \(N\). Already at \(N=2\),
\[
 \sum_{k=0}^{2}(-1)^kq^{\binom k2}\GB2k k=q-1,
\]
which is nonzero for generic \(q\). What the product theorem gives is
cancellation at specified \emph{geometric spectral points}. The Jackson stencil
and the extrapolation filters below include the necessary additional powers of
\(q\). Replacing every ordinary binomial coefficient in an additive
finite-difference identity by a Gaussian coefficient is not a valid rule.
\end{auditbox}

\section{An all-order differentiation of central Vandermonde}
The central convolution proved earlier is
\[
 G_{2n,n}(q)=\sum_{k=0}^{n}q^{k^2}G_{n,k}(q)^2.
\]
Differentiating this identity at \(q=1\) produces infinitely many finite
identities involving ordinary binomial coefficients, Bernoulli numbers, and
Bell polynomials. The logarithmic coordinate makes the formula uniform.

\begin{theorem}[Differentiated central convolution]
\label{q1:qcc:differentiated-central}
For \(r\ge0\), define, for each \(k\),
\[
 v_1^{(k)}=nk,\qquad
 v_a^{(k)}=2\kappa_a(n,k)\quad(a\ge2).
\]
Then
\begin{equation}\label{q1:qcc:differentiated-central-eq}
 \binom{2n}{n}\Y r\bigl(\kappa_1(2n,n),\ldots,\kappa_r(2n,n)\bigr)
 =\sum_{k=0}^{n}\binom nk^2
           \Y r\bigl(v_1^{(k)},\ldots,v_r^{(k)}\bigr).
\end{equation}
\end{theorem}
\begin{proof}
Substitute \(q=e^t\). On the \(k\)-th summand the logarithm, after
removing the constant \(\binom nk^2\), is
\[
 k^2t+2\sum_{a\ge1}\kappa_a(n,k)\frac{t^a}{a!}.
\]
Its linear coefficient is
\(k^2+2\kappa_1(n,k)=k^2+k(n-k)=nk\), and its remaining cumulants
are \(2\kappa_a(n,k)\). Expand both sides by
\cref{q1:qcc:cumulants} and compare coefficients of \(t^r/r!\).
\end{proof}

For \(r=1\) and \(n\ge1\), this gives
\[
 \sum_{k=0}^{n} k\binom nk^2=\frac n2\binom{2n}{n}.
\]
For \(r=2\), it gives the less immediate identity
\[
 \sum_{k=0}^{n}\binom nk^2
    \left(n^2k^2+\frac{k(n-k)(n+1)}6\right)
 =\binom{2n}{n}\left(\frac{n^4}{4}
                      +\frac{n^2(2n+1)}{12}\right).
\]
Higher identities require no new differentiation argument: the same finite Bell
formula gives their complete coefficients. The same procedure works for every
block convolution in \cref{q1:qcc:convolution}; each explicit power of \(q\)
adds only to the first cumulant.
% ---- Q1 lines 1967--2210 ----
\chapter{Gaussian filters on geometric grids}
\label{q1:qcc:filters}
\section{Spectral design instead of recursive extrapolation}
Let \(0<\rho<1\) and define the dilation operator
\((T_\rho f)(h)=f(\rho h)\). Its eigenfunctions are the monomials:
\(T_\rho h^m=\rho^m h^m\). To preserve a constant and annihilate the
first \(N\) positive powers, construct the normalized spectral polynomial
\begin{equation}\label{q1:qcc:filter-def}
 R_N=\prod_{a=1}^{N}\frac{T_\rho-\rho^a}{1-\rho^a}.
\end{equation}
Here \(R_0\) is the identity operator. This is a geometric version of the
source's principle of choosing a basis in which an operator is diagonal.

\begin{theorem}[Closed Gaussian extrapolation weights]
\label{q1:qcc:filter-weights}
For every \(N\ge0\),
\begin{align}
 (R_Nf)(h)&=\sum_{j=0}^{N}w_{N,j}f(\rho^jh),
          \label{q1:qcc:filter-sum}\\
 w_{N,j}&=
 \frac{(-1)^{N-j}\rho^{\binom{N-j+1}{2}}
                \GB[\rho]Nj}{(\rho;\rho)_N}.
          \label{q1:qcc:weights}
\end{align}
The filter preserves constants and annihilates \(h^m\) for
\(1\le m\le N\). For every integer \(m\ge N+1\), its complete
remaining response is
\begin{equation}\label{q1:qcc:filter-response}
 R_Nh^m=(-1)^N\rho^{\binom{N+1}{2}}
              \GB[\rho]{m-1}{N}h^m.
\end{equation}
\end{theorem}
\begin{proof}
Expand the numerator polynomial in \eqref{q1:qcc:filter-def}. The coefficient
of \(T_\rho^j\) is
\((-1)^{N-j}e_{N-j}(\rho,\ldots,\rho^N)\). Principal specialization
of the elementary symmetric function gives
\[
 e_{N-j}(\rho,\ldots,\rho^N)
 =\rho^{\binom{N-j+1}{2}}\GB[\rho]{N}{N-j},
\]
and Gaussian symmetry gives \eqref{q1:qcc:weights}. Evaluate the spectral
polynomial at \(1,\rho,\ldots,\rho^N\) to prove preservation and
annihilation. For \(m\ge N+1\), direct factorization gives
\[
 \prod_{a=1}^{N}\frac{\rho^m-\rho^a}{1-\rho^a}
 =(-1)^N\rho^{N(N+1)/2}
      \prod_{a=1}^{N}\frac{1-\rho^{m-a}}{1-\rho^a}
 =(-1)^N\rho^{\binom{N+1}{2}}\GB[\rho]{m-1}{N}.
\]
This is \eqref{q1:qcc:filter-response}.
\end{proof}

\begin{corollary}[Exact analytic error series and a bound]
\label{q1:qcc:filter-error}
Suppose \(f(h)=L+\sum_{m\ge1}c_mh^m\) is analytic for \(|h|<R\).
Then, throughout that disk,
\begin{equation}\label{q1:qcc:filter-error-series}
 (R_Nf)(h)-L=(-1)^N\rho^{\binom{N+1}{2}}
       \sum_{m\ge N+1}\GB[\rho]{m-1}{N}c_mh^m.
\end{equation}
If, additionally, \(f\) is analytic on a neighborhood of the closed disk
\(|h|\le R\) and \(|f(h)|\le M\) there, then, for \(x=|h|/R<1\),
\begin{equation}\label{q1:qcc:filter-error-bound}
 |(R_Nf)(h)-L|
 \le M\rho^{\binom{N+1}{2}}
                  \frac{x^{N+1}}{(x;\rho)_{N+1}}.
\end{equation}
\end{corollary}
\begin{proof}
The operator is a finite combination of evaluations at points inside the disk,
so it can be applied term by term to the convergent power series. Use
\eqref{q1:qcc:filter-response}. Cauchy's coefficient bound gives
\(|c_m|\le M R^{-m}\). All Gaussian values are nonnegative at
\(0<\rho<1\), and the reciprocal finite-product theorem gives
\[
 \sum_{m\ge N+1}\GB[\rho]{m-1}{N}x^m
   =\frac{x^{N+1}}{(x;\rho)_{N+1}}.
\]
Taking absolute values therefore proves \eqref{q1:qcc:filter-error-bound}.
\end{proof}

\section{Noise amplification is bounded at a fixed geometric ratio}
\begin{theorem}[Exact absolute-weight norm]\label{q1:qcc:filter-noise}
For \(0<\rho<1\),
\begin{equation}\label{q1:qcc:noise-norm}
 \sum_{j=0}^{N}|w_{N,j}|
       =\frac{(-\rho;\rho)_N}{(\rho;\rho)_N}
       \le\prod_{a\ge1}\frac{1+\rho^a}{1-\rho^a}<\infty.
\end{equation}
Thus the bound is uniform in \(N\) when \(\rho\) is fixed. One explicit,
though not sharp, estimate is
\begin{equation}\label{q1:qcc:noise-simple-bound}
 \sum_{j=0}^{N}|w_{N,j}|
       \le\exp\left(\frac{2\rho}{(1-\rho)^2}\right).
\end{equation}
If each sample has absolute error at most \(\varepsilon\), the resulting
filtered error is at most \(\varepsilon\) times
\eqref{q1:qcc:noise-norm}.
\end{theorem}
\begin{proof}
After taking absolute values, the numerator sum in \eqref{q1:qcc:weights} is
\[
 \sum_{b=0}^{N}\rho^{\binom{b+1}{2}}\GB[\rho]Nb
       =(-\rho;\rho)_N
\]
by the finite product theorem. For \(0\le u<1\),
\[
 \log\frac{1+u}{1-u}
 =2\sum_{r\ge0}\frac{u^{2r+1}}{2r+1}
 \le\frac{2u}{1-u}.
\]
Consequently the logarithm of the infinite product is at most
\(\sum_{a\ge1}2\rho^a/(1-\rho^a)
 \le2\rho/(1-\rho)^2\), which proves convergence and the bound.
The final assertion is the triangle inequality applied to
\eqref{q1:qcc:filter-sum}.
\end{proof}

The fixed-ratio hypothesis matters. The bound need not remain moderate as
\(\rho\) approaches one. Nor does it control the error in generating the
samples themselves; the theorem assumes a bound on those errors and describes
only the linear combination's amplification.

\begin{example}[A dyadic three-point formula]
For \(\rho=1/2\) and \(N=2\),
\begin{equation}\label{q1:qcc:dyadic-filter}
 (R_2f)(h)=\frac13 f(h)-2f(h/2)+\frac83f(h/4).
\end{equation}
It preserves constants, cancels the linear and quadratic terms, and has leading
error \(c_3h^3/8\). Its exact absolute-weight norm is five. More generally,
\eqref{q1:qcc:filter-error-series} specifies \emph{every} remaining power,
not just the leading one.
\end{example}

\section{Fractional powers and arithmetic exponent lattices}
For positive \(h\), let the unwanted exponents be
\(\alpha,\alpha+\delta,\ldots,\alpha+(N-1)\delta\), where
\(\alpha,\delta>0\). Define
\begin{equation}\label{q1:qcc:general-filter}
 R_{N;\alpha,\delta}
 =\prod_{a=0}^{N-1}
       \frac{T_\rho-\rho^{\alpha+a\delta}}
            {1-\rho^{\alpha+a\delta}}.
\end{equation}
It preserves constants and annihilates exactly those prescribed power modes.
The same elementary-symmetric specialization gives the closed weights
\begin{equation}\label{q1:qcc:general-weights}
 [T_\rho^j]R_{N;\alpha,\delta}
 =\frac{(-1)^{N-j}
        \rho^{\alpha(N-j)+\delta\binom{N-j}{2}}
        \GB[\rho^\delta]Nj}
       {(\rho^\alpha;\rho^\delta)_N}.
\end{equation}
This is proved exactly as \eqref{q1:qcc:weights}, now with the alphabet
\(\rho^\alpha,\rho^{\alpha+\delta},\ldots\). For a general exponent
\(\beta\), the full response is the explicit product
\[
 R_{N;\alpha,\delta}h^\beta
  =h^\beta\prod_{a=0}^{N-1}
      \frac{\rho^\beta-\rho^{\alpha+a\delta}}
           {1-\rho^{\alpha+a\delta}}.
\]
In particular, an even-power error expansion in the original step size uses
\(\alpha=\delta=2\), and hence the Gaussian base is \(\rho^2\),
not \(\rho\).

\section{Logarithmic errors and repeated spectral roots}
The geometric operator also handles polynomial--logarithmic errors. The key
identity is
\begin{equation}\label{q1:qcc:log-spectral-identity}
 p(T_\rho)\bigl(h^\beta(\log h)^r\bigr)
  =\frac{\partial^r}{\partial\beta^r}
            \bigl(h^\beta p(\rho^\beta)\bigr),\qquad h>0.
\end{equation}
Indeed, \(h^\beta(\log h)^r=\partial_\beta^rh^\beta\), and
\(p(T_\rho)\) is a finite sum of substitutions independent of \(\beta\).

\begin{proposition}[Confluent geometric cancellation]
\label{q1:qcc:confluent-filter}
If \(p(z)\) has a zero of multiplicity at least \(s\) at
\(z=\rho^\beta\), then \(p(T_\rho)\) annihilates
\(h^\beta(\log h)^r\) for every \(0\le r<s\).
For distinct positive exponents \(\beta_1,\ldots,\beta_A\) and positive
integers \(s_1,\ldots,s_A\), the normalized filter
\begin{equation}\label{q1:qcc:confluent-polynomial}
 p(z)=\prod_{a=1}^{A}
        \left(\frac{z-\lambda_a}{1-\lambda_a}\right)^{s_a},
 \qquad\lambda_a=\rho^{\beta_a},
\end{equation}
preserves constants and cancels all the corresponding logarithmic modes.
Its coefficients have the recurrence-free Bell formula
\begin{align}
 C_0&=\prod_{a=1}^{A}
       \frac{(-\lambda_a)^{s_a}}{(1-\lambda_a)^{s_a}},
       \label{q1:qcc:confluent-C}\\
 [z^\ell]p(z)&=\frac{C_0}{\ell!}
   \Y\ell\left(-0!\sum_a s_a\lambda_a^{-1},
                -1!\sum_a s_a\lambda_a^{-2},\ldots,
                -(\ell-1)!\sum_a s_a\lambda_a^{-\ell}\right).
       \label{q1:qcc:confluent-Bell}
\end{align}
For \(\ell>\sum_a s_a\), the latter expression is identically zero.
\end{proposition}
\begin{proof}
The derivative of \(\rho^\beta\) with respect to \(\beta\) is
\((\log\rho)\rho^\beta\ne0\). A zero of order \(s\) in the
spectral coordinate therefore remains a zero of order \(s\) in the
exponent coordinate. Applying \eqref{q1:qcc:log-spectral-identity} with
\(r<s\) gives zero. Formula \eqref{q1:qcc:confluent-polynomial} has all
these zeros and satisfies \(p(1)=1\).

For its coefficients, write
\[
 p(z)=C_0\prod_{a=1}^{A}(1-z/\lambda_a)^{s_a},\qquad
 \log\frac{p(z)}{C_0}
   =-\sum_{r\ge1}\left(\sum_as_a\lambda_a^{-r}\right)\frac{z^r}{r}.
\]
Exponentiation is exactly the Bell master formula. The degree assertion follows
because \(p\) is a polynomial of degree \(\sum_as_a\).
\end{proof}

For an asymptotic expansion rather than a convergent series, the cancellation
claims apply to any finite collection of terms for which a remainder estimate
has been established. A finite filter transports that remainder by the triangle
inequality. One should not apply an infinite formal expansion as an equality of
functions without a separate convergence or asymptotic argument.

\section{Connection with the Fabius-function setting}
The relevant bridge to the surrounding repository is the geometric grid and its
spectral algebra. Whenever an approximation to a Fabius- or Rvachev-related
quantity is known to have an expansion in powers of a dyadic scale, its known
error terms can be canceled by the corresponding filter. For a grid ratio of
\(1/2\), integral powers use Gaussian base \(1/2\); an even-power
expansion uses base \(1/4\). Logarithmic corrections require repeated roots,
as just proved.

This application is deliberately conditional on the expansion of the particular
approximation being used. No assertion here identifies a specific Fourier
normalization, spline sequence, or exact dyadic Fabius value with such an error
series without proving that additional bridge. The unconditional result is the
filter algebra, including explicit weights, complete power response, analytic
remainder bound, and noise norm. That separation makes the result reusable
without importing hidden normalization assumptions.
% ---- hand chunk 23-uniform-head.tex ----

\section{Application: moments of geometric uniform convolutions}
\label{q3:sec:uniform}
This application returns to the Fabius--Rvachev setting of the repository
and makes the conditional remark of the previous section concrete. Bell
polynomials make the finite-convolution moments of a geometric uniform sum
polynomial in one geometric tail parameter, and the filter of
\cref{q1:qcc:filter-weights} then removes that parameter exactly. In the
notation used below, \(w_{d,j}(s)\) is the weight \(w_{N,j}\) of
\eqref{q1:qcc:weights} with \(N=d\) and \(\rho=s\), and the sample point
\(x_0\) plays the role of the step \(h\). The same moment sequence is
computed by different means, through the Thue--Morse form of the dyadic
values, in the companion volume on recurrence-free dyadic values of the
Fabius function filed beside the sampling volume of this corpus.


% ---- Q3 lines 2134--2205 ----
\section{Moments from finite cumulant formulae}
Let $0<q<1$, let $U_0,U_1,\ldots$ be independent uniform random variables
on $[-1,1]$, and set
\begin{equation}
 X_q=(1-q)\sum_{j\ge0}q^jU_j,\qquad
 X_{q,N}=(1-q)\sum_{j=0}^{N-1}q^jU_j.
 \label{q3:eq:uniform-model}
\end{equation}
The first series converges absolutely, even pointwise for all choices
of the $U_j$, since its absolute value is bounded by one. Independence
and $\mathbb E e^{tU}=\sinh(t)/t$ give
\begin{equation}
 \mathbb E e^{tX_q}
   =\prod_{j\ge0}\frac{\sinh((1-q)q^jt)}{(1-q)q^jt}.
 \label{q3:eq:uniform-mgf}
\end{equation}
The finite products converge uniformly on compact $t$-sets: each
factor is $1+O(q^{2j})$ there, and the tails of $X_q$ also converge
uniformly, justifying the expectation limit directly.

\begin{theorem}[Closed even moments]
Let $\gamma_j$ be the cumulants of $X_q$. Then
\begin{equation}
 \gamma_{2r+1}=0\quad(r\ge0),\qquad
 \gamma_{2r}=\frac{2^{2r}B_{2r}}{2r}
                \frac{(1-q)^{2r}}{1-q^{2r}}
       \quad(r\ge1).
 \label{q3:eq:uniform-cumulants}
\end{equation}
Consequently,
\begin{align}
 \mathbb E X_q^{2d}
   &=\Yc_{2d}(0,\gamma_2,0,\gamma_4,\ldots,\gamma_{2d})
 \label{q3:eq:uniform-moments-bell}\\
   &=(2d)!\sum_{\sum_{r=1}^d r m_r=d}
          \prod_{r=1}^d\frac1{m_r!}
                 \left(\frac{\gamma_{2r}}{(2r)!}\right)^{m_r}.
 \label{q3:eq:uniform-moments-profile}
\end{align}
For $X_{q,N}$ replace $\gamma_{2r}$ by
$\gamma_{2r}(1-q^{2rN})$.
\end{theorem}
\begin{proof}
From the Bernoulli generating function,
\[
 \log\frac{\sinh z}{z}
    =\sum_{r\ge1}\frac{2^{2r}B_{2r}}{2r(2r)!}z^{2r}.
\]
For example, it follows by writing
$\sinh z/z=e^{-z}(e^{2z}-1)/(2z)$ and applying the previously proved
logarithm expansion; the linear terms cancel. Sum this identity over
$z=(1-q)q^jt$ in a sufficiently small disk around zero. The geometric
sum is $(1-q^{2r})^{-1}$ for the infinite model, or
$(1-q^{2rN})/(1-q^{2r})$ for the finite model. This proves the cumulant
formula. Exponentiating the cumulant series gives the moment formula
by the Bell identity. Each coefficient uses finitely many cumulants,
so the final formulae have no remaining limiting or recurrence step.
\end{proof}

For $q=1/2$, this centered geometric-uniform sum has density
$\operatorname{up}$ on $[-1,1]$ under the Rvachev normalization of
\cite{proveit-primary}. Thus, in particular,
\begin{equation}
 \int_{-1}^1x^2\operatorname{up}(x)\dd x=\frac19,
 \qquad
 \int_{-1}^1x^4\operatorname{up}(x)\dd x=\frac{19}{675}.
 \label{q3:eq:up-moment-examples}
\end{equation}
The fourth cumulant is $-2/225$ and the second is $1/9$, so the second
number is $3(1/9)^2-2/225$. More generally every even moment at a
rational base $0<q<1$ is rational by
\cref{q3:eq:uniform-moments-profile}.
% ---- Q3 lines 2268--2294 ----
\begin{corollary}[Exact finite recovery of geometric-convolution moments]
For $d,N\ge0$, put $s=q^2$. Then
\begin{equation}
 \boxed{\displaystyle
 \mathbb E X_q^{2d}
       =\sum_{j=0}^d w_{d,j}(q^2)
                   \mathbb E X_{q,N+j}^{2d}.}
 \label{q3:eq:moment-extrapolation}
\end{equation}
Here $X_{q,0}=0$ and the zeroth moment is one, including at zero.
\end{corollary}
\begin{proof}
In \cref{q3:eq:uniform-moments-profile}, the finite-$N$ replacement is
$\gamma_{2r}\mapsto\gamma_{2r}(1-y^r)$ with $y=q^{2N}$.
A profile with $\sum rm_r=d$ therefore has degree at most $d$ in $y$.
Thus $\mathbb E X_{q,N}^{2d}=P_d(s^N)$ for a polynomial $P_d$ of
degree at most $d$, whose constant term is $\mathbb E X_q^{2d}$.
Apply \cref{q1:qcc:filter-weights} with $x_0=s^N$.
\end{proof}
For $d=1$ the filter is
$P(0)=(P(sx_0)-sP(x_0))/(1-s)$. The higher filters remove every
geometric error $s^{rN}$ with $1\le r\le d$ at once. This exactness
has been proved for the moment sequence through its polynomial tail
parameter; it does not assert exactness for arbitrary pointwise
samples of a density or for a general asymptotic expansion. That
qualification is essential in applications to spline approximations
and Fabius values.
% ---- Q1 lines 2212--2354 ----
\chapter{Multinomial coefficients and factorial ratios}
\label{q1:qcc:factorial-ratios}
\section{The Gaussian multinomial extension}
For \(k_1+\cdots+k_s=n\), define
\[
 G_{k_1,\ldots,k_s}(q)
  =\frac{(q;q)_n}{\prod_{i=1}^{s}(q;q)_{k_i}}
  =\GB n{k_1}\GB{n-k_1}{k_2}\cdots
                  \GB{n-k_1-\cdots-k_{s-1}}{k_s}.
\]
The product expression proves polynomiality, nonnegative integer coefficients,
and the value \(n!/\prod_i k_i!\) at one. The degree is
\[
 D=\sum_{i<j}k_ik_j=\frac12\left(n^2-\sum_i k_i^2\right).
\]
The last expression follows by adding the degrees in the product; reciprocity
follows by multiplying the Gaussian reciprocity identities.

\begin{proposition}[Multinomial coefficient and parameter kernels]
\label{q1:qcc:multinomial-kernels}
For \(r\ge1\), put
\[
 A_r=\sum_{i=1}^{s}\sum_{\substack{d\mid r\\d\le k_i}}d
           -\sum_{\substack{d\mid r\\d\le n}}d.
\]
Then
\begin{equation}\label{q1:qcc:multinomial-coefficients}
 [q^R]G_{k_1,\ldots,k_s}(q)
   =\frac1{R!}\Y R\bigl(0!A_1,1!A_2,\ldots,(R-1)!A_R\bigr).
\end{equation}
Its logarithmic-parameter cumulants at one are
\begin{equation}\label{q1:qcc:multinomial-cumulants}
 \kappa_1=\frac D2,\qquad
 \kappa_r=\frac{\beta_r}{r}
       \left(P_r(n)-\sum_{i=1}^{s}P_r(k_i)\right)\quad(r\ge2),
\end{equation}
where \(P_r(M)=\sum_{j=1}^{M}j^r\). Its multiplicity at a primitive
\(d\)-th root is
\begin{equation}\label{q1:qcc:multinomial-carry}
 \nu_d=\left\lfloor\frac nd\right\rfloor
            -\sum_{i=1}^{s}\left\lfloor\frac{k_i}{d}\right\rfloor
          \in\{0,1,\ldots,s-1\}.
\end{equation}
\end{proposition}
\begin{proof}
For the coefficient formula, take the formal logarithm at zero. Each
\(\log(1-q^j)\) contributes \(-j/r\) to the coefficient of \(q^r\)
when \(j\mid r\), and zero otherwise. Subtracting denominator
logarithms gives \(A_r/r\), and Bell expansion proves
\eqref{q1:qcc:multinomial-coefficients}. The calculation at one is the
factor-by-factor regular logarithm of \cref{q1:qcc:regular-log}, with
\(P_r(n)-\sum_iP_r(k_i)\) in place of the two-block power-sum difference.
For the multiplicity, count multiples of \(d\) in the factorials.
Writing \(k_i=a_id+b_i\), \(0\le b_i<d\), gives
\(\nu_d=\lfloor\sum_i b_i/d\rfloor\), proving the stated range.
\end{proof}

Unlike a Gaussian binomial polynomial, a Gaussian multinomial need not be
squarefree. For example, the polynomial \(G_{1,1,1,1}(q)=[4]_q!\) has multiplicity
two at \(-1\):
\begin{equation}\label{q1:qcc:multinomial-repeat-example}
 G_{1,1,1,1}(q)=(1+q)^2(1+q+q^2)(1+q^2).
\end{equation}
For its local jets, define \(C_\zeta\) and \(\lambda_r\) by taking
all exponents \(1,\ldots,n\) in the numerator and all exponents
\(1,\ldots,k_i\), with multiplicities, in the denominator of
\eqref{q1:qcc:root-C} and \eqref{q1:qcc:root-lambda}. Then
\eqref{q1:qcc:root-factorization} and
\eqref{q1:qcc:root-ordinary-derivatives} hold verbatim with
\(\nu=\nu_d\). Their proof was factorwise and never required
\(\nu\le1\). In this way the same Bell formula handles higher-order zeros
without a new inversion calculation.

\subsection{A multinomial Lucas formula}
Let \(n=ad+b\), \(k_i=a_id+b_i\), with all remainders in
\(\{0,\ldots,d-1\}\). If \(\sum_i b_i>b\), then \(\nu_d>0\)
and the multinomial vanishes at a primitive \(d\)-th root. If
\(\sum_i b_i=b\), there is no carry and
\begin{equation}\label{q1:qcc:multinomial-lucas}
 G_{k_1,\ldots,k_s}(\zeta)
 =\binom{a}{a_1,\ldots,a_s}\,
                 G_{b_1,\ldots,b_s}(\zeta).
\end{equation}
Indeed, use the product of Gaussian binomials and apply \(q\)-Lucas at
each step. Without a carry, the ordinary binomial factors multiply to the
ordinary multinomial, and the remainder factors multiply to the smaller
Gaussian multinomial. With a carry, the multiplicity calculation already
proves vanishing.

\section{A polynomiality criterion for balanced factorial ratios}
Let \(a_1,\ldots,a_u\) and \(b_1,\ldots,b_v\) be nonnegative integers
with \(\sum_i a_i=\sum_j b_j\), and set
\[
 F(q)=\frac{\prod_{i=1}^{u}(q;q)_{a_i}}
                 {\prod_{j=1}^{v}(q;q)_{b_j}}.
\]
Let \(M\) be the largest index, taking \(M=0\) if every index is zero.

\begin{theorem}[Finite cyclotomic criterion]\label{q1:qcc:ratio-criterion}
The rational function \(F(q)\) is a polynomial if and only if
\begin{equation}\label{q1:qcc:ratio-floors}
 E_d:=\sum_i\left\lfloor\frac{a_i}{d}\right\rfloor
       -\sum_j\left\lfloor\frac{b_j}{d}\right\rfloor\ge0
       \qquad(2\le d\le M).
\end{equation}
When these conditions hold,
\begin{equation}\label{q1:qcc:ratio-cyclotomic}
 F(q)=\prod_{d=2}^{M}\Phi_d(q)^{E_d}\in\Z[q],\qquad F(0)=1.
\end{equation}
Its degree and value at one are
\begin{equation}\label{q1:qcc:ratio-degree-value}
 D=\frac12\left(\sum_i a_i^2-\sum_j b_j^2\right),\qquad
 F(1)=\frac{\prod_i a_i!}{\prod_j b_j!}.
\end{equation}
It satisfies \(F(q^{-1})=q^{-D}F(q)\). All its parameter jets follow
from the same factorwise deflation as \cref{q1:qcc:root-jets}.
\end{theorem}
\begin{proof}
Cyclotomic factorization gives \(F=\prod_{d=1}^{M}\Phi_d^{E_d}\).
The balancing hypothesis makes \(E_1=0\) and cancels the signs from
\(1-q^a\). Distinct cyclotomic factors have disjoint roots. A negative
\(E_d\) would therefore give a genuine pole and prevent polynomiality;
nonnegative exponents give the integer polynomial
\eqref{q1:qcc:ratio-cyclotomic}. The degree is the difference of the sums
\(1+\cdots+a_i\) and \(1+\cdots+b_j\); balancing cancels their linear
parts. Cancelling \(1-q\) from every individual factor before taking
\(q\to1\) gives the factorial value. Finally,
\(1-q^{-r}=-q^{-r}(1-q^r)\), so the same degree and sign calculation
gives reciprocity. The proof of \cref{q1:qcc:root-jets} applies to this
multiset of factors without change.
\end{proof}

Polynomiality does not imply coefficientwise positivity. A small balanced
example is
\[
 \frac{(q;q)_6(q;q)_1^2}{(q;q)_5(q;q)_3}
 =\frac{(1-q^6)(1-q)}{(1-q^3)(1-q^2)}
 =1-q+q^2=\Phi_6(q).
\]
Thus the finite-field and partition interpretations of Gaussian coefficients
supply information beyond the cyclotomic criterion. Conversely, the criterion
extends the local-series machinery to examples for which no positive counting
interpretation is available.
% ---- Q3 lines 1899--2123 ----
\chapter{Three distinct asymptotic regimes}
\label{q3:sec:asymptotic}
A formula that is uniform for fixed $q$ need not be uniform as $q\to1$.
This section separates three useful regimes instead of silently mixing
them. Finite coefficient formulae give an exact expansion in the first,
infinite products give convergent exponential corrections in the
second, and Euler--Maclaurin gives a uniform Poincar\'e expansion in the
third. The last calculation extends the canonical manuscript's asymptotic
coefficient machinery to a Gaussian input.

\section{Fixed lower index and fixed base}
For fixed $k\ge0$ and $0<|q|<1$, the finite $q$-binomial theorem gives
\begin{equation}
 \boxed{\displaystyle
 \G nk=\frac1{(q;q)_k}
      \sum_{j=0}^k(-1)^jq^{j(n-k+1)+\binom j2}\G kj,
       \qquad n\ge k.}
 \label{q3:eq:fixed-k-exact}
\end{equation}
This is simply the product
$(q^{n-k+1};q)_k/(q;q)_k$. It is an \emph{exact finite} expansion in
powers of $q^n$, not an infinite asymptotic ansatz. In particular,
\[
 \G nk\longrightarrow\frac1{(q;q)_k}\quad(n\to\infty).
\]
For $k\ge1$ the first correction is
$-q^{n-k+1}[k]_q/(q;q)_k$, with the remaining finite terms explicitly
given by \cref{q3:eq:fixed-k-exact}.

\section{Both sides large and fixed base}
Fix $0<|q|<1$, and let $k,n-k\to\infty$. Put
$a=q^{k+1}$, $b=q^{n-k+1}$, and $c=q^{n+1}=ab/q$.
Since $(q;q)_m=(q;q)_\infty/(q^{m+1};q)_\infty$,
\begin{equation}
 \G nk=\frac1{(q;q)_\infty}
              \frac{(a;q)_\infty(b;q)_\infty}{(c;q)_\infty}.
 \label{q3:eq:bulk-fixed-q-product}
\end{equation}
The infinite $q$-binomial theorem therefore gives the exact convergent
expansion
\begin{equation}
 \G nk=\frac1{(q;q)_\infty}
 \sum_{r,s,t\ge0}
   \frac{(-1)^{r+s}q^{\binom r2+\binom s2}a^rb^sc^t}
        {(q;q)_r(q;q)_s(q;q)_t}.
 \label{q3:eq:bulk-fixed-q-triple}
\end{equation}
The two numerator series converge absolutely for every finite $a,b$;
the reciprocal series converges absolutely for $|c|<1$, which holds
here. Hence their product and rearrangement are justified.
The logarithm of the relative correction is
\begin{equation}
 \log\bigl((q;q)_\infty\G nk\bigr)
   =-\sum_{j\ge1}\frac{a^j+b^j-c^j}{j(1-q^j)}.
 \label{q3:eq:bulk-fixed-q-log}
\end{equation}
It follows by expanding the logarithm of each product; the absolute
convergence for $|a|,|b|,|c|<1$ justifies interchanging sums. In
particular,
\begin{equation}
 \G nk=\frac1{(q;q)_\infty}
       \left(1-\frac{a+b}{1-q}
                       +O(|a|^2+|b|^2)\right).
 \label{q3:eq:bulk-fixed-q-leading}
\end{equation}
The constant may depend on the fixed base $q$; the $c$ term is of order
$ab$. Higher corrections can be generated either by the triple sum or
by the ordinary Bell expansion of \cref{q3:eq:bulk-fixed-q-log}. None of
these fixed-base formulae asserts uniformity near $q=1$.

\section{The double scaling \texorpdfstring{$q=e^{-\tau/n}$}{q = exp(-tau/n)}}
Let $k=k_n$ be an integer, set $\alpha=k/n$ and $b=1-\alpha$, and
assume
\begin{equation}
 \varepsilon\le\alpha\le1-\varepsilon,\qquad
 0<\tau_0\le\tau\le\tau_1<\infty,
 \qquad q=e^{-\tau/n}.
 \label{q3:eq:double-scaling-domain}
\end{equation}
Thus $\alpha$ denotes the actual ratio $k/n$, including any rounding;
there is no unrecorded floor correction. Define a smooth function on
$[0,1]$ by
\begin{equation}
 h_\tau(x)=\log\frac{1-e^{-\tau x}}{\tau x},\qquad h_\tau(0)=0,
 \label{q3:eq:h-tau}
\end{equation}
and set
\begin{align}
 \mathscr S_\tau(\alpha)
 &=-\alpha\log\alpha-b\log b
      +\int_0^1h_\tau(x)\dd x
      -\int_0^\alpha h_\tau(x)\dd x
      -\int_0^b h_\tau(x)\dd x
 \label{q3:eq:rate-integral}\\
 &=\frac{\Li_2(e^{-\tau})-\Li_2(e^{-\tau\alpha})
                  -\Li_2(e^{-\tau b})+\pi^2/6}{\tau}.
 \label{q3:eq:rate-dilog}
\end{align}
The equality follows from
$\frac{\dd}{\dd x}\Li_2(e^{-\tau x})
=\tau\log(1-e^{-\tau x})$ and elementary integration of
$\log(\tau x)$. One can take \cref{q3:eq:rate-integral} as the definition
without using a special function.

\begin{theorem}[Uniform all-order logarithmic expansion]\label{q3:thm:double-scaling}
Under \cref{q3:eq:double-scaling-domain}, for every fixed $M\ge0$,
\begin{align}
 \log\Gb nk{e^{-\tau/n}}
 ={}&n\mathscr S_\tau(\alpha)-\frac12\log(2\pi n\alpha b)
       +\frac12\bigl(h_\tau(1)-h_\tau(\alpha)-h_\tau(b)\bigr)
       \notag\\
 &+\sum_{r=1}^M\frac{C_{2r-1}(\alpha,\tau)}{n^{2r-1}}
       +O\bigl(n^{-(2M+1)}\bigr),
 \label{q3:eq:double-scaling-log}
\end{align}
with an error uniform on the indicated compact parameter set. The
coefficients are explicit:
\begin{align}
 C_{2r-1}(\alpha,\tau)
 ={}&\frac{B_{2r}}{2r(2r-1)}
            \bigl(1-\alpha^{1-2r}-b^{1-2r}\bigr) \notag\\
 &+\frac{B_{2r}}{(2r)!}
  \bigl(h_\tau^{(2r-1)}(1)-h_\tau^{(2r-1)}(\alpha)
       -h_\tau^{(2r-1)}(b)+h_\tau^{(2r-1)}(0)\bigr).
 \label{q3:eq:double-scaling-C}
\end{align}
There are no even negative powers in this logarithmic expansion.
\end{theorem}
\begin{proof}
The product identity has the exact decomposition
\begin{equation}
 \log\Gb nk{e^{-\tau/n}}
  =\log\binom nk+
      \sum_{j=1}^n h_\tau(j/n)
      -\sum_{j=1}^k h_\tau(j/n)
      -\sum_{j=1}^{n-k}h_\tau(j/n).
 \label{q3:eq:double-scaling-exact-decomposition}
\end{equation}
Indeed $1-e^{-\tau j/n}=(\tau j/n)e^{h_\tau(j/n)}$, and the elementary
factors cancel to the ordinary binomial coefficient.

For an endpoint $u\in\{1,\alpha,b\}$ with $un$ integral,
Euler--Maclaurin on mesh $1/n$ reads
\begin{align*}
 \sum_{j=1}^{un}h_\tau(j/n)
 ={}&n\int_0^u h_\tau(x)\dd x+\frac{h_\tau(u)-h_\tau(0)}2\\
 &+\sum_{r=1}^M\frac{B_{2r}}{(2r)!n^{2r-1}}
       \bigl(h_\tau^{(2r-1)}(u)-h_\tau^{(2r-1)}(0)\bigr)
       +O(n^{-(2M+1)}).
\end{align*}
To justify this error with the displayed truncation, apply the usual
periodic-Bernoulli remainder formula one order farther, then absorb
both the next boundary term and the remainder into $O(n^{-(2M+1)})$.
The derivatives of $h_\tau$ of every fixed order are uniformly bounded
on $[0,1]\times[\tau_0,\tau_1]$, because its apparent singularity at
zero is removable. Hence the bounds are uniform. The underlying
Euler--Maclaurin formula follows by integration by parts successively
with periodic Bernoulli functions; its standard remainder form is
recorded in \cite[\S2.10(i)]{dlmfem}.

Finally the factorial Stirling expansion
\cite[\S5.11]{dlmfgamma} gives
\begin{align*}
 \log\binom nk={}&n(-\alpha\log\alpha-b\log b)
                  -\frac12\log(2\pi n\alpha b)\\
 &+\sum_{r=1}^M\frac{B_{2r}}{2r(2r-1)n^{2r-1}}
       (1-\alpha^{1-2r}-b^{1-2r})+O(n^{-(2M+1)}).
\end{align*}
Its remainder is uniform since $k,n-k\ge\varepsilon n$.
Insert these three Euler--Maclaurin expansions and the Stirling
expansion into \cref{q3:eq:double-scaling-exact-decomposition}; collect
the integral, endpoint, and odd-power terms. This gives precisely
\cref{q3:eq:double-scaling-log,q3:eq:double-scaling-C}.
\end{proof}

\section{Explicit derivatives and multiplicative coefficients}
No recursive differentiation is needed in \cref{q3:eq:double-scaling-C}.
For $x>0$ and $s\ge1$,
\begin{equation}
 h_\tau^{(s)}(x)=(-1)^{s-1}
  \left(\tau^s\Li_{1-s}(e^{-\tau x})-\frac{(s-1)!}{x^s}\right).
 \label{q3:eq:h-tau-derivative}
\end{equation}
It follows by differentiating
$\log(1-e^{-\tau x})-\log(\tau x)$ and using
\cref{q3:eq:Rr}. At zero the nonsingular values are
\begin{equation}
 h_\tau'(0)=-\frac\tau2,\quad
 h_\tau^{(2r)}(0)=\frac{B_{2r}\tau^{2r}}{2r},\quad
 h_\tau^{(2r+1)}(0)=0\quad(r\ge1),
 \label{q3:eq:h-tau-zero}
\end{equation}
by the Bernoulli logarithm expansion. Equations
\cref{q3:eq:Rr-stirling,q3:eq:h-tau-derivative,q3:eq:h-tau-zero}
therefore make every $C_j$ a finite expression in rational functions
of $e^{-\tau}$, $e^{-\tau\alpha}$, $e^{-\tau b}$ and the parameters.

Exponentiating the leading terms simplifies the prefactor:
\begin{equation}
 \Gb nk{e^{-\tau/n}}
 \sim e^{n\mathscr S_\tau(\alpha)}
     \sqrt{\frac{\tau(1-e^{-\tau})}
         {2\pi n(1-e^{-\tau\alpha})(1-e^{-\tau b})}}
     \sum_{m\ge0}\frac{H_m(\alpha,\tau)}{n^m},
 \label{q3:eq:double-scaling-multiplicative}
\end{equation}
where $C_{2r}=0$ and
\begin{equation}
 H_m=\frac1{m!}\Yc_m(1!C_1,2!C_2,\ldots,m!C_m)
    =\sum_{\sum jm_j=m}\prod_{j=1}^m\frac{C_j^{m_j}}{m_j!}.
 \label{q3:eq:double-scaling-H}
\end{equation}
The first terms are
\[
 1+\frac{C_1}{n}+\frac{C_1^2}{2n^2}
   +\frac{C_3+C_1^3/6}{n^3}
   +\frac{C_1C_3+C_1^4/24}{n^4}+\cdots.
\]
This is an all-order Poincar\'e expansion with the compact-parameter
uniformity stated above, not a claim of convergence of the inverse-$n$
series. Its formal $\tau\to0$ limit recovers ordinary binomial Stirling
asymptotics: $h_\tau\to0$, the rate tends to the entropy, and the
prefactor tends to $(2\pi n\alpha b)^{-1/2}$. Endpoint regimes
$\alpha\to0,1$ require a different uniform analysis and are not included
in the theorem.
% ---- Q1 lines 2356--2444 ----
\chapter{Computation, verification, and integration}
\label{q1:qcc:computation}
\section{Choosing a representation for the task}
The formulae above are equal algebraically, but they serve different
computational purposes. A recurrence-free expression is a useful mathematical
specification; it is not automatically the fastest or best-conditioned numerical
algorithm.

For an exact Gaussian polynomial, the division-free Pascal recurrence is usually
a natural baseline: it uses only integer polynomial addition and shifts. For a
single coefficient in the base, \cref{q1:qcc:gaussian-coefficients} supplies a
finite divisor-sum Bell expression or a restricted-partition expression. For a
coefficient of a product in an auxiliary variable, the geometric-product master
formula avoids expanding all factors separately. One first computes a finite
list of power sums and then evaluates one Bell polynomial.

For symbolic derivatives at a generic parameter, the logarithmic derivatives
\(L_r\) in \cref{q1:qcc:generic-jets} compress large product-rule expressions.
At one, Bernoulli cumulants are simpler. At another root of unity, use the
constant, multiplicity, and regular logarithmic data \((C_\zeta,\nu,\lambda_r)\)
from \cref{q1:qcc:root-jets}. Do not use a product quotient as a numerical
\(0/0\) and hope that floating-point cancellation will reconstruct the limit.

For basis conversion, the triangular \(q\)-Stirling recurrences give a
convenient algorithm, while the finite Gaussian sum and Bell diagonal formulae
give independent specifications. For implicit series, compute the logarithmic
power sums of the defining product and apply \cref{q1:qcc:inverse-product}; the
ordinary Lagrange factor remains \(1/n\). For extrapolation, compute the
weights directly from \eqref{q1:qcc:weights} and use the error and noise formulae
together, rather than increasing the cancellation order without controlling
sample accuracy.

\section{Arithmetic cost and finite verification}
Suppose the power sums \(P_1,\ldots,P_R\) are available. The coefficients
\(f_0,\ldots,f_R\) of
\(F(z)=\exp(\sum_{r\ge1}P_rz^r/r)\) can be evaluated through
\begin{equation}\label{q1:qcc:coefficient-algorithm}
 f_0=1,\qquad nf_n=\sum_{r=1}^{n}P_rf_{n-r}.
\end{equation}
This uses \(O(R^2)\) arithmetic operations and \(O(R)\) stored scalars.
It follows by differentiating \(F\) and comparing coefficients in
\(zF'(z)=F(z)\sum_{r\ge1}P_rz^r\). The algorithm is recurrent, but
its output is specified independently by the recurrence-free Bell sum
\eqref{q1:qcc:master-coefficient}. The distinction permits efficient calculation
and an independent formula audit at the same time. Bit complexity also depends
on coefficient heights and is not described by the arithmetic-operation count.

The companion \texttt{verify\_identities.py} uses only Python's standard
library, integer polynomial arithmetic, and exact rational numbers. Its successful
run for this volume reports
\begin{center}
\fbox{\large\bfseries 10,912 exact checks in 36 groups.}
\end{center}
The tests include Gaussian products and convolution, both terminating
\(q\)-Chu--Vandermonde sums, divisor identities, Jackson stencils and Taylor
reconstruction, \(q\)-Stirling sums and Bell diagonals, weighted composition,
Lagrange coefficients, cyclotomic factors, parameter cumulants, all-root local
jets, and ordinary and logarithmic extrapolation filters.

The checks deliberately use different representations. Gaussian polynomials are
constructed by polynomial Pascal addition and compared with product values;
Bell expressions are checked against direct finite products or moments;
Lagrange coefficients are compared with ordinary truncated fixed-point
substitution; and \(q\)-Lucas is checked as exact divisibility by a
cyclotomic polynomial. Root-of-unity jets are checked not just at \(-1\),
but in the exact fields \(\Q(\zeta_d)\), \(3\le d\le8\), using
polynomials modulo \(\Phi_d\). For \(u\) of exact order \(L>1\),
the finite certificate
\[
 (1-u)^{-1}=-\frac1L\sum_{j=1}^{L-1}ju^j
\]
provides exact division in these calculations: multiplication by \(1-u\)
reduces the numerator sum to \(-L\).

Polynomial coefficient and moment tests cover \(0\le n\le12\), all
\(0\le k\le n\), and orders through eight. The \(-1\) tests cover
orders through nine. The cyclotomic-field jet tests cover \(0\le n\le9\),
all lower indices, and orders through seven. Other test bounds and the exact
rational parameter values are specified directly in the program. The archive
includes the complete output as \texttt{validation\_report.txt}.

\begin{auditbox}[title={What this verification establishes}]
These deterministic finite checks are reproducible audits of signs, shifts,
normalizations, and coefficient formulae. They are not proofs for unbounded
indices, do not establish floating-point stability outside the stated bounds,
and do not assert Lean verification. The mathematical proofs in the preceding
chapters establish the general statements. No current declaration inventory or
formalization status of the linked repository is inferred from these tests.
\end{auditbox}
% ---- hand chunk 24-verifiers.tex ----

\section{The three retained verification programs}
\label{q3:sec:verifiers}
Each source shipped an exact, deterministic verification program, and all
three are retained unchanged under \path{verification/} in the package
directory, together with the reports of the runs that accompanied the
sources. They overlap in what they check, and they were written
independently, so agreement between them is evidence about the
normalizations rather than about a shared implementation.
\begin{itemize}
\item \path{verification/source1/verify_identities.py} is the program
 described in the preceding section: \(10{,}912\) exact checks in
 \(\Q(q)\) at symbolic and rational bases, covering the identities of the
 spine of this volume, with the run recorded in
 \path{validation_report.txt}.
\item \path{verification/source2/verify_identities.py} performs
 \(2{,}693\) exact assertions in four groups (algebra, coefficients,
 inversion, roots), with Gaussian polynomials built from subset weights
 rather than from the identities under test; product expansions are compared
 with partition kernels, fixed-point series with the Lagrange and coupled
 formulae, and polynomial derivatives with the Bernoulli--Bell and
 regularized root formulae. Its machine-readable reports are the five
 \path{verification_*.json} files beside it.
\item \path{verification/source3/verify_identities.py} performs
 \(10{,}267\) exact checks using only the standard library: Gaussian
 products against Pascal recurrences and subset enumeration,
 \(q\)-Stirling finite sums against node recurrences, \(q\)-Eulerian
 coefficients against direct permutation enumeration, Jackson chain rules
 against barycentric divided differences, inverse coefficients against
 independent series substitution, moments from cumulants against direct
 products of elementary uniform moment-generating series, Lucas evaluations
 as polynomial congruences, and local derivatives at \(q=-1\) after removing
 the zero factor. Its run is recorded in \path{verification_results.txt}.
 An optional final experiment, requiring \texttt{mpmath}, checks the
 double-scaling orders numerically; its output is reproduced next.
\end{itemize}
Finite testing cannot establish a theorem for arbitrary indices; the proofs
carry that burden. Conversely, a failure would expose a transcription,
convention, or implementation error. These complementary roles are
particularly valuable for the powers of \(q\) in inversion formulae and for
the distinction between \(S_q(n,k)\) and \(q^{\binom k2}S_q(n,k)\).


% ---- Q3 lines 2328--2353 ----
\section{A numerical check of the double-scaling order}
An optional final experiment uses \texttt{mpmath} with 70 decimal
digits. For $\tau=1.3$ and $\alpha=1/3$, let $L_0(n)$ be the leading
line of \cref{q3:eq:double-scaling-log}, and define
\[
 R_1(n)=\log\Gb n{n/3}{e^{-1.3/n}}-L_0(n)-C_1/n,
 \qquad R_2(n)=R_1(n)-C_3/n^3.
\]
The actual output, rounded here, is
\begin{center}
\begin{tabular}{@{}rrr@{}}
\toprule
$n$ & $n^3R_1(n)$ & $n^5R_2(n)$\\
\midrule
60  & $0.0815396566910922$ & $-0.1977273873208401$\\
120 & $0.0815808310074465$ & $-0.1979993937823302$\\
240 & $0.0815911422914091$ & $-0.1980676188797450$\\
\midrule
Predicted limit & $C_3=0.0815945809653480$ & $C_5=-0.1980903805565352$\\
\bottomrule
\end{tabular}
\end{center}
The approach to the two predicted constants checks the signs and
orders in the formulae. These are high-precision numerical observations,
not interval-certified inequalities. The uniform error estimate is
proved independently in \cref{q3:sec:asymptotic}.
% ---- Q2 lines 1608--1627 ----
\section{Normalization audit}
\begin{longtable}{@{}>{\RaggedRight\arraybackslash}p{0.25\textwidth}>{\RaggedRight\arraybackslash}p{0.34\textwidth}>{\RaggedRight\arraybackslash}p{0.34\textwidth}@{}}
\toprule
\textbf{Operation} & \textbf{Correct object} & \textbf{Incorrect shortcut}\\
\midrule
\endfirsthead
\toprule
\textbf{Operation} & \textbf{Correct object} & \textbf{Incorrect shortcut}\\
\midrule
\endhead
Finite product & $q^{\binom k2}\qb Nk$ & Omitting $q^{\binom k2}$\\[1mm]
Gaussian inverse & $(-1)^r q^{\binom r2}\qb nr$ & Replacing the inverse by $T_{-c}$\\[1mm]
Ordinary composition & Ordinary $k!/\prod m_j!$ inside \eqref{q1:qcc:weighted-bell-def} & Replacing every profile multiplicity by a Gaussian multinomial\\[1mm]
Lagrange--B\"urmann & $m/n$ and an ordinary derivative & Replacing $n$ by $\qi n$ without a new theorem\\[1mm]
Newton limit & $q$-integer nodes tend to integers & Treating geometric nodes as if they did not coalesce\\[1mm]
Root-of-unity values & Polynomial evaluation or regularization & Substituting into an uncancelled factorial quotient\\[1mm]
$q$-Catalan family & Its defining equation and normalization & Assuming all Gaussian and area refinements coincide\\[1mm]
Infinite product & $q$-adic, logarithmic, or analytic interpretation & Claiming stabilization in the $z$-adic topology over $\Q(q)$\\
\bottomrule
\end{longtable}
% ---- Q2 lines 1629--1649 ----
\section{Where the extension attaches}
The following crosswalk uses labels observed in the canonical source, not its
mutable page or theorem numbers. Labels are pointers to mathematical interfaces,
not assertions that the new theorems have existing Lean counterparts.
\begin{longtable}{@{}>{\RaggedRight\arraybackslash}p{0.36\textwidth}>{\RaggedRight\arraybackslash}p{0.57\textwidth}@{}}
\toprule
\textbf{Base-source interface} & \textbf{Extension developed here}\\
\midrule
\endfirsthead
\toprule
\textbf{Base-source interface} & \textbf{Extension developed here}\\
\midrule
\endhead
\texttt{thm:newton-\allowbreak{}expansion} & Two node systems, explicit Gaussian interpolation, and the specified $q$-Stirling inverse arrays (\cref{q3:sec:newton}).\\[1mm]
\texttt{thm:bell-\allowbreak{}symmetric-\allowbreak{}functions} & Geometric-alphabet specialization and the ramified shifted-factorial Bell kernel (\cref{q1:qcc:bell-master}).\\[1mm]
\texttt{thm:merged-\allowbreak{}binomial-\allowbreak{}inversion} & Gaussian inverse weights and subspace-lattice M\"obius interpretation (\cref{q2:sec:transforms}).\\[1mm]
\texttt{thm:merged-\allowbreak{}weighted-\allowbreak{}binomial-\allowbreak{}translation} & The corrected composition law; the ordinary additive law does not persist (\cref{q2:sec:transforms}).\\[1mm]
\texttt{eq:merged-\allowbreak{}bell-\allowbreak{}normalization} & Coordinate-specific Bell conversion, with ordinary profile multiplicities retained (\cref{q1:qcc:composition}).\\[1mm]
\texttt{thm:lagrange-\allowbreak{}burmann} & Recurrence-free product-weighted inverse coefficients and coupled inverses (\cref{q1:qcc:lagrange}).\\
\bottomrule
\end{longtable}
% ---- Q1 lines 2446--2482 ----
\section{A theorem-level insertion map}
The following map identifies the mathematical interface with the source without
assuming that its chapter numbering will remain fixed.
\begin{longtable}{@{}p{0.31\textwidth}p{0.63\textwidth}@{}}
\toprule
Source mechanism & Extension supplied here\\
\midrule
\endfirsthead
\toprule
Source mechanism & Extension supplied here\\
\midrule
\endhead
Newton expansion and factorial bases &
Geometric Jackson--Newton interpolation, explicit Gaussian stencils, and the
\([j]_q\)-node \(q\)-Stirling arrays
(\cref{q1:qcc:jackson,q1:qcc:qstirling}).\\[2mm]
Exponential formula and Bell polynomials &
Finite-alphabet power sums, arbitrary geometric-product exponents, determinants,
and inverse logarithmic coordinates (\cref{q1:qcc:bell-master}).\\[2mm]
Binomial inversion and incidence algebra &
Gaussian inversion, subspace-lattice M\"obius coefficients, and positive
linearization (\cref{q1:qcc:convolution}).\\[2mm]
Stirling, Bernoulli, and Eulerian families &
Parameter derivatives and cumulants, ordinary/logarithmic derivative conversion,
and regular root-of-unity factors
(\cref{q1:qcc:parameter-jets,q1:qcc:roots}).\\[2mm]
Riordan and composition calculus &
A specified \(q\)-factorial normalization of ordinary composition and its
matrix group law (\cref{q1:qcc:composition}).\\[2mm]
Lagrange--B\"urmann and Good inversion &
Product-defined inverses with finite Bell coefficients and a fully computed
cross-coupled two-variable case (\cref{q1:qcc:lagrange}).\\[2mm]
Asymptotic coefficient operators &
Geometric spectral filters, complete error coefficients, noise norms, and
confluent cancellation of logarithmic modes (\cref{q1:qcc:filters}).\\
\bottomrule
\end{longtable}
% ---- hand chunk 26-analytic-head.tex ----

\section{Formal identities, analytic domains, and specialization}
\label{q2:sec:analytic}
The domain remarks scattered through the audit boxes of this volume are
collected here in one place, together with a checklist of what survives
specialization of \(q\).


% ---- Q2 lines 1536--1577 ----
\subsection{What convergence is, and is not, required}
Finite Gaussian identities, triangular transforms, and coefficient extraction in
$R[[z]]$ need no analytic convergence. Formal exponentials, logarithms, and divided
Lagrange formulas do require the indicated rational denominators. For products
with infinitely many geometric factors, two valid settings have been used:
coefficientwise $q$-adic products, or rational logarithmic coefficients such as
$S_\infty(q^j)=1/(1-q^j)$. Neither should be mislabeled as ordinary $z$-adic
convergence of the finite products.

For a fixed complex $|q|<1$, $\poch{az}q\infty$ converges locally uniformly in
$z$: on each compact set, $\sum_{j\ge0}|azq^j|$ converges uniformly. Away from its
zeros, its reciprocal is analytic. Fractional powers of a unit product have a
unique analytic branch near $z=0$ with value $1$, agreeing with the formal power
used here. The logarithm computation is justified in a sufficiently small disk
by absolute convergence, and the identities elsewhere follow by analytic
continuation where both sides are defined. These conditions recover the standard
analytic $q$-binomial theorem \cite{dlmf172,dlmf175}.

If $\Phi$ in \cref{q1:qcc:lagrange-theorem} is analytic near zero and $\Phi(0)=1$, the
analytic inverse-function theorem makes its formal solution an analytic germ.
The displayed coefficient formulas are then its actual Taylor coefficients.
No claim is made that they converge past the first singularity of that germ.
The two-variable system has the analogous local analytic interpretation because
its implicit Jacobian at the origin is the identity.

\subsection{A specialization checklist}
Polynomial identities such as \eqref{q1:qcc:finite}, \eqref{q1:qcc:vandermonde-formula},
\eqref{q2:eq:transform-inverse}, and \eqref{q1:qcc:lucas-eq} survive specialization at any
$q$. By contrast, three types of rational formulas need attention.
First, interpolation on geometric or $q$-integer nodes requires distinct nodes;
colliding nodes require a confluent limit rather than raw substitution.
Second, $q$-divided-power coordinates need invertible $\qf j$, so a coordinate
formula can have genuine poles even when a separately defined polynomial family
does not. Third, infinite-factor logarithms involve $1-q^j$ in denominators,
and a root-of-unity specialization is not available without a separate limiting
argument or a proved cancellation in the final coefficients.

The safe procedure for a Gaussian polynomial is to use its subset expansion,
its cyclotomic factorization, or \cref{q1:qcc:root-jets}. The fact that an intermediate
formula has $0/0$ terms is not evidence that the Gaussian polynomial is undefined.
Conversely, polynomiality of a Gaussian coefficient does not regularize every
expression built from reciprocals of $q$-factorials.
% ---- Q3 lines 2374--2396 ----
\section{Directions not claimed as results here}
The framework suggests several further investigations. First, one can
replace one geometric alphabet by several alphabets with different bases.
The Bell-product theorem continues immediately after replacing
$[N]_{q^j}$ by the corresponding sum of geometric power sums, but a
single Gaussian connection array need no longer exist. A useful problem
is to classify which node changes preserve particularly simple inverse
matrices.

Second, ordinary multivariate Lagrange inversion can accept the same
finite $q$-product inputs. The determinant and Jacobian factors remain
ordinary; only their coefficients acquire Gaussian or geometric-power
structure. Working out a systematic multivariate analogue of
\cref{q1:qcc:inverse-product-formula} would extend the canonical manuscript's coefficient
calculus without conflating it with noncommutative substitution.

Third, the regimes $q\to\zeta$ with $n\to\infty$, or
$q=e^{-\tau/n}$ with $k/n\to0$, lie outside the uniform theorem proved
here. The local cyclotomic decomposition and the separate fixed-$k$
formula provide starting points for matched expansions, but no uniform
transition theorem for those regimes is asserted. In particular, the
present inverse-$n$ expansion is not advertised as a complete
exponentially improved transseries.
% ---- Q1 lines 2484--2499 ----
\section{Concluding perspective}
Gaussian coefficients occur here in three complementary roles. They are
coefficients of elementary and complete symmetric functions on geometric
alphabets; they are kernels for triangular basis changes and inversions; and
they are the explicit weights of spectral polynomials on geometric grids.
Ordinary Bell polynomials connect these roles by turning logarithmic power sums
into finite coefficients. Stirling numbers change the derivative coordinate,
Bernoulli numbers remove the apparent singularity at one, and Eulerian
polynomials make the remaining logarithmic derivatives rational.

This common mechanism is more informative than a formal replacement of every
integer by a \(q\)-integer. It explains which structures survive unchanged,
which acquire a Gaussian kernel, and which require a different operation.
It also separates exact finite algebra from local analysis and asymptotic
remainder estimates. Those separations are what allow the extension to serve
as a reusable part of the canonical manuscript's coefficient calculus.
% ---- Q1 lines 2501--2568 ----
\appendix
\chapter{A compact formula atlas}
\label{q1:qcc:atlas}
This atlas records the input and output of the principal coefficient mechanisms.
All symbols and domains are defined in the referenced results.

\section*{Finite geometric products}
For
\(F(z)=\prod_\ell(a_\ell z;q^{d_\ell})_{N_\ell}^{-\alpha_\ell}\), put
\(P_r=\sum_\ell\alpha_\ell a_\ell^r[N_\ell]_{q^{d_\ell r}}\). Then
\[
 [z^n]F(z)=\frac1{n!}\Y n\bigl(((r-1)!P_r)_{r=1}^{n}\bigr)
 =\sum_{\sum rm_r=n}\prod_{r=1}^{n}\frac{P_r^{m_r}}{r^{m_r}m_r!}.
\]
This is \cref{q1:qcc:master-theorem}. For an implicit solution \(y=zF(y)\),
replace \(P_r\) by \(nP_r\), use Bell degree \(n-1\), and divide by
\(n!\), as in \cref{q1:qcc:inverse-product}.

\section*{Gaussian polynomial coefficients}
For \(m=n-k\), set
\[
 A_r=\sum_{\substack{d\mid r\\1\le d\le k}}d
     -\sum_{\substack{d\mid r\\m<d\le n}}d.
\]
Then
\[
 [q^R]\GB nk=\frac1{R!}
       \Y R\bigl(((r-1)!A_r)_{r=1}^{R}\bigr).
\]
The two-block divisor kernels are the multiset difference of denominator and
numerator exponents; see \cref{q1:qcc:gaussian-coefficients}.

\section*{The \(q=1\) expansion}
Let
\(\kappa_1=k(n-k)/2\) and
\(\kappa_r=(\beta_r/r)\sum_{j=1}^{k}((n-k+j)^r-j^r)\), \(r\ge2\).
Then
\[
 \GB[e^t]nk=\binom nk\sum_{r\ge0}\Y r(\kappa_1,\ldots,\kappa_r)
                                      \frac{t^r}{r!}.
\]
For ordinary \(q\)-derivatives, apply signed Stirling coefficients as in
\eqref{q1:qcc:ordinary-jets-one}.

\section*{All roots of unity}
Deflate every factor \(1-\zeta^ae^{at}\), using
\(-at\) when \(\zeta^a=1\) and \(1-\zeta^a\) otherwise. The constant
product is \(C_\zeta\), the net order is \(\nu\), and the regular
logarithmic derivatives are \(\lambda_r\). Then
\[
 \GB[\zeta e^t]nk=C_\zeta t^\nu
      \sum_{r\ge0}\Y r(\lambda_1,\ldots,\lambda_r)\frac{t^r}{r!}.
\]
Equations \eqref{q1:qcc:eta}--\eqref{q1:qcc:root-lambda} supply every input by
finite sums, and \eqref{q1:qcc:root-ordinary-derivatives} converts to ordinary
Taylor coefficients.

\section*{Geometric extrapolation}
For \(0<\rho<1\), the weights
\[
 w_{N,j}=\frac{(-1)^{N-j}\rho^{\binom{N-j+1}{2}}
                  \GB[\rho]Nj}{(\rho;\rho)_N}
\]
preserve constants, kill the first \(N\) powers, and have exact absolute norm
\(( -\rho;\rho)_N/(\rho;\rho)_N\). Their response to every higher power
is \((-1)^N\rho^{\binom{N+1}{2}}\GB[\rho]{m-1}{N}\).
Repeated spectral roots cancel logarithmic modes; see
\cref{q1:qcc:confluent-filter}.
% ---- hand chunk 27-guide-head.tex ----

\chapter{A formula-selection guide}
\label{q3:app:guide}
The following guide lists the finite mechanisms most useful for exact
calculations. Each entry points to a formula with all indices and
normalizations specified in the main text; it complements the atlas of
\cref{q1:qcc:atlas}, which records inputs and outputs.

% ---- Q3 lines 2425--2478 ----
The following guide lists the finite mechanisms most useful for exact
calculations. Each entry points to a formula with all indices and
normalizations specified in the main text.
\begin{longtable}{@{}>{\raggedright\arraybackslash}p{0.32\textwidth}>{\raggedright\arraybackslash}p{0.60\textwidth}@{}}
\toprule
Task & Formula and applicable setting\\
\midrule
\endfirsthead
\toprule
Task & Formula and applicable setting\\
\midrule
\endhead
Evaluate a Gaussian polynomial & Subset or rectangle sums
\eqref{q3:eq:subset}--\eqref{q3:eq:rectangle}; no division, hence valid at every
complex $q$.\\[3pt]
Convolve or invert a Gaussian transform &
Vandermonde \eqref{q1:qcc:vandermonde-formula}, multiblock
\eqref{q1:qcc:block-vandermonde}, inversion \eqref{q1:qcc:inverse-pair}.\\[3pt]
Change from geometric to affine-geometric nodes &
The finite mixed ordinary/Gaussian sum \eqref{q3:eq:affine-geometric}.\\[3pt]
Compute $q$-Stirling entries without recurrence &
\eqref{q3:eq:qstirling-gaussian}, \eqref{q1:qcc:qstirling-sum},
\eqref{q3:eq:qstirling-first-gaussian}; polynomial cancellation at exceptional
bases.\\[3pt]
Convert Stirling and Eulerian data &
\eqref{q3:eq:stirling-eulerian} and \eqref{q3:eq:eulerian-stirling}, using the
major-index convention \eqref{q3:eq:qeulerian}.\\[3pt]
Extract a coefficient of mixed $q$-products &
The Bell or finite profile formulae
\eqref{q1:qcc:master-coefficient}.\\[3pt]
Differentiate on a geometric grid &
Finite Jackson sum \eqref{q1:qcc:q-stencil}; product
\eqref{q1:qcc:q-leibniz}; path chain rule \eqref{q3:eq:dd-chain}.\\[3pt]
Compose $q$-normalized ordinary series &
\eqref{q1:qcc:weighted-bell-finite} and \eqref{q1:qcc:composition-formula}; the partial
coefficients can be rational in $q$.\\[3pt]
Invert an ordinary series with finite $q$-product input &
\eqref{q1:qcc:inverse-product-formula}; the prefactor is
ordinary $r/n$.\\[3pt]
Compute derivatives at $q=1$ &
Bernoulli cumulants \eqref{q1:qcc:cumulant-expansion} and the finite
Stirling--Bell sum \eqref{q1:qcc:ordinary-jets-one}.\\[3pt]
Evaluate or expand at a root of unity &
Lucas \eqref{q1:qcc:lucas-eq}; regularized local jet
\eqref{q3:eq:root-local-expansion}.\\[3pt]
Extract moving-base asymptotic coefficients &
\eqref{q3:eq:double-scaling-C} and \eqref{q3:eq:double-scaling-H}, under the
compactness assumptions \eqref{q3:eq:double-scaling-domain}.\\[3pt]
Remove a finite geometric tail polynomial &
Gaussian weights \eqref{q1:qcc:weights}, exact residual
\eqref{q1:qcc:filter-error-series}, moment application
\eqref{q3:eq:moment-extrapolation}.\\
\bottomrule
\end{longtable}
% ---- hand chunk 28-provenance.tex ----

\chapter{Source provenance, notation dictionary, and reproducibility}
\label{q1:qcc:provenance}
\section{Sources and consolidation}
This volume consolidates the three \(q\)-binomial coefficient-calculus
reports filed in the repository's frontier intake of September 4, 2026,
under \path{drafts/combinatorial-coefficient-calculus/}: the packages
\path{q_binomial_coefficient_calculus}, \path{q_binomial_coefficient_calculus-2},
and \path{q_binomial_coefficient_calculus-3}. Their \LaTeX{} sources, PDFs,
and READMEs were removed when this volume was filed; the repository history
retains them. The verification programs were retained, unchanged.

The first report supplies the spine: its chapter structure, notation,
theorem environments, and the statements and proofs of every result that
two or three reports shared, since it was the most complete on the shared
core. The second report contributed the partition-profile kernel and the
ramified master theorem (\cref{q2:sec:kernel,q2:sec:ramified}), the negative
upper index (\cref{q2:sec:negative}), the weighted Gaussian transform and the
failure of ordinary translation together with the quantum-plane binomial
theorem (\cref{q2:sec:transforms}), the \(q\)-integer Newton expansion, the
worked coordinate-Bell examples, the reciprocal Lagrange example, the
area-type \(q\)-Catalan equation, the Good-determinant entries, the
normalization audit and crosswalk tables, and the analytic checklist
(\cref{q2:sec:analytic}). The third report contributed the finite models and
flag counts, the Gaussian annihilator, the general node calculus with the
affine-geometric master connection (\cref{q3:sec:newton}), the
affine-geometric \(q\)-Stirling forms, the colored and type-\(B\) arrays, the
\(q\)-Dobinski identity, the \(q\)-Eulerian chapter, the iterated product
rule, the path form of the chain rule (\cref{q3:sec:chain}), the \(q\)-Appell
check, the reciprocal-power binomial form, the slot-tree theorem, the fourth
central moment and the Stirling form of the Eulerian rational functions, the
additive-coordinate root prescription, the three asymptotic regimes
(\cref{q3:sec:asymptotic}), the geometric-uniform moment application
(\cref{q3:sec:uniform}), the numerical double-scaling check, the
formula-selection guide, and the list of directions not claimed.

Where the reports proved the same theorem, the following identifications
were checked before one statement was kept: the three master formulae (the
geometric-product coefficient formula, the ramified kernel at
\(d_r=1\), and the third report's Bell-product theorem) agree term by term
with \(P_r=-L_r\); the three \(q\)-Stirling conventions coincide, all three
defining \(S_q\) through the monic basis on the nodes \([j]_q\) and the
normal-ordering array as \(q^{\binom k2}S_q\); the \(q=1\) cumulants
\(\kappa_r=\beta_r\Delta_r/r\) are identical in all three; the coordinate
Bell family \(\qB nk\) of the spine is the second report's
\(\mathcal B^{(q)}_{n,k}\) and the third report's \(q\)-normalized partial
profile; the slot-tree coefficients of the second and third reports are the
same polynomials; and the two root-of-unity input systems (polylogarithmic
and Stirling-form) are the same rational functions. No discrepancy was found.

\section{Notation dictionary}
The consolidated text uses the spine's notation throughout. Passages
imported verbatim from the other two reports were transformed by macro
aliases, so that the printed symbols are uniform; their original macro names
are recorded here for anyone consulting the retained verifiers, whose
identifiers follow their own report.
\begin{center}\small
\begin{tabular}{@{}lll@{}}
\toprule
This volume & Second report & Third report\\
\midrule
\(\GB nk\) & \verb|\qb{n}{k}|, \verb|\qbb{n}{k}{q}| & \verb|\G{n}{k}|, \verb|\Gb{n}{k}{q}|, \(G_{n,k}(q)\)\\
\([n]_q\), \([n]_q!\) & \verb|\qi|, \verb|\qf| & \verb|\qnum|, \verb|\qfac|, \(Q_n=[n]_q!\)\\
\((a;q)_N\) & \verb|\poch{a}{q}{N}|, \(D_n=(q;q)_n\) & \verb|\qp{a}{N}|\\
\([z^n]\) & \verb|\coef| & \verb|\coef|\\
\(\Y n\), \(\B nk\) & \(B_n\), \(\Bhat_{n,k}\), \(\Ecal_n=\Y n((r-1)!L_r)/n!\) & \(\mathcal Y_n\), \(\Bcal_{n,k}\)\\
\(\qB nk\) & \(\mathcal B^{(q)}_{n,k}\) & \(C^{(q)}_{n,k}\)\\
\(P_r\) (power sums) & \(-L_j\) & \(-L_j\)\\
\(\widetilde S_q(n,k)\) & \(q^{\binom k2}S_q(n,k)\) & \(C_q(n,k)\)\\
\(\beta_r\) (Bernoulli) & \(B_j\) & \(B_j\)\\
\(\logq=q\,\partial_q\) & \(q\,\partial_q\) & \(\theta\)\\
\(\Li_{-s}(u)\) & \(-h_{s+1}\) & \(R_{s+1}(x)\)\\
\(w_{N,j}\), \(\rho\) & \(w_{N,j}\), \(q\) & \(w_{d,j}(s)\)\\
\bottomrule
\end{tabular}
\end{center}

\section{Formalization status}
No Lean files were added or compiled for this volume. The canonical
manuscript's formalization register covers its own results; the results here
are candidates for that register, in the order suggested at the end of
\cref{q1:qcc:computation}: polynomial Gaussian identities and
denominator-free transforms first, then finite partition kernels, then
normalized composition and Lagrange inversion, and finally cyclotomic
regularization and the analytic regimes.

\section{Files and commands}
The package directory contains \path{Gaussian_Coefficient_Calculus.tex},
its compiled PDF, \path{README.md}, and \path{verification/} with the three
retained programs and their reports. The volume loads
\path{docs/fabius-notation.tex} by relative path and has no other
dependency; with a \TeX{} installation containing Libertinus (or, as a
fallback, Latin Modern) and the listed standard packages, compile with
three passes of
\begin{verbatim}
pdflatex -interaction=nonstopmode -halt-on-error \
  Gaussian_Coefficient_Calculus.tex
\end{verbatim}
The first and third verifiers need Python 3.9 or later and no third-party
packages; the second needs SymPy (see its \path{requirements.txt}); the
third's optional double-scaling experiment needs \texttt{mpmath}. Run each
with \texttt{python verify\_identities.py} inside its directory. The
programs recompute their certificates; they do not read the saved reports.

% ---- hand chunk 29-bibliography.tex ----

\backmatter
\begin{thebibliography}{99}
\addcontentsline{toc}{chapter}{Bibliography}

\bibitem{ccc}
Vladimir Reshetnikov, \emph{Combinatorial Coefficient Calculus: Stirling,
Bell, Eulerian, and Bernoulli--Euler Families; Bell-Polynomial and Fa\`a di
Bruno Calculus; Lagrange--B\"urmann and Multivariate Good Inversion},
canonical \LaTeX{} monograph in the ProveIt repository. Consulted
September 4, 2026.
\url{https://github.com/VladimirReshetnikov/ProveIt/tree/main/Analysis/FabiusFunction/docs/semi-formalized-research-frontiers/drafts/combinatorial-coefficient-calculus/Combinatorial_Coefficient_Calculus}.

\bibitem{ccc-provenance}
ProveIt repository, \emph{Provenance and consolidation boundary},
\texttt{PROVENANCE.md} in the canonical package cited above. Consulted
September 4, 2026.

\bibitem{proveit-primary}
Vladimir Reshetnikov, \emph{The Fabius Function and Rvachev's Up-Function:
Exact arithmetic, probability, Fourier analysis, and corrected sharp
asymptotics}, ProveIt repository. Consulted September 4, 2026.
\url{https://github.com/VladimirReshetnikov/ProveIt/tree/main/Analysis/FabiusFunction/docs/Fabius_Function_and_Rvachev_Up}.

\bibitem{dlmf172}
NIST Digital Library of Mathematical Functions, \S17.2,
\emph{q-Calculus}. Definitions of shifted factorials, Gaussian coefficients,
finite and infinite binomial identities, and Jackson differentiation.
\url{https://dlmf.nist.gov/17.2}.

\bibitem{dlmf175}
NIST Digital Library of Mathematical Functions, \S17.5,
\emph{\({}_0\phi_0,{}_1\phi_0,{}_1\phi_1\) Functions}.
Infinite \(q\)-binomial identities and \(q\)-exponential conventions.
\url{https://dlmf.nist.gov/17.5}.

\bibitem{dlmf176}
NIST Digital Library of Mathematical Functions, \S17.6,
\emph{\({}_2\phi_1\) Function}. Terminating \(q\)-Chu--Vandermonde
summations.
\url{https://dlmf.nist.gov/17.6}.

\bibitem{dlmf26}
NIST Digital Library of Mathematical Functions, \S26.9,
\emph{Integer Partitions: Restricted Number and Part Size}.
\url{https://dlmf.nist.gov/26.9}.

\bibitem{dlmf24}
NIST Digital Library of Mathematical Functions, \S24.2,
\emph{Definitions and Generating Functions} for Bernoulli numbers and
polynomials.
\url{https://dlmf.nist.gov/24.2}.

\bibitem{dlmfem}
NIST Digital Library of Mathematical Functions, \S2.10,
\emph{Sums and Sequences}. Euler--Maclaurin formula and remainder.
\url{https://dlmf.nist.gov/2.10}.

\bibitem{dlmfgamma}
NIST Digital Library of Mathematical Functions, \S5.11,
\emph{Asymptotic Expansions}. Logarithmic Stirling expansion and error
control.
\url{https://dlmf.nist.gov/5.11}.

\bibitem{cai}
Yue Cai, Richard Ehrenborg, and Margaret Readdy,
``\(q\)-Stirling identities revisited,''
\emph{Electronic Journal of Combinatorics} \textbf{25}(1) (2018), Paper P1.37.
\url{https://doi.org/10.37236/6829}.

\bibitem{cigler}
Johann Cigler,
\emph{\(q\)-Catalan numbers and \(q\)-Narayana polynomials},
arXiv:math/0507225 (2005).
\url{https://arxiv.org/abs/math/0507225}.

\bibitem{dilcher}
Karl Dilcher, ``Some \(q\)-series identities related to divisor functions,''
\emph{Discrete Mathematics} \textbf{145} (1995), 83--93.
\url{https://www.sciencedirect.com/science/article/pii/0012365X9500092B}.

\bibitem{dingzeng}
Ming-Jian Ding and Jiang Zeng,
``Some identities involving \(q\)-Stirling numbers of the second kind in
type B,'' \emph{Electronic Journal of Combinatorics} \textbf{31}(1) (2024),
Paper P1.36.
\url{https://doi.org/10.37236/12147}.

\bibitem{floaterlyche}
Michael S. Floater and Tom Lyche,
``Two chain rules for divided differences and Fa\`a di Bruno's formula,''
\emph{Mathematics of Computation} \textbf{76} (2007), 867--877.
\url{https://doi.org/10.1090/S0025-5718-06-01916-8}.

\bibitem{formichella}
Sam Formichella and Armin Straub,
``Gaussian binomial coefficients with negative arguments,''
\emph{Annals of Combinatorics} \textbf{23}(3) (2019), 725--748.
\url{https://arxiv.org/abs/1802.02684}.

\bibitem{gessel}
Ira M. Gessel, ``Lagrange inversion,''
\emph{Journal of Combinatorial Theory, Series A} \textbf{144} (2016), 212--249.
\url{https://arxiv.org/abs/1609.05988}.

\bibitem{harmanhopkins}
Nate Harman and Sam Hopkins,
\emph{Quantum integer-valued polynomials},
arXiv:1601.06110 (2016).
\url{https://arxiv.org/abs/1601.06110}.

\bibitem{voigt}
Bernd Voigt,
\emph{A common generalization of binomial coefficients, Stirling numbers
and Gaussian coefficients},
Publications de l'Institut de Recherche Math\'ematique Avanc\'ee,
Strasbourg, 229/S-08 (1984), Actes du 8e S\'eminaire Lotharingien de
Combinatoire, pp.~87--89.
\url{https://www.mat.univie.ac.at/~slc/opapers/s08voigt.pdf}.

\end{thebibliography}
\end{document}

